% VNA–GIR: cơ sở lý thuyết, chứng minh và thuật toán (bản tổng hợp chặt chẽ)
% Biên dịch: xelatex vna_gir_theory.tex (hai lần)
\documentclass[11pt,a4paper]{article}

\usepackage{fontspec}
\setmainfont{Latin Modern Roman}
\setsansfont{Latin Modern Sans}
\setmonofont{Latin Modern Mono}

\usepackage[margin=2.3cm]{geometry}
\usepackage{amsmath,amssymb,amsthm,mathtools,bm}
\usepackage{booktabs,tabularx,array,longtable}
\usepackage{enumitem}
\usepackage{algorithm}
\usepackage[noend]{algpseudocode}
\usepackage{tikz}
\usetikzlibrary{positioning,arrows.meta,calc,fit,shapes.geometric}
\usepackage{xcolor}
\usepackage{footmisc}
\usepackage[colorlinks=true,linkcolor=blue!50!black,citecolor=green!35!black,urlcolor=blue!50!black]{hyperref}
\usepackage[capitalise,nameinlink]{cleveref}

\setlist{itemsep=2pt,topsep=4pt}
\renewcommand{\arraystretch}{1.18}
\setlength{\LTpre}{6pt}\setlength{\LTpost}{6pt}

% ---------- môi trường định lý ----------
% Định lý và bổ đề mang tên cố định (T1–T24, L1–L10) để khớp với mã kiểm số.
\newcommand{\thmtag}{}
\theoremstyle{plain}
\newtheorem*{namedthminner}{\thmtag}
\newenvironment{namedthm}[2]{\renewcommand{\thmtag}{#1}\begin{namedthminner}[#2]\phantomsection\label{thm:#1}}{\end{namedthminner}}
\newcommand{\Tref}[1]{\hyperref[thm:#1]{#1}}
\theoremstyle{definition}
\newtheorem*{namedassinner}{\thmtag}
\newenvironment{namedass}[1]{\renewcommand{\thmtag}{#1}\begin{namedassinner}\phantomsection\label{thm:#1}}{\end{namedassinner}}
\newtheorem{definition}{Định nghĩa}
\theoremstyle{remark}
\newtheorem*{remark}{Nhận xét}
\newenvironment{hequa}{\par\smallskip\noindent\textbf{Hệ quả cho thuật toán.}\ }{\par\smallskip}
\newenvironment{kiem}{\par\noindent\textbf{Kiểm số.}\ \small}{\normalsize\par\smallskip}
\renewcommand{\proofname}{Chứng minh}

\crefname{section}{Mục}{Mục}
\crefname{subsection}{Mục}{Mục}
\crefname{appendix}{Phụ lục}{Phụ lục}
\crefname{equation}{}{}
\crefformat{equation}{(#2#1#3)}
\crefname{algorithm}{Thuật toán}{Thuật toán}
\crefname{figure}{Hình}{Hình}
\crefname{table}{Bảng}{Bảng}
\crefname{definition}{Định nghĩa}{Định nghĩa}
\renewcommand{\figurename}{Hình}
\renewcommand{\tablename}{Bảng}
\renewcommand{\abstractname}{Tóm tắt}
\renewcommand{\refname}{Tài liệu tham khảo}
\renewcommand{\contentsname}{Mục lục}
\renewcommand{\appendixname}{Phụ lục}
\floatname{algorithm}{Thuật toán}
\algrenewcommand\algorithmicrequire{\textbf{Vào:}}
\algrenewcommand\algorithmicensure{\textbf{Ra:}}

% ---------- ký hiệu ----------
\newcommand{\R}{\mathbb{R}}
\newcommand{\T}{^{\top}}
\newcommand{\range}{\operatorname{range}}
\newcommand{\nullsp}{\operatorname{null}}
\newcommand{\tr}{\operatorname{tr}}
\newcommand{\diag}{\operatorname{diag}}
\newcommand{\blkdiag}{\operatorname{blkdiag}}
\newcommand{\sym}{\operatorname{sym}}
\newcommand{\rank}{\operatorname{rank}}
\newcommand{\argmin}{\operatorname*{arg\,min}}
\newcommand{\argmax}{\operatorname*{arg\,max}}
\newcommand{\eO}{\tilde e_O}
\newcommand{\dN}{d_N^{\beta}}
\newcommand{\drep}{d_{\mathrm{rep}}}
\newcommand{\leak}{\operatorname{leak}}
\newcommand{\rperp}{r_{\perp}}
\newcommand{\xGN}{x_{\mathrm{GN}}}
\newcommand{\Wcal}{\mathcal{W}}
\newcommand{\Jcal}{\mathcal{J}}
\newcommand{\Vcal}{\mathcal{V}}
\newcommand{\norm}[1]{\lVert #1\rVert}
\newcommand{\st}[1]{\textsf{\small #1}}   % nhãn trạng thái
\newcommand{\unv}{\footnote{Tài liệu này được ghi theo trí nhớ, chưa được kiểm lại tại nguồn trong các vòng survey của dự án; cần kiểm trước khi nộp.}}

\title{\textbf{VNA--GIR: chiếu động lực của mạng mở rộng ảo về kiến trúc cố định}\\[4pt]
\large Cơ sở lý thuyết, chuỗi suy diễn, chứng minh và thuật toán}
\author{Bản tổng hợp kỹ thuật của dự án}
\date{18/10/2026 --- bản lý thuyết v3 + T19--T24 (đã qua 5 vòng phản biện độc lập)}

\begin{document}
\maketitle

\begin{abstract}
Tài liệu này tổng hợp toàn bộ cơ sở toán học của phương pháp VNA--GIR. Ý tưởng là: một mạng có kích thước cố định được huấn luyện bằng optimizer gốc của nó; định kỳ, ta \emph{tưởng tượng} mạng có thêm vài neuron (neuron ảo), tính vận tốc trong không gian output mà các neuron này sẽ đóng góp, và chiếu vận tốc đó về các tham số đang có theo nguyên lý biến phân Dirac--Frenkel với metric là Hessian của loss theo output. Hiệu chỉnh thu được, ký hiệu $d$, là \emph{phép chiếu chính xác phần ảo} của động lực mạng mở rộng, tính trên cùng cotangent $z$ mà bước Gauss--Newton có damping tách ra từ residual. Trong học tăng dần theo lớp (CIL), hiệu chỉnh được tách thành phần sửa lớp cũ và phần học mới; phần học mới bị chặn rò rỉ sang output lớp cũ bằng một nhân tử $\beta$ là nhân tử KKT của một ngân sách quên không có tham số tự do. Trong bài toán mất plasticity, phần residual mà phép chiếu không với tới được xử lý bằng GIR: lắp neuron ảo vào slot đã chết với trọng số ra bằng $0$, nên hàm số được bảo toàn chính xác tại lúc lắp.

Mọi khẳng định trong tài liệu được gắn một trong năm trạng thái: \st{proven}, \st{assumption}, \st{empirical}, \st{interpretation}, \st{measured}. Mỗi định lý đi kèm phần ``Hệ quả cho thuật toán'' chỉ rõ nó làm cơ sở cho bước nào. Các đồng nhất và bất đẳng thức được kiểm số ở float64 (144/144 phép kiểm cho T1--T18 với 29 đối chứng đột biến; 15/15 cho T19--T24), và đã qua năm vòng phản biện độc lập. Tài liệu cũng ghi rõ những khẳng định đã bị rút lại và những câu hỏi thực nghiệm mà lý thuyết không trả lời được.
\end{abstract}

\tableofcontents
\newpage

\section{Giới thiệu và phạm vi đóng góp}\label{sec:intro}

\subsection{Bài toán}

Ta xét một mạng có kiến trúc \emph{cố định} được huấn luyện liên tục trên một chuỗi task. Hai hiện tượng làm việc này khó:
\begin{itemize}
  \item \textbf{Mất plasticity} \cite{dohare2024}: sau nhiều task, mạng mất khả năng học tiếp; nhiều neuron rơi vào trạng thái chết (dormant) \cite{sokar2023redo}.
  \item \textbf{Quên thảm họa} trong học tăng dần theo lớp: khi học lớp mới, output trên lớp cũ bị trôi; hiện tượng mạnh nhất ngay sau khi chuyển task \cite{delange2023gap}.
\end{itemize}
Các phương pháp \emph{mở rộng thật} (thêm neuron, thêm nhánh) giải quyết dung lượng bằng cách làm model lớn dần. Câu hỏi của chúng tôi là: có thể lấy được \emph{động lực} mà một mạng lớn hơn sẽ có, nhưng giữ nguyên kích thước mạng, không?

\subsection{Ý tưởng trong một đoạn}

Ở mỗi lần áp dụng (mỗi $K$ iteration), ta:
\begin{enumerate}
  \item dùng một bước Gauss--Newton có damping để tách residual $e$ thành phần mà mạng hiện tại với tới được, $z$, và phần còn lại, $\rperp=e-z$ (\Tref{T20});
  \item chọn vài neuron ảo theo tiêu chí grow (tiny\_V) và tính vận tốc output $K_E z$ mà chúng sẽ đóng góp;
  \item chiếu vận tốc đó về tham số đang có theo Dirac--Frenkel với metric $H$, được hiệu chỉnh $d$ (\Tref{T19}); trong CIL, tách $d$ thành phần sửa lớp cũ và phần học mới bị chặn rò rỉ (\Tref{T1}, \Tref{T3});
  \item kiểm chứng chỉ trên đúng vector sẽ áp dụng; bỏ qua nếu chứng chỉ âm (\Tref{T5}, \Tref{T9});
  \item nếu có slot chết, lắp neuron mới chọn theo $\rperp$ vào slot với trọng số ra bằng $0$ (\Tref{T7}).
\end{enumerate}
Optimizer gốc (SGD của iCaRL, BP hay CBP) chạy ở mọi iteration như bình thường, \emph{không} bị thay đổi.

\subsection{Đóng góp (phạm vi trung thực)}

\begin{enumerate}[label=(C\arabic*)]
  \item \textbf{Một họ cập nhật duy nhất.} Hiệu chỉnh $d$ là phép chiếu Dirac--Frenkel có damping \emph{chính xác} của phần ảo của vận tốc mạng mở rộng tại $w=0$, trên cotangent $z$ (\Tref{T19}). Mọi ablation (thay $K_E$ bằng $K$, chọn $A$ ngẫu nhiên, $\beta=0$, không swap, $\lambda=0$) là thành viên của cùng họ đó; chúng chỉ khác nhau ở \emph{kernel} đưa vào cùng một phép chiếu (\cref{sec:ablation}).
  \item \textbf{Chứng chỉ cho từng lần áp dụng.} Với $\beta=0$, $g\T d=V\ge 0$ (\Tref{T2}); với $\beta>0$, chứng chỉ được tính trực tiếp trên vector áp dụng, có dung sai từ sai số CG (\Tref{T9}); quy tắc bỏ qua bảo đảm không làm tăng loss và KD ở bậc một (\Tref{T5}).
  \item \textbf{Bảo vệ lớp cũ có nguyên tắc.} Phân rã $d^*=\drep+\dN$; $\beta$ là nhân tử KKT duy nhất của ngân sách quên $\varepsilon^*$ không có tham số tự do (\Tref{T3}); giới hạn cứng đúng là $\Lambda C(d-\drep)=0$ (\Tref{T2}).
  \item \textbf{GIR có chứng chỉ cho slot chết.} Lắp với $B_v=0$ bảo toàn hàm chính xác; lợi được cho bởi dạng Schur theo $\rperp$ cho mọi phi tuyến (\Tref{T7}); có quy tắc nhiều slot, đặt lại optimizer, BN (\Tref{T13}, \Tref{T17}).
  \item \textbf{Các trường hợp suy biến và phạm vi} (V$_O=0$, $e_O=0$, BN, momentum, CG không chính xác) được xử lý tường minh (\Tref{T9}--\Tref{T11}, \Tref{T18}).
\end{enumerate}

\textbf{Những gì tài liệu \emph{không} khẳng định:}
\begin{itemize}
  \item không khẳng định mạng ``train như một mạng lớn hơn'' theo nghĩa quỹ đạo; bảo đảm là bậc một, cho từng lần áp dụng (\Tref{T8});
  \item không khẳng định bước SGD của optimizer gốc là một phần của cùng phép chiếu; đó chỉ là một \st{interpretation} xấp xỉ (\Tref{T20}, \Tref{T22});
  \item không khẳng định tương đương với DER hay ``growth với hàm cũ đóng băng'';
  \item không khẳng định gì về accuracy trên tập test; câu hỏi ``kernel ảo $K_E$ có hơn kernel của chính mạng $K$ không'' là câu hỏi \emph{thực nghiệm} (\cref{sec:open}).
\end{itemize}

\subsection{Quy ước trạng thái}

\begin{tabularx}{\linewidth}{@{}lX@{}}
\toprule
Trạng thái & Ý nghĩa \\
\midrule
\st{proven} & có chứng minh đầy đủ dưới các giả thiết đã liệt kê \\
\st{assumption} & giả thiết được nêu rõ, có cách kiểm lúc chạy hoặc lý do \\
\st{empirical} & chỉ có bằng chứng số (thường trên toy) \\
\st{interpretation} & diễn giải hay liên hệ, không phải định lý \\
\st{measured} & đại lượng đo chính xác lúc chạy nhưng không có bảo đảm về hệ quả \\
\bottomrule
\end{tabularx}

\medskip
Tên các định lý (T1--T24) và bổ đề (L1--L10) giữ đúng như trong tài liệu nghiên cứu và mã kiểm số, để có thể đối chiếu trực tiếp với \texttt{final\_check\_v3.py} và \texttt{check\_unified\_gn.py} (\cref{app:verification}).

\section{Thiết lập, ký hiệu và giả thiết}\label{sec:setting}

\subsection{Cửa sổ, loss và residual}

\textbf{Cửa sổ.} $\Wcal$ gồm $N$ mẫu (dữ liệu mới cộng exemplar, sau augmentation), forward ở train mode. $f(\theta)\in\R^n$ là logit xếp chồng, $n=N c$ với $c$ là số lớp hiện tại; $\theta_S\in\R^p$ là các tham số được hiệu chỉnh; $J=\partial f/\partial\theta_S\in\R^{n\times p}$ và $K=JJ\T$.

\textbf{Loss} (đúng như iCaRL trong PyCIL \cite{rebuffi2017icarl,zhou2023pycil}): $L=L_{CE}+L_{KD}$ với
\[
L_{CE}=\tfrac1N\textstyle\sum_i \mathrm{CE}(f_i,y_i),\qquad
L_{KD}=-\tfrac1N\textstyle\sum_i p_{t,i}\T\log\operatorname{softmax}(f_{i,\mathrm{old}}/T),
\]
tức trung bình, \emph{không} nhân $T^2$. $O\in\{0,1\}^{m\times n}$ chọn $m=Nc_{\mathrm{old}}$ logit cũ, $OO\T=I$.

\textbf{Residual và Hessian theo output.} $e=\partial L/\partial f=e_N+\eO$, với $e_{N,i}=(p_i-\mathbf 1_{y_i})/N$, $\eO=O\T e_O$, $e_{O,i}=(p_{s,i}-p_{t,i})/(NT)$ và $p_{s,i}=\operatorname{softmax}(f_{i,\mathrm{old}}/T)$ là xác suất của \emph{student}. Hessian khối chéo theo mẫu:
\[
H_i=\tfrac1N\big(D(p_i)-p_ip_i\T\big)+O_i\T W_iO_i,\qquad W_i=\tfrac{1}{NT^2}\big(D(p_{s,i})-p_{s,i}p_{s,i}\T\big),
\]
$W=\blkdiag(W_i)$.

\subsection{Bảng ký hiệu}

\begin{longtable}{@{}p{0.27\linewidth}p{0.68\linewidth}@{}}
\toprule
Ký hiệu & Định nghĩa \\
\midrule
\endhead
$C=OJ$, $B=C\T WC$, $F=J\T HJ$ & Jacobian logit cũ; metric KD trên tham số; Gauss--Newton \\
$g=J\T e$, $g_O=J\T\eO=C\T e_O$ & gradient toàn phần; gradient KD \\
$\mu>0$, $A_\beta=F+\beta B+\mu I$ & damping; ma trận của bài toán có bảo vệ \\
$G=A_0^{-1}$, $G_\beta=A_\beta^{-1}$ & nghịch đảo có damping \\
$z(r)=HJGJ\T r$; $z=z(e)$, $z_O=z(\eO)$, $z_N=z-z_O$ & cotangent đã tách (\Tref{T20}) \\
$\xGN=Gg$ & bước Gauss--Newton có damping (lần giải đầu tiên ở bước 1) \\
$\rperp(r)=r-z(r)=\mu(HK+\mu)^{-1}r$ & phần bù theo metric $H$ \\
$A=[a_1,\dots,a_r]$, $J_{a_j}=\partial f/\partial w_j|_{w_j=0}$ & hướng vào và Jacobian theo trọng số ra của neuron ảo \\
$K_E=\sum_jJ_{a_j}J_{a_j}\T$ & kernel ảo \\
$u=K_Ez$, $u_N=K_Ez_N$, $u_R=K_Ez_O$ & vận tốc ảo (toàn phần, học mới, sửa) \\
$V=z\T K_Ez$, $V_O=z_O\T K_Ez_O$ & năng lượng ảo; năng lượng sửa \\
$\rho=V/\norm g^2$ & tỉ số ảo trên nền (\Tref{T21}) \\
$\drep=GJ\T Hu_R$, $\dN=G_\beta J\T Hu_N$, $d^*(\beta)=\drep+\dN$ & hiệu chỉnh \\
$\Phi(d)=\norm{Jd-u}_H^2+\mu\norm d^2$, $\Phi_N(x)=\norm{Jx-u_N}_H^2+\mu\norm x^2$ & độ trung thành của phép chiếu \\
$\leak(\beta)=\norm{C\dN}_W^2$, $h(\beta)=g_O\T d^*(\beta)$ & rò rỉ; chứng chỉ KD bậc một \\
$\chi^2=\norm{e_O}^2_{W^+}$ (\Tref{L6}) & khoảng cách teacher--student \\
$\varepsilon^*=V_O^2/\chi^2$, $\beta_{\mathrm{suf}}=\chi^2\norm{u_N}_H^2/V_O^2$ & ngân sách quên; cận trên của $\beta$ \\
$\Lambda$, $\Lambda\T\Lambda=W$; $\hat C=\Lambda C$ & thừa số chính xác của $W$ (\Tref{L7}) \\
$X$, $\tilde X=X-\mathbf 1\bar x\T$; $S_x=\tilde X\T\tilde X/N+\rho_xI$ & đầu vào của site; ma trận Gram có ridge (\Tref{L9}) \\
$N_s$: $J_a\T s=N_s\T a$ (trong một vùng mặt nạ) & đặc trưng tuyến tính hóa của neuron ảo \\
$f_u$, $J_u$, $B_u$; $J_v=\partial f/\partial B_v|_{B_v=0}$ & đóng góp, Jacobian, trọng số ra của slot $u$; neuron mới $v$ \\
$q^*_r(\Jcal)=-\tfrac12(\Jcal\T r)\T G(\Jcal)(\Jcal\T r)$ & mô hình GN của mức giảm loss trong span$(\Jcal)$ \\
$H_\perp(\Jcal)=H-H\Jcal G(\Jcal)\Jcal\T H$ & metric phần bù \\
$E(\Jcal,t)=\min_x\norm{\Jcal x-t}_H^2+\mu\norm x^2$ & hàm chiếu \\
$j_u^\ell=\ell(f-f_u)-\ell(f)$ & bước nhảy khi bỏ slot sống \\
\bottomrule
\end{longtable}

\textbf{Neuron ảo và BN.} Không có BN theo sau: đặc trưng $\varphi(a)=D(a)Xa$ với $D(a)$ là mặt nạ ReLU thực tế. Có BN train mode theo sau (kênh mới, $\gamma=1$, $\beta_{BN}=0$): $\varphi(a)=D(a)\tilde Xa/\sigma(a)$, $\sigma(a)^2=a\T\tilde X\T\tilde Xa/N$.

\textbf{Slot hợp lệ.} Slot $u$ chỉ đi tiếp qua trọng số riêng $B_u$, đặt được về $0$ mà không ảnh hưởng đơn vị khác. Trong ResNet, kênh ra của \texttt{conv1} trong một basic block là hợp lệ; kênh ra của \texttt{conv2} đi vào phép cộng residual thì không.

\subsection{Giả thiết}\label{sec:assumptions}

\begin{longtable}{@{}p{0.07\linewidth}p{0.47\linewidth}p{0.40\linewidth}@{}}
\toprule
\# & Giả thiết & Kiểm lúc chạy, hoặc lý do đúng \\
\midrule
\endhead
A1 & $e,J,H$ và bước cập nhật dùng cùng một cửa sổ và cùng forward train-mode; buffer running-stat của BN không bị forward của VNA làm đổi. & So buffer trước và sau mỗi lần áp dụng (phải bằng nhau bitwise). \\
A2 & Loss là CE trung bình cộng KD của PyCIL; mọi xác suất $>0$. & Theo định nghĩa. Ở float32 xác suất có thể underflow; \Tref{T1}(c,d) không cần range nên vẫn đúng. \\
A3 & $\mu>0$. & $\mu=\max(\mu_{\mathrm{rel}}\tr(F)/n,10^{-12})$. \\
A4$'$ & Vùng mặt nạ cố định: các hướng $a$ xét nằm trong cùng vùng $R_D=\{a:D(a)=D\}$. & Mặt nạ đo được. \emph{Chỉ} dùng cho tính tối ưu của bước chọn (\Tref{T6}, \Tref{T7}(d), \Tref{T14}). Tối ưu toàn cục khi mặt nạ đổi là A4-glob, không chứng minh. \\
A5 & Slot chết trên cửa sổ ($f_u=0$, $J_u=0$) tại đúng tham số lúc chèn, tức sau bước 5. & Đếm mẫu có tiền kích hoạt $>0$ tại $\theta'$ (phải bằng $0$). \\
A6 & $\nabla L$, $\nabla L_{KD}$ lần lượt $L_s$-, $L_K$-Lipschitz trong không gian tham số. & \emph{Không} kiểm được; chỉ dùng cho phát biểu bậc hai (\Tref{T5}(e,f)). \\
A6$'$ & $\partial^2L/\partial f^2$ là $L_3$-Lipschitz trong không gian output. & \st{proven}: dạng đóng, chặt (\Tref{L10}); không còn là giả thiết. \\
A7 & Tham chiếu tự trị: mô hình mở rộng được thay bằng $\dot y=U(y)$ trong không gian output, $U$ là $L_U$-Lipschitz. & Chỉ dùng cho \Tref{T8}. \\
A8 & CG dừng theo $\rho=\norm{b-A\hat x}/\norm b$. & Tính trực tiếp $\rho$ cuối. \\
A9 & Hiệu chỉnh cộng \emph{ngoài} momentum buffer; khi swap, trạng thái optimizer của slot đặt bằng $0$. & Kiểm bằng \texttt{torch.optim.SGD} thật (\Tref{T10}, \Tref{T13}). \\
A10 & Hồi phục trên $L$: sau swap, mạng sống hồi phục bằng một bước GN có damping trên $L$ trong chế độ tuyến tính. & \st{assumption, interpretive}. Chỉ dùng để \emph{diễn giải} lợi về sau; không đúng cho train thật (\Tref{T16}). \\
A11 & Mọi bảo đảm là cục bộ cho từng lần áp dụng, trên loss thực nghiệm của cửa sổ. & Phạm vi. \\
A12 & Slot hợp lệ như trên. & Kiểm theo kiến trúc. \\
A13 & Hiệu chỉnh lười: $\hat d$ thay đổi ít trong khoảng $K$ bước. & \st{assumption, interpretive}; chỉ dùng cho quy tắc $\lambda_{\mathrm{faithful}}=K$ (\Tref{T24}); giống giả thiết staleness của K-FAC/Shampoo. Không chứng chỉ nào dùng nó. \\
\bottomrule
\end{longtable}

\section{Cơ sở lý thuyết từ các nghiên cứu trước}\label{sec:foundations}

Mỗi mục dưới đây nêu: nguồn, \emph{chính xác} điều chúng tôi dùng, và nơi nó đi vào chuỗi suy diễn. Những nguồn chỉ được ghi theo trí nhớ (chưa kiểm lại tại nguồn trong các vòng survey của dự án) được đánh dấu bằng chú thích.

\subsection{Phép chiếu biến phân Dirac--Frenkel và các dạng số của nó}

\paragraph{Nguyên lý Dirac--Frenkel--McLachlan.} Cho một đa tạp xấp xỉ $\mathcal M$ và một động lực $\dot y=U(y)$ trong không gian lớn, nguyên lý chọn vận tốc trên không gian tiếp tuyến của $\mathcal M$ gần $U$ nhất theo một metric \cite{lubich2008book,kochlubich2007}. \emph{Điều chúng tôi dùng:} với tham số hóa $\theta\mapsto f(\theta)$, không gian tiếp tuyến là $\range(J)$, metric là $H$, và có thêm damping $\mu$; vận tốc tham số là
\[
P(t)=\argmin_x\norm{Jx-t}_H^2+\mu\norm x^2=GJ\T Ht.
\]
Đây là định nghĩa của hiệu chỉnh $d$ khi $t=K_Ez$ (\Tref{T19}).

\paragraph{Neural Galerkin \cite{bruna2024neuralgalerkin}.} Chiếu Dirac--Frenkel liên tục động lực của PDE lên đa tạp tham số của mạng, giải bằng bình phương tối thiểu có regularization. \emph{Điều chúng tôi dùng:} khuôn mẫu ``chiếu liên tục từng bước nhỏ'' và việc giải hệ chuẩn tắc có damping. Khác biệt: động lực được chiếu là của một \emph{mạng mở rộng ảo}, không phải PDE.

\paragraph{Dynamical low-rank và BUG thích nghi hạng \cite{ceruti2022bug,schotthofer2022dlrt}.} Lược đồ ``mở rộng cơ sở $\to$ bước Galerkin $\to$ cắt về hạng cố định''. \emph{Điều chúng tôi dùng:} chỉ ở mức \st{interpretation}: thêm neuron ảo tương ứng mở rộng, VNA tương ứng bước Galerkin, và GIR tương ứng việc tái dùng bậc tự do. Không có định lý nào của chúng tôi dựa vào kết quả hội tụ của các lược đồ này.

\subsection{Gauss--Newton, natural gradient và kernel}

\paragraph{Gauss--Newton và natural gradient \cite{amari1998,martens2020}.} Bước $(J\T HJ+\mu I)^{-1}J\T e$ là bước Gauss--Newton có damping (Levenberg--Marquardt), đồng thời là natural gradient theo metric output. \emph{Điều chúng tôi dùng:} bước $\xGN=Gg$ làm bộ tách residual $e=z+\rperp$ (\Tref{T20}), và mô hình bậc hai $q^*$ của mức giảm loss (\cref{sec:setting}).

\paragraph{Energy natural gradient \cite{muller2023engd}.} Natural gradient trong không gian hàm như một phép chiếu lên không gian tiếp tuyến, kèm damping, chạy được trên mô hình nhỏ. \emph{Điều chúng tôi dùng:} tiền lệ cho tính khả thi; đồng thời ENGD/GN chính là trường hợp $K_E=0$ của khung chiếu, nên là đối chứng bắt buộc.

\paragraph{Neural tangent kernel \cite{jacot2018ntk}.} $K=JJ\T$ là kernel tiếp tuyến thực nghiệm. \emph{Điều chúng tôi dùng:} ngôn ngữ kernel để phát biểu bổ đề chiếu lồng nhau ($Jd\in\range(K)$) và phần bù Schur $K_{NS}$ (\Tref{T2}). Không dùng kết quả nào về giới hạn độ rộng vô hạn.

\paragraph{K-FAC, Shampoo \cite{martens2015kfac,gupta2018shampoo,purifyingshampoo}.} Các preconditioner bậc hai này xấp xỉ độ cong của \emph{chính} mạng hiện có và cập nhật preconditioner mỗi vài chục đến vài trăm bước. \emph{Điều chúng tôi dùng:} (i) tiền lệ cho việc áp hiệu chỉnh mỗi $K\approx100$ bước, với giả thiết staleness tương tự A13; (ii) điểm phân biệt: VNA là preconditioner $P=GJ\T HK_EHJG\succeq0$ dựng từ kernel $K_E$ của neuron \emph{không tồn tại}, thứ mà không phương pháp độ cong nào tạo ra.

\subsection{Tiêu chí grow}

\paragraph{TINY \cite{verbockhaven2024tiny}, GradMax \cite{evci2022gradmax}, Firefly \cite{wu2020firefly}.} Các phương pháp này \emph{thêm} neuron thật và khởi tạo chúng theo hướng làm giảm loss nhanh nhất (cực đại chuẩn gradient theo trọng số ra, hoặc giải bài toán trị riêng trên thống kê đặc trưng). \emph{Điều chúng tôi dùng:} bài toán chọn hướng vào $A$ là bài toán trị riêng suy rộng $(N_sN_s\T,S_x)$, cùng họ với TINY; chúng tôi dùng nó \emph{ảo} (VNA) hoặc để \emph{khởi tạo lại slot chết} (GIR). Ở $\beta=0$, chúng tôi chứng minh tiny\_V với cotangent $z$ là lựa chọn bậc một tối ưu trong một vùng mặt nạ (\Tref{T6}).

\subsection{Mất plasticity và reset}

\paragraph{Mất plasticity \cite{dohare2024}; CBP \cite{dohare2024}; ReDo \cite{sokar2023redo}.} Mạng mất khả năng học tiếp; nhiều đơn vị chết. CBP và ReDo chọn đơn vị theo utility hoặc theo độ hoạt động rồi \emph{khởi tạo lại ngẫu nhiên}, trọng số ra bằng $0$. \emph{Điều chúng tôi dùng:} cơ chế ``trọng số ra bằng 0 thì bảo toàn hàm'' và tiêu chí đơn vị chết; điểm khác của GIR là hướng vào được chọn theo $\rperp$ (\Tref{T7}). Các phương pháp reset khác tìm được trong survey (ReGraMa, SWR, Calibrated Partial Resets) đều khởi tạo ngẫu nhiên hoặc bằng $0$.

\subsection{Học liên tục}

\paragraph{iCaRL \cite{rebuffi2017icarl}.} Exemplar cộng KD với model cũ. \emph{Điều chúng tôi dùng:} đúng loss của PyCIL làm $L$; model cũ làm teacher $p_t$; exemplar nằm trong cửa sổ.

\paragraph{Chiếu gradient: GPM \cite{saha2021gpm}, Adam-NSCL \cite{wang2021nscl}, TRGP \cite{lin2022trgp}, SGP \cite{saha2023sgp}.} Chiếu cập nhật vào phần bù của không gian đặc trưng cũ, cứng (GPM, Adam-NSCL) hoặc mềm (TRGP, SGP). TRGP chỉ ra ràng buộc cứng hại việc học task mới. \emph{Điều chúng tôi dùng:} (i) lý do chọn ràng buộc \emph{mềm} với $\beta$ hữu hạn; (ii) phép chiếu null-space ở giới hạn $\beta\to\infty$ (\Tref{T2}(b,d)). Phần null space \emph{không} phải đóng góp của chúng tôi.

\paragraph{NCL \cite{kao2021ncl}.} Natural gradient cộng chiếu, bảo vệ trong suốt quá trình tối ưu; chiếu kiểu OWM xấp xỉ một preconditioner Fisher. \emph{Điều chúng tôi dùng:} tiền lệ cho ``preconditioner + chiếu'' (không mới), và lập luận áp bảo vệ ở \emph{mỗi} lần hiệu chỉnh.

\paragraph{FROMP \cite{pan2020fromp}, EFC \cite{magistri2024efc}.} Regularize trong không gian hàm qua kernel GP/GGN (FROMP); phạt độ trôi đặc trưng bậc hai theo các hướng quan trọng (EFC). \emph{Điều chúng tôi dùng:} tiền lệ cho phạt \emph{theo metric KD}, gần với $\beta\norm{C(d-\drep)}_W^2$. EFC là hàng xóm hình thức gần nhất của số hạng bảo vệ.

\paragraph{Stability gap \cite{delange2023gap,harun2024gap}.} Quên mạnh ngay đầu task rồi hồi phục; giữ output cũ trong lúc học giảm gap. \emph{Điều chúng tôi dùng:} dự đoán kiểm được (lợi tập trung sau khi chuyển task) và giao thức đánh giá theo iteration.

\paragraph{C-Flat \cite{bian2024cflat,cflatpp}.} Plug-in optimizer trên PyCIL cùng setting (ResNet32, CIFAR-100); lợi trên iCaRL khoảng $+0.3$ đến $+0.95$ điểm với chi phí khoảng $2.9\times$. \emph{Điều chúng tôi dùng:} mốc hiệu quả và chi phí của một đóng góp cùng loại.

\subsection{Liên hệ ngoài ngành (chỉ là diễn giải)}

\paragraph{Nguyên lý ràng buộc tối thiểu của Gauss và phương trình Udwadia--Kalaba \cite{udwadia1992,udwadia2005}.} Bài toán $\min\Phi(d)$ với ràng buộc $\norm{C(d-\drep)}_W^2\le\varepsilon^*$ (\Tref{T3}(f)) là một bài toán ``vận tốc khả thi gần nhất'' theo metric $(H,\mu)$. \st{interpretation}.

\paragraph{Điều khiển ưu tiên tác vụ \cite{nakamura1987,siciliano1991,chiaverini1997}.} Ưu tiên 1 là giữ hàm cũ, ưu tiên 2 là học mới trong null space của ưu tiên 1; ở giới hạn cứng đúng là \Tref{T2}(d). Khi null space cạn, robot thêm bậc tự do; ở đây là GIR. \st{interpretation}.

\subsection{Hàng xóm cần phân biệt}

\begin{tabularx}{\linewidth}{@{}lX@{}}
\toprule
Công trình & Khác biệt với VNA \\
\midrule
LoRA-Pro \cite{wang2025lorapro} & Chiếu gradient của fine-tune toàn phần lên cấu trúc low-rank (tuyến tính, metric Frobenius trên gradient). VNA: neuron ảo \emph{phi tuyến}, chiếu trong không gian output với metric $H$, kiến trúc bất kỳ, có chứng chỉ và bảo vệ lớp cũ. \\
NetAug \cite{cai2022netaug} & Train mạng nhỏ nhờ các bản mở rộng dùng chung trọng số; có tham số thật lúc train và loss phụ. VNA không có tham số thêm và chiếu giải tích. \\
Virtual Width Networks \cite{vwn2025}, FPE \cite{fpe2025} & Trùng thuật ngữ; mở rộng biểu diễn hoặc chia neuron thật, không có động lực ảo được chiếu. \\
CL through Control Minimization \cite{dehaan2026control} & Cân bằng tín hiệu học và tín hiệu giữ; gần số hạng $\beta$; không có dung lượng ảo, không tái dùng slot. \\
FOSTER, Progress \& Compress & Mở rộng thật rồi nén ở cuối; VNA chiếu liên tục từng bước, không train model phụ. \\
\bottomrule
\end{tabularx}

\section{Chuỗi suy diễn}\label{sec:chain}

\Cref{fig:chain} cho sơ đồ phụ thuộc: giả thiết $\to$ bổ đề $\to$ định lý $\to$ bước thuật toán. \Cref{tab:chain} liệt kê cùng thông tin dưới dạng bảng, kèm câu trả lời cho câu hỏi ``bước này suy ra từ đâu''.

\begin{figure}[ht]
\centering
\resizebox{\linewidth}{!}{%
\begin{tikzpicture}[
  every node/.style={font=\footnotesize},
  ass/.style={draw=orange!70!black,fill=orange!8,rounded corners=2pt,minimum width=1.9cm,minimum height=0.55cm,align=center},
  lem/.style={draw=teal!70!black,fill=teal!8,rounded corners=2pt,minimum width=1.9cm,minimum height=0.55cm,align=center},
  thm/.style={draw=blue!60!black,fill=blue!6,rounded corners=2pt,minimum width=2.6cm,minimum height=0.55cm,align=center},
  stp/.style={draw=black!70,fill=black!4,minimum width=3.3cm,minimum height=0.55cm,align=center},
  ar/.style={-{Stealth[length=2mm]},draw=black!45,thin},
  node distance=0.32cm]
% cột tiêu đề
\node[font=\bfseries\small] at (0,1.0) {Giả thiết};
\node[font=\bfseries\small] at (3.3,1.0) {Bổ đề};
\node[font=\bfseries\small] at (7.4,1.0) {Định lý};
\node[font=\bfseries\small] at (12.2,1.0) {Bước thuật toán};
% giả thiết
\node[ass] (A2) at (0,0) {A2 loss CE+KD};
\node[ass,below=of A2] (A3) {A3 $\mu>0$};
\node[ass,below=of A3] (A1) {A1 cùng cửa sổ};
\node[ass,below=of A1] (A4) {A4$'$ vùng mặt nạ};
\node[ass,below=of A4] (A8) {A8 dừng CG};
\node[ass,below=of A8] (A5) {A5 slot chết};
\node[ass,below=of A5] (A12) {A12 slot hợp lệ};
\node[ass,below=of A12] (A9) {A9 momentum};
\node[ass,below=of A9] (A10) {A10 hồi phục (diễn giải)};
\node[ass,below=of A10] (A13) {A13 lười (diễn giải)};
% bổ đề
\node[lem] (L1) at (3.3,0) {L1 push-through};
\node[lem,below=of L1] (L2) {L2--L4 thứ tự, đồng nhất};
\node[lem,below=of L2] (L5) {L5 range};
\node[lem,below=of L5] (L6) {L6 $\chi^2$ dạng đóng};
\node[lem,below=of L6] (L7) {L7 thừa số $\Lambda$};
\node[lem,below=of L7] (L8) {L8 Hessian student};
\node[lem,below=of L8] (L9) {L9 ridge $S_x$};
\node[lem,below=of L9] (L10) {L10 $L_3$ chặt};
% định lý
\node[thm] (T20) at (7.4,0) {T20 tách $e=z+r_\perp$};
\node[thm,below=of T20] (T19) {T19 $d$ = chiếu phần ảo};
\node[thm,below=of T19] (T1) {T1 dạng đóng, phân rã};
\node[thm,below=of T1] (T2) {T2 $\beta=0$; giới hạn cứng};
\node[thm,below=of T2] (T3) {T3 sửa, rò rỉ, KKT};
\node[thm,below=of T3] (T6) {T6/T14 chọn $A$};
\node[thm,below=of T6] (T5) {T5 giảm loss, bỏ qua};
\node[thm,below=of T5] (T9) {T9/T18 CG không chính xác};
\node[thm,below=of T9] (T10) {T10 BN, momentum};
\node[thm,below=of T10] (T7) {T7/T13/T17/T23 GIR};
\node[thm,below=of T7] (T24) {T24 cỡ bước};
% bước
\node[stp] (S1) at (12.2,0) {S1 residual; tách $z$, $r_\perp$};
\node[stp,below=0.9cm of S1] (S2) {S2 chọn neuron ảo $A$};
\node[stp,below=0.9cm of S2] (S3) {S3 $\beta$ và $d^*=d_{\mathrm{rep}}+d_N^\beta$};
\node[stp,below=0.9cm of S3] (S4) {S4 chứng chỉ, bỏ qua};
\node[stp,below=0.9cm of S4] (S5) {S5 áp dụng ngoài momentum};
\node[stp,below=0.9cm of S5] (S6) {S6 GIR trên slot chết};
\node[stp,below=0.9cm of S6] (S7) {S7 log, chẩn đoán};
% cạnh giả thiết -> bổ đề
\foreach \a/\b in {A2/L5,A2/L6,A2/L7,A2/L8,A2/L10,A3/L1,A3/L2,A4/L9}{\draw[ar] (\a.east)--(\b.west);}
% bổ đề -> định lý
\foreach \a/\b in {L1/T20,L1/T19,L1/T1,L2/T1,L2/T5,L5/T1,L5/T3,L6/T3,L7/T2,L7/T3,L8/T3,L9/T6,L10/T7,L2/T6}{\draw[ar] (\a.east)--(\b.west);}
% giả thiết -> định lý
\foreach \a/\b in {A8/T9,A5/T7,A12/T7,A9/T10,A1/T10,A10/T7,A13/T24,A4/T6}{\draw[ar,dashed] (\a.east) to[out=0,in=180] (\b.west);}
% định lý -> bước
\foreach \a/\b in {T20/S1,T19/S3,T1/S3,T2/S3,T3/S3,T6/S2,T5/S4,T9/S4,T10/S5,T7/S6,T24/S5,T3/S7,T20/S6}{\draw[ar] (\a.east)--(\b.west);}
\end{tikzpicture}}
\caption{Sơ đồ phụ thuộc. Mũi tên liền: suy ra trực tiếp. Mũi tên đứt: giả thiết được dùng trực tiếp trong định lý. A10 và A13 chỉ là giả thiết diễn giải; không chứng chỉ nào dựa vào chúng.}
\label{fig:chain}
\end{figure}

\begin{longtable}{@{}p{0.2\linewidth}p{0.34\linewidth}p{0.4\linewidth}@{}}
\caption{Mỗi bước thuật toán suy ra từ đâu.}\label{tab:chain}\\
\toprule
Bước & Định lý / giả thiết làm cơ sở & Nội dung suy diễn \\
\midrule
\endfirsthead
\toprule
Bước & Định lý / giả thiết làm cơ sở & Nội dung suy diễn \\
\midrule
\endhead
S0 optimizer gốc & (baseline, không đổi) & Không có định lý nào về bước này; nó chỉ được \emph{diễn giải} là phần $K$ (\Tref{T20}(c), \Tref{T22}). \\
S1 tách residual & \Tref{T20}, \Tref{L1}, \Tref{L5} & Bước GN có damping tách chính xác $e=z+\rperp$; $z$ là phần mà bước GN của mạng hiện tại loại bỏ. \\
S2 chọn $A$ & \Tref{T6}, \Tref{T14}, \Tref{L9}; A4$'$ & Lợi bậc một là một dạng vết theo $A$; tối ưu là eigvec suy rộng; ở $\beta=0$ là tiny\_V$(z)$. \\
S3 $\beta$, $d^*$ & \Tref{T19}, \Tref{T1}, \Tref{T2}, \Tref{T3}, \Tref{T11}, \Tref{L6}, \Tref{L7} & $d$ là phép chiếu chính xác của phần ảo; phân rã sửa/học mới; $\beta$ là nhân tử KKT của ngân sách $\varepsilon^*$; bisection trên $\leak$ hợp lệ, trên $h$ thì không. \\
S4 chứng chỉ & \Tref{T5}, \Tref{T9}, \Tref{T18}; A8 & Chứng chỉ tính trực tiếp trên vector áp dụng; dung sai từ sai số CG; bỏ qua bảo đảm không làm tăng $L$, $L_{KD}$ ở bậc một. \\
S5 áp dụng & \Tref{T10}(d), \Tref{T24}; A9, A13 & Cộng ngoài momentum buffer giữ chứng chỉ; ghi vào buffer thì $\hat d$ bị phát lại $1/(1-\nu)$ lần. \\
S6 GIR & \Tref{T7}, \Tref{T13}, \Tref{T17}, \Tref{T23}; A5, A12 & Slot chết: bảo toàn hàm chính xác; lợi Schur theo $\rperp$; swap sau hiệu chỉnh; đặt lại optimizer và BN. \\
S7 log & \Tref{T3}, \Tref{T8}, \Tref{T21} & Các đại lượng chẩn đoán có nghĩa được chứng minh ($\rho$, $\leak$, $h$, defect). \\
\bottomrule
\end{longtable}

\section{Bổ đề}\label{sec:lemmas}

Các bổ đề dưới đây là công cụ đại số dùng lại trong mọi định lý. Ta viết $R=H^{1/2}$, $M=RKR$.

\begin{namedthm}{L1}{push-through}
$RJGJ\T R=M(M+\mu)^{-1}\preceq I$. Hệ quả: $\norm{J\T Hu}_G^2\le\norm u_H^2$ với mọi $u$.
\end{namedthm}
\begin{proof}
Đồng nhất push-through $X(X\T X+\mu)^{-1}X\T=(XX\T+\mu)^{-1}XX\T$ với $X=RJ$ cho $(RJ)(J\T R^2J+\mu)^{-1}(RJ)\T=(RJJ\T R+\mu)^{-1}RJJ\T R=M(M+\mu)^{-1}$, có trị riêng trong $[0,1)$. Do đó $u\T HJGJ\T Hu=(Ru)\T M(M+\mu)^{-1}(Ru)\le\norm{Ru}^2$.
\end{proof}

\begin{namedthm}{L2}{thứ tự}
$G_\beta\preceq G\preceq\mu^{-1}I$, vì $B,F\succeq0$.
\end{namedthm}

\begin{namedthm}{L3}{chặn bảo vệ}
$\beta G_\beta^{1/2}BG_\beta^{1/2}\preceq I$, vì $\beta B\preceq A_\beta$.
\end{namedthm}

\begin{namedthm}{L4}{đồng nhất cotangent}
$(J\T r)\T GJ\T H=z(r)\T$ với mọi $r$, vì $H$ và $G$ đối xứng.
\end{namedthm}

\begin{namedthm}{L5}{range}
Dưới A2: $\nullsp(H_i)=\operatorname{span}(\mathbf 1_c)$, $\nullsp(W_i)=\operatorname{span}(\mathbf 1_{c_{\mathrm{old}}})$, và $\mathbf 1\T e_{N,i}=\mathbf 1\T e_{O,i}=0$. Do đó $e_N,\eO,e\in\range(H)$ và $e_O\in\range(W)$.
\end{namedthm}
\begin{proof}
$x\T(D(p)-pp\T)x=\operatorname{Var}_p(x)$, bằng $0$ khi và chỉ khi $x$ là hằng (vì $p>0$). $H_i$ là tổng của phần CE (bằng $0$ khi và chỉ khi $x$ hằng) và phần KD (bằng $0$ khi và chỉ khi $x_{\mathrm{old}}$ hằng), nên $\nullsp(H_i)=\operatorname{span}(\mathbf 1)$. Các residual có tổng bằng $0$ vì $\sum p=\sum\mathbf 1_y=1$ và $\sum p_s=\sum p_t=1$. Với $\eO$: $\range(H)\supseteq\range(O\T WO)=O\T\range(W)\ni O\T e_O$.
\end{proof}

\begin{namedthm}{L6}{dạng đóng của $\chi^2$}
$\norm{e_O}_{W^+}^2=\frac1N\sum_i\chi^2(p_{t,i}\,\|\,p_{s,i})$, với $\chi^2(p_t\|p_s)=\sum_j(p_{s,j}-p_{t,j})^2/p_{s,j}$.
\end{namedthm}
\begin{proof}
Với $v\perp\mathbf 1$, $x=D^{-1}v$ thỏa $(D-pp\T)x=v-p(\mathbf 1\T v)=v$. Hình chiếu của $x$ lên range khác $x$ một bội của $\mathbf 1$, nên $v\T(D-pp\T)^+v=v\T D^{-1}v$. Thay hệ số: $e_i\T W_i^+e_i=(NT)^{-2}(NT^2)\sum_j(p_s-p_t)^2/p_s=\chi^2_i/N$.
\end{proof}

\begin{namedthm}{L7}{thừa số chính xác của $W$}
$\Lambda_i:=(NT^2)^{-1/2}\big(D(p_s)^{1/2}-\sqrt{p_s}\,p_s\T\big)$ thỏa $\Lambda_i\T\Lambda_i=W_i$ và $\nullsp(\Lambda_i)=\operatorname{span}(\mathbf 1)$.
\end{namedthm}
\begin{proof}
Đặt $S=D^{1/2}$; vì $S\sqrt p=p$, ta có $\Lambda\T\Lambda\propto(S-p\sqrt p\T)(S-\sqrt pp\T)=D-pp\T-pp\T+p(\sqrt p\T\sqrt p)p\T=D-pp\T$. Và $\Lambda x=0\iff Sx=\sqrt p(p\T x)\iff x_j=p\T x\ \forall j\iff x$ là hằng.
\end{proof}
\noindent\emph{Hệ quả cài đặt.} Mọi chuẩn $\norm\cdot_W$ và ràng buộc ở giới hạn cứng dùng $\Lambda$; không dùng căn bậc hai qua trị riêng, vì trị riêng xấp xỉ $0$ sinh hạng giả.

\begin{namedthm}{L8}{đạo hàm KD}
$\partial L_{KD}/\partial f=O\T e_O$ và $\partial^2L_{KD}/\partial f^2=O\T WO$, với $p_s$ của \emph{student}. Hessian dựng từ xác suất teacher là sai, trừ khi $p_s=p_t$.
\end{namedthm}
\begin{proof}
$L_{KD}$ của mẫu $i$ là $\frac1N\big(\operatorname{LSE}(f_{i,\mathrm{old}}/T)-p_{t,i}\T f_{i,\mathrm{old}}/T\big)$; gradient của LSE là softmax của student, Hessian là $D(p_s)-p_sp_s\T$; quy tắc dây chuyền cho thêm $1/T$ và $1/T^2$.
\end{proof}

\begin{namedthm}{L9}{$S_x$ có ridge}
Cho $S_0=\tilde X\T\tilde X/N$, $S_x=S_0+\rho_xI$ với $\rho_x>0$, và $M$ đối xứng với $\range(M)\subseteq\range(\tilde X\T)$. Gọi $\mathcal B$ là một cơ sở trực chuẩn của $\range(\tilde X\T)$. Khi đó:
(1) $S_x\succ0$ với mọi $d_a$, kể cả $d_a>N$;
(2) mọi eigvec suy rộng của $(M,S_x)$ với trị riêng khác $0$ nằm trong $\range(\tilde X\T)$;
(3) các trị riêng khác $0$ đó là trị riêng của bút chì thu hẹp $(\mathcal B\T M\mathcal B,\mathcal B\T S_0\mathcal B+\rho_xI)$, hội tụ khi $\rho_x\to0$;
(4) với $a$ $S_x$-chuẩn hóa trong $\range(\tilde X\T)$: $a\T S_0a\in[1-\rho_x/\lambda^+_{\min}(S_0),1]$.
\end{namedthm}
\begin{proof}
(1) $S_0\succeq0$, $\rho_x>0$. (2) Viết $a=a_r+a_n$, $a_r\in\range(\tilde X\T)$, $a_n\in\nullsp(\tilde X)$; $S_0$ và $I$ giữ cả hai không gian. Chiếu $Ma=\lambda S_xa$ lên $\nullsp(\tilde X)$ cho $0=\lambda\rho_xa_n$, nên $\lambda\ne0\Rightarrow a_n=0$. (3) Thay $a=\mathcal Bc$ và nhân $\mathcal B\T$; $\mathcal B\T S_0\mathcal B\succ0$ nên trị riêng liên tục tại $\rho_x=0$. (4) $a\T S_xa=1\Rightarrow a\T S_0a=1-\rho_x\norm a^2$, và $\norm a^2\le1/\lambda^+_{\min}(S_0)$ trên $\range(\tilde X\T)$.
\end{proof}

\begin{namedthm}{L10}{$L_3$ dạng đóng, chặt}
Đặt $c_3=2\sqrt2/(6\sqrt3)=\sqrt2/(3\sqrt3)\approx0.2722$. Với $L$ là CE trung bình cộng KD của PyCIL, với mọi $f$ và $h\in\R^n$:
\[
|D^3L(f)[h,h,h]|\le c_3\,\frac{1+T^{-3}}{N}\,\norm h^3 .
\]
Hằng số $c_3$ là chặt cho từng thành phần: đạt ở giới hạn một mẫu, hai lớp mang xác suất $q^*=(3-\sqrt3)/6$ và $1-q^*$, các lớp khác có xác suất $\to0$, $h=(1,-1,0,\dots)/\sqrt2$.
\end{namedthm}
\begin{proof}
Với $p=\operatorname{softmax}(x)$: $D^3\operatorname{LSE}(x)[h,h,h]=\mathbb E_p[(X-\mathbb E_pX)^3]$, trong đó $X$ nhận giá trị $h_j$ với xác suất $p_j$ (đạo hàm bậc ba của hàm sinh tích lũy là tích lũy bậc ba).

\emph{Cực đại của mô-men trung tâm bậc ba.} Gọi $Y=X-\mathbb EX$, nhận giá trị trong $[-\alpha,\beta]$ với $\alpha+\beta\le R:=\max_jh_j-\min_jh_j$. Bài toán $\max\mathbb EY^3$ trên mọi phân phối có $\mathbb EY=0$ là một bài toán mô-men tuyến tính với hai ràng buộc; theo định lý Richter--Rogosinski \cite{richter1957,rogosinski1958}, cực trị đạt tại phân phối có tối đa hai điểm. Với hai điểm $-s<0<t$ (xác suất $t/(s+t)$ tại $-s$): $\mathbb EY^2=st$, $\mathbb EY^3=st(t-s)$, $s+t\le R$. Biểu thức tăng khi phóng to $(s,t)$, nên lấy $s+t=R$, $s=Rx$: $\mathbb EY^3=R^3x(1-x)(1-2x)$, cực đại tại $x=(3-\sqrt3)/6$ với giá trị $R^3/(6\sqrt3)$. Đối xứng cho cực tiểu. Vậy $|\mathbb EY^3|\le R^3/(6\sqrt3)$.

$R\le\sqrt2\norm h$ vì $(h_i-h_j)^2\le2(h_i^2+h_j^2)$, nên $|D^3\operatorname{LSE}[h,h,h]|\le c_3\norm h^3$. CE của mẫu $i$ là $\operatorname{LSE}(f_i)-f_{i,y_i}$ chia $N$; KD của mẫu $i$ là $(\operatorname{LSE}(f_{i,\mathrm{old}}/T)-p_{t,i}\T f_{i,\mathrm{old}}/T)/N$, có thêm hệ số $T^{-3}$ và chỉ phụ thuộc $h_{i,\mathrm{old}}$. Cộng theo mẫu và dùng $\sum_i\norm{h_i}^3\le(\sum_i\norm{h_i}^2)^{3/2}$.

\emph{Đạt cận.} Với $h=(1,-1,0,\dots)/\sqrt2$, $\norm h=1$, $R=\sqrt2$ và phân phối của $X$ dồn về $\pm1/\sqrt2$ với xác suất $q^*$, $1-q^*$: cả hai bất đẳng thức thành đẳng thức ở giới hạn.
\end{proof}
\begin{hequa}
A6$'$ không còn là giả thiết. Cận Taylor của bước nhảy $j_u$ khi bỏ một slot sống có hằng số tính được (\Tref{T7}(i.2)), dùng để lọc ứng viên ở bước S6(b).
\end{hequa}
\begin{kiem}
T0a--T0f (L5--L8, sai số $\le2.5\cdot10^{-11}$; T0c-NEG loại mẫu số $p_t$, T0e-NEG loại Hessian teacher lệch 0.69); T14c, T14c-WIT (L9); T12a, T12a2 (L10, cấu hình hai điểm đạt $1.0000$ cận; NEG loại hằng số ${\times}0.1$ và chuẩn hóa $T^{-1}$).
\end{kiem}

\section{Hiệu chỉnh: dạng đóng, bảo vệ lớp cũ và chứng chỉ}\label{sec:core}

Mục này chứng minh rằng hiệu chỉnh của bước S3 được xác định duy nhất, tách được thành phần sửa và phần học mới, và có chứng chỉ không làm tăng KD ở bậc một.

\subsection{T1: dạng đóng và phân rã}

\begin{namedthm}{T1}{dạng đóng, phân rã, các dạng không cần range}
Giả thiết A2, A3, $\beta\ge0$.
\begin{enumerate}[label=(\alph*)]
\item Bài toán $(P_\beta)$: $\min_d\Phi(d)+\beta\norm{C(d-\drep)}_W^2$ có nghiệm duy nhất $d^*(\beta)=\drep+\dN$.
\item $\Phi(\drep+x)=\Phi_N(x)+\mathrm{const}$; do đó $(P_\beta)$ tương đương $\min_x\Phi_N(x)+\beta x\T Bx$, nghiệm $x=\dN$.
\item \emph{Dạng output-space.} Với $\hat J=[J;C]$, $\hat H=\blkdiag(H,\beta W)=\hat R\T\hat R$ (một thừa số bất kỳ, ví dụ $\hat R=\blkdiag(H^{1/2},\sqrt\beta\Lambda)$), $\hat t=[u;C\drep]$, $\hat M=\hat R\hat J\hat J\T\hat R\T$: $d^*=\hat J\T\hat R\T(\hat M+\mu)^{-1}\hat R\hat t$. Không có pseudo-inverse.
\item \emph{$z$ không cần range.} Với mọi thừa số $H=\Lambda_h\T\Lambda_h$: $z(r)=\Lambda_h\T(\Lambda_hK\Lambda_h\T+\mu)^{-1}\Lambda_h(Jg_r)$, $g_r=J\T r$.
\end{enumerate}
\end{namedthm}
\begin{proof}
(a) Hessian của mục tiêu là $2A_\beta\succ0$, nên nghiệm duy nhất. Điều kiện dừng: $A_\beta d=J\T Hu+\beta B\drep$. Thay $d=\drep+x$ và dùng $(F+\mu)\drep=J\T Hu_R$, được $A_\beta x=J\T Hu_N$.
(b) Khai triển; số hạng chéo là $-2x\T[J\T Hu_R-(F+\mu)\drep]=0$.
(c) Mục tiêu là $\norm{\hat R(\hat Jd-\hat t)}^2+\mu\norm d^2$, nên $d=(\hat J\T\hat H\hat J+\mu)^{-1}\hat J\T\hat R\T\hat R\hat t$; push-through với $X=\hat R\hat J$.
(d) $z(r)=\Lambda_h\T\Lambda_hJ(J\T\Lambda_h\T\Lambda_hJ+\mu)^{-1}J\T r$; push-through với $X=\Lambda_hJ$.
\end{proof}
\begin{hequa}
Bước S3 tính $\drep$ và $\dN$ bằng hai hệ tuyến tính riêng; phân rã (b) cho phép áp $\beta$ chỉ lên phần học mới mà không đụng phần sửa. Dạng (c, d) thay công thức cũ $RM(M+\mu)^{-1}R^+e$, vốn sai khi softmax bão hòa vì ngưỡng trị riêng của $R^+$.
\end{hequa}
\begin{kiem}
T1a $3.7\cdot10^{-12}$ (T1a-NEG: bỏ $\beta$ lệch $6\cdot10^{-2}$); T1b $4.0\cdot10^{-12}$; T1c $1.4\cdot10^{-11}$ ở logit scale 1 đến 30 (T1c-NEG: công thức pinv có ngưỡng lệch 0.81 ở scale $\ge20$).
\end{kiem}

\subsection{T2: trường hợp $\beta=0$, giới hạn cứng và bổ đề chiếu lồng nhau}

\begin{namedthm}{T2}{rút gọn, giới hạn $\beta\to\infty$, chiếu lồng nhau}
Giả thiết A2, A3.
\begin{enumerate}[label=(\alph*)]
\item $\beta=0$: $d^*(0)=GJ\T HK_Ez$ (VNA thuần). Ngoài ra $g\T d^*(0)=V\ge0$ và $d^*(0)=Qg$ với $Q=L\T K_EL\succeq0$, $L=HJG$.
\item \emph{Giới hạn cứng, $C$ hạng tùy ý.} Với $A:=A_0$, $b:=J\T Hu_N$, $\hat C=\Lambda C$: khi $\beta\to\infty$,
\[
\dN\to d_N^\infty:=\argmin\{\Phi_N(x):\hat Cx=0\}=A^{-1}b-A^{-1}\hat C\T(\hat CA^{-1}\hat C\T)^+\hat CA^{-1}b,
\]
với $\norm{\dN-d_N^\infty}_A\le\norm{\Pi A^{-1/2}b}/(1+\beta\sigma_{\min}^2)$, trong đó $\sigma_{\min}$ là giá trị suy biến khác $0$ nhỏ nhất của $\hat CA^{-1/2}$ và $\Pi$ là phép chiếu lên không gian hàng của nó. Nếu $\hat C$ có hạng cột đầy thì $d_N^\infty=0$.
\item Ràng buộc ở giới hạn là $\Lambda Cx=0$, \emph{không phải} $Cx=0$: mỗi mẫu vẫn được dịch đều toàn bộ logit cũ, và KD bất biến với phép dịch đó.
\item \emph{Chiếu lồng nhau ở giới hạn cứng.} Với $P=I-\hat C^+\hat C$: $Jd_N^\infty\in\range(K_{NS})$, $K_{NS}:=JPJ\T=K-J\hat C\T(\hat C\hat C\T)^+\hat CJ\T$ (phần bù Schur theo khối $[J;\hat C]$). Do đó vận tốc học mới nằm trong $\range(K_{NS})\subseteq\range(K)$.
\end{enumerate}
\end{namedthm}
\begin{proof}
(a) \Tref{L4} cho $g\T GJ\T H=z\T$, nên $g\T d=z\T K_Ez$; viết $d=GJ\T HK_EHJGg=L\T K_ELg$.
(b) Đổi biến $y=A^{1/2}x$: $\min\norm{y-A^{-1/2}b}^2+\beta\norm{\hat CA^{-1/2}y}^2$. Lấy SVD rút gọn $\hat CA^{-1/2}=U\Sigma V\T$; phần của $y$ ngoài $\range(V)$ không phụ thuộc $\beta$; tọa độ thứ $i$ trong range là $b_i/(1+\beta\sigma_i^2)\to0$. Giới hạn là nghiệm KKT, đúng với $(\cdot)^+$ vì $\hat CA^{-1}b\in\range(\hat CA^{-1}\hat C\T)$.
(c) $\hat Cx=0\iff Cx\in\nullsp(\Lambda)=\nullsp(W)$ (\Tref{L7}).
(d) $\hat Cd_N^\infty=0$ nên $d_N^\infty=Pd_N^\infty$ và $Jd_N^\infty\in\range(JP)=\range(JPJ\T)$.
\end{proof}
\begin{hequa}
(a) là cơ sở của chứng chỉ ở bước S4 khi $\beta=0$ (B5 và nhánh $\beta=0$ của CIL): bước VNA thuần luôn là hướng giảm ở bậc một. (b, c) cho biết $\beta$ lớn tới đâu thì vẫn còn chỗ học: đúng ràng buộc là $\Lambda C$, và giới hạn cứng chỉ đóng băng hoàn toàn khi $\hat C$ có hạng cột đầy (không xảy ra với ResNet32, vì $p\gg N(c_{\mathrm{old}}-1)$). (d) là \emph{bổ đề chiếu lồng nhau}: phép chiếu không tạo hướng mới ngoài $\range(K)$; đây là lý do cần GIR.
\end{hequa}
\begin{kiem}
T2a $5.2\cdot10^{-16}$; T2b $5.5\cdot10^{-13}$; T2c ($Q\succeq0$); T2d--e (tốc độ $O(1/\beta)$ lệch $\le1.8\%$); T2f (T2f-NEG: ràng buộc $Cx=0$ cách $\ge0.265$); T2g ($C$ hạng thiếu); T2h, T2i (Schur, $1.6\cdot10^{-14}$); T11f (hạng cột đầy: $\norm{\dN}/\norm{d_N^0}=1.3\cdot10^{-8}$ tại $\beta=10^{10}$).
\end{kiem}

\subsection{T3: sửa, rò rỉ và chứng chỉ không quên thêm}

\begin{namedthm}{T3}{sửa, rò rỉ, nhân tử KKT}
Giả thiết A2, A3; có lớp cũ (CIL).
\begin{enumerate}[label=(\alph*)]
\item \emph{Phần sửa đúng chiều:} $g_O\T\drep=V_O\ge0$.
\item $h(\beta)=V_O+e_O\T C\dN$ và $|e_O\T C\dN|\le\norm{e_O}_{W^+}\sqrt{\leak(\beta)}$.
\item \emph{Rò rỉ:} $\leak'(\beta)=-2(B\dN)\T G_\beta(B\dN)\le0$ và $\leak(\beta)\le\min(\leak(0),\norm{u_N}_H^2/\beta)$; $\leak$ giảm chặt khi còn dương và $\to0$.
\item \emph{Mức quên bậc hai.} Với bước $\theta-\varepsilon d$: $L_{KD}(\theta-\varepsilon d)-L_{KD}(\theta)=-\varepsilon g_O\T d+\tfrac12\varepsilon^2[\norm{Cd}_W^2+\eO\T\nabla^2f[d,d]]+O(\varepsilon^3)$.
\item \emph{Chứng chỉ.} Với $V_O>0$: nếu $\leak(\beta)\le\varepsilon^*$ thì $h(\beta)\ge0$. Tập $\{\beta:\leak(\beta)\le\varepsilon^*\}$ là khoảng $[\beta_{\mathrm{cert}},\infty)$, $\beta_{\mathrm{cert}}\le\beta_{\mathrm{suf}}$, và $\beta_{\mathrm{cert}}=0\iff\leak(0)\le\varepsilon^*$.
\item \emph{Nhân tử KKT.} Với $V_O>0$, $d^*(\beta_{\mathrm{cert}})$ là nghiệm duy nhất của $\min\Phi(d)$ với $\norm{C(d-\drep)}_W^2\le\varepsilon^*$, và $\beta_{\mathrm{cert}}$ là nhân tử Lagrange duy nhất (đối ngẫu mạnh). Ngân sách $\varepsilon^*$ không có tham số: đó là ngân sách lớn nhất mà Cauchy--Schwarz ở (b) \emph{đủ} để bảo đảm $h\ge0$ (điều kiện đủ, có tính bảo thủ).
\item \emph{Biên và suy biến.} (i) Nếu $\beta_{\mathrm{cert}}>0$ thì $0\le h(\beta_{\mathrm{cert}})\le2V_O$. (ii) $h(\infty)=V_O$. (iii) Nếu $V_O=0$ và $e_O\ne0$ thì $h(0)=0$ chính xác, nên $\beta=0$ có chứng chỉ. (iv) Nếu $e_O=0$ thì $h\equiv0$ và $\varepsilon^*=0/0$ không được tính; quy tắc là $\beta=0$.
\item $h$ \emph{không} đơn điệu: $h(\beta)=V_O+\sum_ic_i/(1+\beta\lambda_i)$ trong cơ sở riêng suy rộng của $(A_0,B)$ với $c_i$ có thể trái dấu. Bisection trên $h$ không hợp lệ; trên $\leak$ thì hợp lệ.
\item \emph{Phạt về đích $0$ làm mất phần sửa:} với phạt $\beta\norm{Cd}_W^2$ trên toàn bộ $u$:
\[|g_O\T d_\beta|\le\norm{e_O}_{W^+}\norm u_H/\sqrt\beta\to0.\]
\end{enumerate}
\end{namedthm}
\begin{proof}
(a) \Tref{L4} với $r=\eO$: $g_O\T\drep=z_O\T K_Ez_O$.
(b) $g_O\T x=e_O\T Cx$; theo \Tref{L5}, $e_O=\Lambda\T w$ với $\norm w=\norm{e_O}_{W^+}$ khi chọn $w$ nhỏ nhất; nên $|e_O\T Cx|\le\norm w\norm{\Lambda Cx}$.
(c) $dG_\beta/d\beta=-G_\beta BG_\beta$, nên $(\dN)'=-G_\beta B\dN$ và $\leak'=2d\T Bd'$. Cận $1/\beta$: $\beta\leak(\beta)\le\Phi_N(\dN)+\beta\leak(\beta)\le\Phi_N(0)=\norm{u_N}_H^2$.
(d) Taylor của $L_{KD}\circ f$: gradient $g_O$, Hessian $J\T O\T WOJ+\sum_j(\eO)_j\nabla^2f_j$.
(e) $h\ge V_O-\norm{e_O}_{W^+}\sqrt{\leak}\ge V_O-\norm{e_O}_{W^+}\sqrt{\varepsilon^*}=0$; tập là khoảng vì $\leak$ không tăng và liên tục; $\beta_{\mathrm{suf}}$ thuộc tập theo (c).
(f) Theo \Tref{T1}(b) bài toán là $\min\Phi_N(x)$ với $x\T Bx\le\varepsilon^*$. Slater thỏa tại $x=0$; mục tiêu lồi chặt, nên nghiệm duy nhất và đối ngẫu mạnh. KKT: $x=\dN$, $\beta\ge0$, $\beta(\leak(\beta)-\varepsilon^*)=0$; theo (c) nghiệm là $\beta_{\mathrm{cert}}$.
(g) (i) Tại $\beta_{\mathrm{cert}}>0$, $\leak=\varepsilon^*$ nên $|e_O\T Cd|\le V_O$. (ii) $\Lambda Cd_N^\infty=0$ và $e_O=\Lambda\T w$. (iii) $V_O=\norm{J_E\T z_O}^2=0\Rightarrow u_R=0$ và $h(0)=z_O\T K_Ez_N=(J_E\T z_O)\T J_E\T z_N=0$. (iv) $g_O=C\T e_O=0$.
(h) Chéo hóa đồng thời $(A_0,B)$; phản ví dụ số ở 62/300 bài.
(i) Như (b), (c) với $u$ thay $u_N$.
\end{proof}
\begin{hequa}
Đây là cơ sở của bước S3 trong CIL. (a) nói phần sửa luôn kéo KD về đúng chiều. (e, f) cho quy tắc chọn $\beta$: tìm đầu mút trên của khoảng có $\leak\le\varepsilon^*$ bằng bisection trên $\leak$ (không trên $h$, theo (h)). (g)(iii, iv) cho hai nhánh suy biến của quy tắc $\beta$: dùng $\beta=0$, không bỏ qua. (i) giải thích vì sao ràng buộc phải đặt về đích $C\drep$ chứ không về $0$.
\end{hequa}
\begin{kiem}
T3a $2.2\cdot10^{-13}$; T3b (T3b-NEG: dùng $\norm{e_O}_W$ vi phạm); T3c (0 vi phạm); T3e (0/6900 điểm lưới; T3e-WIT: VNA thuần làm KD tăng ở 102/300 bài, nên chứng chỉ không rỗng); T3i (SLSQP $10^{-7}$); T3d-WIT ($h$ không đơn điệu 62/300); T11a--g.
\end{kiem}

\subsection{T4: cái giá về plasticity của $\beta$}

\begin{namedthm}{T4}{plasticity}
Giả thiết A2, A3.
(a) $\beta\mapsto\Phi_N(\dN)$ không giảm.
(b) $0\le\Phi_N(\dN)-\Phi_N(d_N^0)\le\beta(\leak(0)-\leak(\beta))\le\beta\leak(0)$.
(c) Riêng độ khớp $\norm{J\dN-u_N}_H^2$ \emph{không} nhất thiết đơn điệu.
\end{namedthm}
\begin{proof}
(a) Điều kiện dừng $\nabla\Phi_N(x_\beta)=-2\beta Bx_\beta$ và $x_\beta'=-G_\beta Bx_\beta$ cho $\frac{d}{d\beta}\Phi_N(x_\beta)=2\beta(Bx_\beta)\T G_\beta(Bx_\beta)\ge0$.
(b) $x_\beta$ cực tiểu $\Phi_N+\beta\leak$, nên $\Phi_N(x_\beta)+\beta\leak(\beta)\le\Phi_N(x_0)+\beta\leak(0)$.
(c) Phản ví dụ: $F=\begin{psmallmatrix}2.077&.563\\.563&.265\end{psmallmatrix}$, $B=\begin{psmallmatrix}12.70&-.027\\-.027&10^{-4}\end{psmallmatrix}$, $J\T Hu_N=(.278,.320)$, $\mu=1$: độ khớp (bỏ hằng số) bằng $-0.1434,-0.1443,-0.1451$ tại $\beta=0,0.05,0.2$ (giảm), trong khi $\Phi_N$ tăng.
\end{proof}
\begin{hequa}
$\beta_{\mathrm{cert}}$ là $\beta$ \emph{nhỏ nhất} có chứng chỉ, nên theo (a, b) nó mất ít plasticity nhất trong các lựa chọn an toàn. Đây là lý do không dùng $\beta_{\mathrm{suf}}$ (trung vị lớn gấp 29 lần) làm mặc định.
\end{hequa}

\subsection{T5: giảm loss và quy tắc bỏ qua}

\begin{namedthm}{T5}{giảm loss, bỏ qua, không có núm $\gamma$}
Giả thiết A2, A3; (e), (f) thêm A6.
\begin{enumerate}[label=(\alph*)]
\item $\norm{\drep}\le\norm{u_R}_H/\sqrt\mu$, $\norm{\dN}\le\norm{u_N}_H/\sqrt\mu$ với mọi $\beta$.
\item $g\T d^*(\beta)\ge V-\norm g_G\min\big(\sqrt{\beta\leak(0)},2\norm{u_N}_H\big)$.
\item \emph{Không có $\gamma$.} Với $u_\gamma=K_E(z_N+\gamma z_O)$, $\beta=0$, và $z_N,z_O$ độc lập tuyến tính: $g\T d\ge0$ với mọi $K_E\succeq0$ khi và chỉ khi $\gamma=1$ (tức VNA hiện tại).
\item Nếu $g\T\hat d\ge0$ thì với mọi $\lambda\ge0$, $p=g+\lambda\hat d$ thỏa $g\T p\ge\norm g^2$.
\item Với $\rho:=\lambda\norm{\hat d}/\norm g$ và $g\T\hat d\ge0$: $L(\theta-\varepsilon p)\le L(\theta)-\varepsilon\norm g^2(1-\tfrac12L_s\varepsilon(1+\rho)^2)$; $L$ giảm chặt khi $0<\varepsilon<2/(L_s(1+\rho)^2)$.
\item Với bước baseline $\delta$ bất kỳ (kể cả momentum): $L(\theta+\delta-\varepsilon\lambda\hat d)-L(\theta+\delta)\le-\varepsilon\lambda g\T\hat d+\varepsilon\lambda L_s\norm\delta\norm{\hat d}+\tfrac12L_s\varepsilon^2\lambda^2\norm{\hat d}^2$, và tương tự cho $L_{KD}$ với $\hat h$, $L_K$. Cải thiện chặt khi $\varepsilon$ nhỏ cần $\hat g>0$ (và $\hat h>0$).
\item Không có bảo đảm vô điều kiện khi $\beta>0$: $g\T d^*(\beta_{\mathrm{cert}})<0$ có thể xảy ra.
\end{enumerate}
\end{namedthm}
\begin{proof}
(a) $\norm{G_\beta J\T Hu}^2\le\norm{G_\beta}\norm{J\T Hu}^2_{G_\beta}\le\mu^{-1}\norm{J\T Hu}_G^2\le\mu^{-1}\norm u_H^2$ (\Tref{L2}, \Tref{L1}).
(b) $g\T(\drep+d_N^0)=V$; $d_N^0-\dN=\beta G_\beta Bd_N^0$; $|g\T\beta G_\beta Bd_N^0|\le\norm g_{G_\beta}\sqrt\beta\norm{\Lambda Cd_N^0}$ (\Tref{L3}); cận thứ hai từ \Tref{L1}, \Tref{L2}.
(c) $g\T d=z\T K_E(z_N+\gamma z_O)$. Với $\gamma=1$ bằng $z\T K_Ez\ge0$. Với $K_E=vv\T$, $g\T d=(v\T z)(v\T(z_N+\gamma z_O))$, âm với một $v$ nào đó trừ khi $z_N+\gamma z_O=cz$, $c\ge0$, buộc $\gamma=1$.
(d) $g\T p=\norm g^2+\lambda g\T\hat d$.
(e) Bất đẳng thức descent $L$-smooth với $g\T p\ge\norm g^2$, $\norm p\le(1+\rho)\norm g$.
(f) $L$-smooth tại $\theta+\delta$ và $\norm{\nabla L(\theta+\delta)-g}\le L_s\norm\delta$.
(g) Phản ví dụ số (8/200 bài).
\end{proof}
\begin{hequa}
Đây là cơ sở của quy tắc bỏ qua ở bước S4: khi chứng chỉ không âm, thêm $\lambda\hat d$ không làm xấu bước của optimizer gốc ở bậc một, với \emph{mọi} $\lambda\ge0$ (d, f); khi âm thì bỏ qua (g). (c) loại hẳn núm $\gamma$: trọng số giữa phần sửa và phần học mới không được tự do.
\end{hequa}
\begin{kiem}
T5a, T5b; T5c-WIT ($g\T d^*<0$ ở 8/200 tại $\beta_{\mathrm{cert}}$); T5d (bổ đề descent) và T5d-NEG; T5e (tìm kiếm đối kháng không tìm được $g\T d<0$ khi $\gamma=1$, min $9.7\cdot10^{-18}$; NEG tìm được với $\gamma\in\{0,2,5\}$); T15a, T15b (mức thuật toán trên MLP thật: $L$ giảm ở mọi lần áp dụng; KD không tăng khi $\hat h\ge0$, 0/24).
\end{kiem}

\section{Chọn neuron ảo và GIR}\label{sec:select}

\subsection{T6: chọn hướng vào $A$}

\begin{namedthm}{T6}{chọn neuron ảo}
Giả thiết A2, A3, A4$'$ (mọi $a_j$ trong vùng mặt nạ cố định $R_D$), $\beta$ cố định, $S_x\succ0$ (\Tref{L9}). Đặt $y_\beta:=HJG_\beta g$ và $M_S(\beta):=\sym(N_zN_{z_O}\T)+\sym(N_{y_\beta}N_{z_N}\T)$.
\begin{enumerate}[label=(\alph*)]
\item Không có BN: $g\T d^*(\beta)=\tr(A\T M_S(\beta)A)$. Có BN: $g\T d^*(\beta)=\sum_ja_j\T M_Sa_j/\sigma(a_j)^2$, trở lại dạng vết khi $\sigma(a_j)=1$.
\item Trên các tập $S_x$-trực chuẩn tối đa $r$ neuron (trong $R_D$), giá trị lớn nhất là tổng các trị riêng \emph{dương} lớn nhất (tối đa $r$) của $(M_S,S_x)$, đạt tại eigvec tương ứng nếu chúng nằm trong $R_D$; nếu không, đó chỉ là cận trên.
\item Với $\beta=0$: $M_S=N_zN_z\T$, nên tiny\_V$(z)$ \emph{chính xác} là tối ưu.
\item Với $\beta>0$, đặt $\Delta=\sym(N_{y_\beta-z}N_{z_N}\T)$: độ dưới tối ưu của tiny\_V$(z)$ không quá
\[2\textstyle\sum_{j\le r}|\lambda_j(S_x^{-1/2}\Delta S_x^{-1/2})|;\]
ngoài ra $y_\beta-z=-\beta HJG_\beta BGg$.
\item Quy tắc xen kẽ $A_1=\argmax M_S(\beta_0)$ tối ưu có điều kiện theo $\beta_0$; không có khẳng định tối ưu đồng thời theo $(A,\beta)$. Mọi chứng chỉ đúng với \emph{mọi} $A$.
\end{enumerate}
\end{namedthm}
\begin{proof}
(a) \Tref{L4} cho $g\T\drep=z\T K_Ez_O$ và $g\T\dN=y_\beta\T K_Ez_N$; với $s\T K_Et=\sum_ja_j\T N_sN_t\T a_j$ (tuyến tính trong vùng mặt nạ), lấy phần đối xứng. (b) Đổi biến $U=S_x^{1/2}A$ và áp Ky Fan \cite{fan1949}; chỉ giữ trị riêng dương. (c) $y_0=z$ và $\sym(N_zN_{z_O}\T)+\sym(N_zN_{z_N}\T)=N_zN_z\T$. (d) $M_S=N_zN_z\T+\Delta$; $A_t$ tối ưu cho $N_zN_z\T$; phần $\Delta$ chặn bằng Ky Fan cho trị tuyệt đối; $y_\beta-z=HJ(G_\beta-G)g$. (e) Theo định nghĩa.
\end{proof}
\begin{hequa}
Bước S2: làm mới $A$ ở đầu task và sau mỗi 10 lần áp dụng; $A_0=\text{tiny\_V}(z)$, rồi (khi $\beta_0>0$) $A$ là các eigvec suy rộng dương của $(M_S(\beta_0),S_x)$. Tối ưu chỉ được khẳng định trong một vùng mặt nạ; chứng chỉ không phụ thuộc lựa chọn này.
\end{hequa}
\begin{kiem}
T6a $4.2\cdot10^{-12}$ (T6a-NEG: dùng $N_zN_z\T$ khi $\beta>0$ lệch $\ge2.1\cdot10^{-3}$); T6b; T6c; T6d (T6d-WIT: dưới tối ưu tới 91\%). Quy tắc xen kẽ tăng $g\T\hat d$ trung vị 22\% (\st{empirical}).
\end{kiem}

\begin{namedthm}{T14}{chọn hướng khi có BN; ridge}
Neuron là một kênh mới, BN train mode theo sau ($\gamma=1$, $\beta_{BN}=0$), có thể có ReLU sau BN. Cho $s\in\R^{N\times q}$ và mặt nạ $D$.
(a) Với mọi $a$ ($\tilde Xa\ne0$): $\norm{\sum_i\varphi_i(a)s_i}^2=N\,a\T\tilde X\T Dss\T D\tilde Xa/a\T\tilde X\T\tilde Xa$, một thương Rayleigh suy rộng với thống kê \emph{đã trừ trung bình}, bất biến theo tỉ lệ của $a$.
(b) Không có ReLU: giá trị lớn nhất là $N\lambda_{\max}(\tilde X\T ss\T\tilde X,\tilde X\T\tilde X)$; với $r$ neuron $S_0$-trực giao là $N$ lần tổng $r$ trị riêng đầu.
(c) Có ReLU: trong $R_D$ giá trị không vượt $N\lambda_{\max}(\tilde X\T Dss\T D\tilde X,\tilde X\T\tilde X)$.
(d) Dùng thống kê không trừ trung bình là sai khi $\bar x\ne0$.
Với $\varepsilon_{BN}>0$, (a) có sai số tương đối $\varepsilon_{BN}/(\sigma^2+\varepsilon_{BN})$; hướng tối ưu không đổi.
\end{namedthm}
\begin{proof}
(a) BN train mode cho $\varphi=D\tilde Xa/\sigma(a)$; thay vào $\varphi\T ss\T\varphi$. (b) Rayleigh--Ritz trên $\range(\tilde X\T)$ (\Tref{L9}) và Ky Fan. (c) Trên $R_D$, $D$ cố định; áp (b) với $D\tilde X$ ở tử. (d) Phản ví dụ số.
\end{proof}
\begin{hequa}
Bước S2 và S6 khi slot có BN theo sau: chuẩn hóa hướng theo $S_0$ đã trừ trung bình và giải bài toán Rayleigh suy rộng tương ứng; dùng thống kê không trừ trung bình là sai.
\end{hequa}
\begin{kiem}
T14a $2.0\cdot10^{-14}$; T14b; T14b-NEG (không trừ trung bình sai $\ge1.7\cdot10^{-4}$); T17d (phạm vi $\varepsilon_{BN}$).
\end{kiem}

\subsection{T7: GIR, đặt neuron vào slot}

Gọi $r\in\range(H)$ là residual mà GIR nhắm tới: $r=\eO$ trong CIL (ưu tiên giữ hàm cũ), $r=e$ trong B5. Đặt $t_r=-H^+r$, và cho mỗi neuron mới $v$: $s_v=J_v\T\rperp$, $S_v=J_v\T H_\perp J_v+\mu I$. Lợi dài hạn của việc thêm cột $J_v$ là $\mathrm{gain}(v\mid J):=q^*_r(J)-q^*_r([J,J_v])$.

\begin{namedthm}{T7}{GIR}
Giả thiết A2, A3, A12. (f)--(h) cần thêm A5 tại tham số lúc chèn; ý ``tiny\_V'' của (d) cần A4$'$; việc \emph{diễn giải} gain như mức thay đổi loss về sau cần A10.
\begin{enumerate}[label=(\alph*)]
\item \emph{Hàm chiếu chung (cấu trúc).} $\mathrm{gain}(v\mid J)=\tfrac12[E(J,t_r)-E([J,J_v],t_r)]$, còn VNA ở $\beta=0$ là nghiệm của $E(J,u)$. Hai bài toán có cùng dạng nhưng \emph{không cùng thang đo}: $u$ là một vận tốc, $t_r$ là một độ dời Newton; VNA tìm hệ số trên span cố định, swap mở rộng span.
\item \emph{Dạng Schur, mọi $a_v,b_v$, mọi phi tuyến.} $\mathrm{gain}=\tfrac12s_v\T S_v^{-1}s_v$, với $\rperp=r-z(r)=\mu(HK+\mu)^{-1}r=H_\perp H^+r$ và $H_\perp=\mu R(M+\mu)^{-1}R\succeq0$. Hệ quả: $\mathrm{gain}>0\iff J_v\T\rperp\ne0$. Không cần A4$'$.
\item Tại cực tiểu của mạng sống $x^*=-GJ\T r$, gradient theo $B_v$ bằng $s_v$.
\item \emph{Chọn $v$.} $\tfrac12\norm{s_v}^2/(\lambda_{\max}(J_v\T HJ_v)+\mu)\le\mathrm{gain}\le\tfrac12\norm{s_v}^2/\mu$. Cận dưới $\ell(v)$ không cần giải hệ và là \emph{chứng chỉ} cho gain. Với $\hat v=\argmax_{v\in\Vcal}\ell(v)$ trên một tập ứng viên hữu hạn: $\mathrm{gain}(\hat v)\ge\frac{\mu}{c_H+\mu}\max_{v\in\Vcal}\mathrm{gain}(v)$, $c_H=\max_v\lambda_{\max}(J_v\T HJ_v)$. Dưới A4$'$, hướng tốt nhất theo $\norm{s_v}$ là tiny\_V với cotangent $\rperp$ (có BN: \Tref{T14}). Vì $\varphi(-a)\ne-\varphi(a)$ với ReLU, $\Vcal$ gồm $\{\sigma a_k+b:\sigma=\pm1,b\in\{0,-\operatorname{median}(\sigma Xa_k)\}\}$, bỏ ứng viên có $\varphi\equiv0$.
\item $\rperp$ là phần bù \emph{xiên} theo metric $H$; nó trùng phần dư Euclid $(I-JJ^+)r$ chỉ khi $H\propto I$. Tương ứng với bước chọn atom của OMP chỉ chính xác khi $H=I$, $\mu\to0$.
\item \emph{Slot chết} (A5 tại tham số lúc chèn): bỏ $u$ và chèn $v$ với $B_v=0$ thì hàm trên cửa sổ không đổi \emph{chính xác}, giá bỏ $u$ bằng $0$, không rò rỉ sang logit cũ tại lúc chèn, và các cột của $a_v$ bằng $0$. A5 phải kiểm lại sau bước hiệu chỉnh: một đơn vị chết tại $\theta$ có thể sống tại $\theta'$.
\item \emph{Swap không giao hoán với VNA.} Sau khi chèn, $K'=K+J_vJ_v\T$; vì vậy $\hat d$ được tính và áp dụng \emph{trước} swap. $\hat d$ không đụng tham số của slot chết: $(\hat d)_u=0$.
\item \emph{Quyết định swap slot chết không cần chân trời:} lợi ròng bằng $\mathrm{gain}(v\mid J)\ge0$, $>0\iff s_v\ne0$. Với $r=\eO$, ``gain = mức giảm KD về sau'' không suy ra từ A10 (A10 là hồi phục trên $L$); chỉ là ghi chú phạm vi vì nhánh này không kích hoạt trong CIL.
\item \emph{Slot sống: đo, không chứng nhận.} (i.1) $\ell(\text{sau})-\ell(\text{trước})=j_u^\ell$ chính xác bằng một forward (\st{measured}). (i.2) $|j_u^\ell-(-r\T f_u+\tfrac12f_u\T H_\ell f_u)|\le(L_3/6)\norm{f_u}^3$, $L_3\le c_3(1+T^{-3})/N$ (\Tref{L10}). (i.3) Dự báo nhất quán A10 là $\Delta L_{\mathrm{pred}},\Delta KD_{\mathrm{pred}}$ dọc bước hồi phục trên $L$, mọi đại lượng sau swap tại $f-f_u$ (\Tref{T16}). (i.5) Nếu $j_u^L\le0$ và $j_u^{KD}\le0$ thì tại lúc chèn $L$ và $L_{KD}$ trên cửa sổ không tăng; đây là điều \emph{duy nhất} được chứng nhận cho slot sống.
\item \emph{GIR trong CIL} (\st{empirical}): tỉ lệ đơn vị chết đo được trên ResNet32 có BN là $0$; với cấu hình mặc định GIR không hoạt động trong CIL.
\end{enumerate}
\end{namedthm}
\begin{proof}
(a) Vì $r\in\range(H)$, $HH^+r=r$, nên $\norm{\Jcal x-t_r}_H^2=x\T\Jcal\T H\Jcal x+2x\T\Jcal\T r+r\T H^+r$, suy ra $E(\Jcal,t_r)=2q^*(\Jcal)+r\T H^+r$; lấy hiệu. Phần VNA là \Tref{T1}(a) ở $\beta=0$.

(b) Nghịch đảo khối cho $[J,J_v]$: $\tilde b\T\tilde A^{-1}\tilde b=b\T A^{-1}b+(b_v-J_v\T HJGb)\T S_v^{-1}(\cdot)$ với $b=J\T r$, $b_v=J_v\T r$; và $b_v-J_v\T HJGJ\T r=J_v\T(r-z(r))$. Theo \Tref{L1}, $H_\perp=R(I-M(M+\mu)^{-1})R=\mu R(M+\mu)^{-1}R$. Đồng nhất $r-z(r)=\mu(HK+\mu)^{-1}r$ suy từ $z(r)=HK(HK+\mu)^{-1}r$ (push-through).

(c) $\nabla_{B_v}q=J_v\T(r+HJx^*)=J_v\T(r-z(r))$.

(d) $\mu I\preceq S_v\preceq J_v\T HJ_v+\mu I$; do đó $\mathrm{gain}(\hat v)\ge\ell(\hat v)\ge\ell(v^*)\ge\frac{\mu}{c_H+\mu}\mathrm{gain}(v^*)$.

(e) $HK$ đồng dạng trên range với $RKR\succeq0$; $\mu(HK+\mu)^{-1}$ tiến tới phép chiếu lên $\nullsp(HK)$ dọc $\range(HK)=H\range(J)$; nếu $H=I$ thì phép chiếu trực giao.

(f) Đóng góp của $v$ là $h_v(x)B_v\T=0$; đạo hàm theo $a_v$ tỉ lệ với $B_v$ nên bằng $0$; cột của $u$ bằng $0$ trên cửa sổ. Nếu A5 không đúng tại $\theta'$ thì hàm đổi đúng một lượng $-f_u(\theta')$.

(g) Cột mới là $J_v$; với cột $u$ bằng $0$, $F$ có hàng và cột $u$ bằng $0$, nên $G_{uu}=\mu^{-1}I$, các khối chéo bằng $0$, và $(J\T Hy)_u=0$.

(h) Từ (b), (f). (i.1) Sau swap output là $f-f_u$. (i.2) Taylor bậc ba với phần dư Lagrange và \Tref{L10}. (i.3) Xem \Tref{T16}. (i.5) Từ (i.1).
\end{proof}
\begin{hequa}
Bước S6(a): sau bước S5, đếm lại slot chết; với mỗi slot còn chết (tối đa $S$), chọn $v=\argmax\ell(v)$ trên tập ứng viên dựng từ eigvec của $(N_{\rperp}N_{\rperp}\T,S_x)$, đặt $B_v=0$. Lợi được bảo đảm không âm, hàm không đổi. Bước S6(b) (slot sống) chỉ chạy khi bật \texttt{LIVE\_SWAP=measured}, và chỉ có chứng chỉ tức thời (i.5).
\end{hequa}
\begin{kiem}
T7a $3.4\cdot10^{-14}$ (T7a-NEG đích $-r$ lệch $\ge0.99$); T7b $4.5\cdot10^{-11}$ (T7b-NEG $s=J_v\T r$ lệch trung vị 1.95); T7c, T7d, T7e; T7f (hệ số $\mu/(c_H+\mu)$ giữ ở mọi trial; T7f-WIT: argmax $\ell$ không phải argmax gain ở 32/60); T7g ($|\cos|$ giữa phần bù xiên và Euclid trung vị 0.763); T7i (hàm giữ nguyên, 0.0); T7j; T7k (phần dư bậc ba; NEG loại dấu cũ). \st{empirical}: trên toy hai task 40 seed, chọn theo $\rperp$ thắng không-swap 72\% trên CE lớp cũ, nhưng CE trung bình bằng cách chọn theo $e$.
\end{kiem}

\subsection{T13, T17: nhiều slot, optimizer và BN}

\begin{namedthm}{T13}{nhiều slot, trạng thái optimizer}
Giả thiết A2, A3. Với $\Jcal$ bất kỳ và khối cột mới $V$, đặt $S=V\T H_\perp(\Jcal)V+\mu I$, $s=V\T\rperp(\Jcal)$.
(a) $\mathrm{gain}(\{v_1,\dots,v_k\}\mid J)=\sum_j\mathrm{gain}(v_j\mid[J,J_{v_1},\dots,J_{v_{j-1}}])$.
(b) $\rperp([\Jcal,V])=\rperp(\Jcal)-H_\perp(\Jcal)VS^{-1}s$ và $H_\perp([\Jcal,V])=H_\perp(\Jcal)-H_\perp(\Jcal)VS^{-1}V\T H_\perp(\Jcal)$.
(c) $\mathrm{gain}(\mathcal S\mid J)\ge\max_{v\in\mathcal S}\mathrm{gain}(v\mid J)$.
(d) Với SGD momentum thường (nesterov tắt, dampening $0$), sau khi đặt trạng thái của slot bằng $0$: bước đầu của $B_v$ là $-\mathrm{lr}\nabla_{B_v}L$ và của $a_v$ là $-\mathrm{lr}\,\omega a_v$. Với nesterov hệ số là $(1+\nu)$; với dampening $\delta$ (buffer bằng $0$) là $(1-\delta)$; cả hai là bội dương của gradient.
\end{namedthm}
\begin{proof}
(a) $q^*_r$ chỉ phụ thuộc span; tổng là telescoping. (b) Nghịch đảo khối của ${\Jcal'}\T H\Jcal'+\mu I$ với $\Jcal'=[\Jcal,V]$ cho $H\Jcal'G(\Jcal'){\Jcal'}\T H=H\Jcal G(\Jcal)\Jcal\T H+H_\perp VS^{-1}V\T H_\perp$; với $r\in\range(H)$, $\rperp(\Jcal)=H_\perp(\Jcal)H^+r$. (c) Thêm cột không làm $q^*$ tăng. (d) Quy tắc cập nhật của \texttt{torch.optim.SGD}: $m\leftarrow\nu m+(1-\delta)(\nabla+\omega\theta)$, $\theta\leftarrow\theta-\mathrm{lr}\,m$ (hoặc $\theta-\mathrm{lr}((\nabla+\omega\theta)+\nu m)$ với nesterov), với $m=0$ trước bước.
\end{proof}
\begin{hequa}
Bước S6(a): với nhiều slot, bản chính xác cập nhật phần bù theo (b); bản rẻ giữ $\rperp$ và chọn hướng $S_x$-trực giao, được bảo đảm bởi (c). Sau swap, đặt trạng thái optimizer của mọi phần tử thuộc slot bằng $0$ (d).
\end{hequa}
\begin{kiem}
T13a-WIT (18/40 đơn vị chết trước S5 sống sau S5), T13a, T13a-NEG (chèn không kiểm lại đổi hàm 0.108); T13b (SGD thật, sai lệch $0$), T13b-NEG, T13b2; T13d $5.0\cdot10^{-14}$, T13d-NEG; T13e, T13e-WIT; T13f (Schur đúng với ReLU, $9.0\cdot10^{-13}$); T13g; T13h-WIT.
\end{kiem}

\begin{namedthm}{T17}{slot có BN theo sau}
Giả thiết A1, A2, A12, BN train mode, kênh $u$ có affine $(\gamma_u,\beta_u)$.
(a) Với $B_v=0$, đặt $a_v,\gamma_v,\beta_v$ và running stats tùy ý không đổi hàm trên cửa sổ.
(b) Nếu kênh chết vì $\beta_u\ll0$, giữ nguyên $(\gamma_u,\beta_u)$ thì neuron mới vẫn chết; phải đặt $\gamma_v=\sigma$, $\beta_v$ theo quy tắc dấu/bias.
(c) Bias trước BN bị triệt tiêu; quy tắc bias tác động lên $\beta_{BN}$: $\beta_v=-\operatorname{median}(\sigma\hat x)$ làm khoảng một nửa cửa sổ kích hoạt.
(d) \Tref{T7}(b) đúng với $J_v=\varphi\otimes I$ tính từ đặc trưng BN thực tế.
(e) Đặt running\_mean$_u=0$, running\_var$_u=1$ là quy ước cài đặt, không ảnh hưởng (a)--(d).
\end{namedthm}
\begin{proof}
(a) Đầu ra kênh $u$ chỉ đi qua $B_u$ (A12), mà $B_v=0$; BN theo kênh; ở train mode running stats không vào đầu ra. (b) $\varphi=\mathrm{ReLU}(\gamma\hat x+\beta)\equiv0$ khi $\gamma\hat x_i+\beta\le0$ với mọi $i$. (c) $\hat x$ không phụ thuộc bias trước BN. (d) \Tref{T7}(b) đúng cho mọi đặc trưng. (e) Định nghĩa.
\end{proof}
\begin{hequa}
Bước S6(a) với slot có BN: đặt lại $\gamma,\beta$ và running stats của kênh, áp quy tắc dấu/bias lên $\beta_{BN}$; hàm vẫn được bảo toàn chính xác.
\end{hequa}

\subsection{T12, T16: slot sống}

\begin{namedthm}{T12}{slot sống: đại lượng chính xác}
Nội dung là \Tref{T7}(i) và \Tref{L10}: $j_u^\ell$ được đo chính xác; cận Taylor của nó có hằng số chặt $c_3(1+T^{-3})/N$.
\end{namedthm}
\begin{kiem}
T12a (hai phía: đúng và chặt, cấu hình đạt $1.0000$ cận; NEG loại chuẩn hóa $T^{-1}$ và hằng số ${\times}0.1$; WIT: hằng số $0.708$ cũ đúng nhưng lỏng 2.6 lần); T12a2; T12b ($T=1,2,4$; NEG bỏ số hạng KD vi phạm 2.0 lần); T12c [structural]; T12d (hiệu mô hình trên mạng ReLU thật bằng $-$gain Schur, $1.8\cdot10^{-14}$; NEG dùng $J_v\T e$ lệch $\ge1.53$); T12e-WIT (chết trên cửa sổ có thể kích hoạt ngoài cửa sổ).
\end{kiem}

\begin{namedthm}{T16}{điểm khai triển và độ tin cậy của dự báo slot sống}
Giả thiết A2. Với mọi $f,f_u,\Delta$ và $\ell\in\{L,L_{KD}\}$:
(a) $|\ell(f-f_u+\Delta)-[\ell(f-f_u)+\nabla\ell(f-f_u)\T\Delta+\tfrac12\Delta\T\nabla^2\ell(f-f_u)\Delta]|\le(L_3/6)\norm\Delta^3$.
(b) Nếu dùng $\nabla^2\ell(f)$ thay cho $\nabla^2\ell(f-f_u)$, sai số chỉ bị chặn bởi $(L_3/6)\norm\Delta^3+\tfrac12L_3\norm{f_u}\norm\Delta^2$, có thêm số hạng bậc hai.
(c) Vì vậy trong dự báo, mọi đại lượng sau swap phải lấy tại $f-f_u$; dưới A10 và trong chế độ tuyến tính, sai số dự báo là $(L_3/6)(\norm\Delta^3+\norm{\Delta'}^3)$ cộng sai số tuyến tính hóa của Jacobian dọc bước hồi phục (thuộc A10).
\end{namedthm}
\begin{proof}
(a) Taylor bậc ba tại $f-f_u$ và \Tref{L10}. (b) Hiệu hai mô hình là $\tfrac12\Delta\T(\nabla^2\ell(f-f_u)-\nabla^2\ell(f))\Delta$, chặn bằng $\tfrac12L_3\norm{f_u}\norm\Delta^2$. (c) Áp (a) cho hai trạng thái rồi lấy hiệu.
\end{proof}
\noindent\textbf{Kết quả thực nghiệm quyết định mặc định} (\st{empirical}, MLP ReLU thật, 198 cặp): ở $\mu_{\mathrm{rel}}=10$ dấu của $\Delta L_{\mathrm{pred}}$, $\Delta KD_{\mathrm{pred}}$ khớp bước GN thực hiện ở 194/198 và 193/198 cặp; ở $\mu_{\mathrm{rel}}=10^{-2}$ (damping dùng thật) chỉ 87/198 và 101/198; so với 200 bước GD thì khớp kém kể cả ở $\mu_{\mathrm{rel}}=10$ (131/198, 116/198).
\begin{hequa}
Swap slot sống là \st{measured}, \emph{không} \st{proven}, và \textbf{tắt mặc định}. Hệ quả: với cấu hình mặc định, GIR không hoạt động trong CIL (\Tref{T7}(j)).
\end{hequa}

\section{Quỹ đạo, giải không chính xác và phạm vi}\label{sec:robust}

\begin{namedthm}{T8}{cận Grönwall rời rạc theo defect thực hiện}
Giả thiết A7. Gọi $x_m=f(\theta_m)$ sau $m$ bước vĩ mô độ dài $\tau$, quỹ đạo tham chiếu $y_{m+1}=y_m+\tau U(y_m)$, $y_0=x_0$, và $\delta_m:=\norm{(x_{m+1}-x_m)/\tau-U(x_m)}$ (đo được lúc chạy). Khi đó $\norm{x_m-y_m}\le\sum_{j<m}(1+\tau L_U)^{m-1-j}\tau\delta_j$.
\end{namedthm}
\begin{proof}
$e_m=\norm{x_m-y_m}$ thỏa $e_{m+1}\le(1+\tau L_U)e_m+\tau\delta_m$; quy nạp.
\end{proof}
\begin{remark}
$\delta_m$ gồm (i) defect chiếu $\lambda\norm{Jd-u}$, đúng đại lượng VNA cực tiểu; (ii) độ trôi kernel; (iii) sai số rời rạc. \textbf{Cận chỉ theo phần (i) là sai:} trên MLP tanh không lazy, defect thực hiện trung bình 0.73 so với phần chiếu 0.02 (T8a-NEG, tỉ số vi phạm 2.7).
\end{remark}
\begin{hequa}
Chỉ khẳng định \emph{tối ưu chiếu cho từng bước}. Khẳng định ``quỹ đạo gần mô hình mở rộng'' chỉ được đưa ra khi $\delta_m$ log ở bước S7 là nhỏ.
\end{hequa}

\begin{namedthm}{T9}{giải không chính xác}
Giả thiết A8.
\begin{enumerate}[label=(\alph*)]
\item $\norm{\hat x-x}\le\rho\norm b/\mu$; warm start không ảnh hưởng nếu $\rho$ đo so với $\norm b$.
\item $\hat h=g_O\T\hat d$, $\hat g=g\T\hat d$ và $\leak(\hat d_N)=\norm{\Lambda C\hat d_N}^2$ được tính trực tiếp trên vector áp dụng, không qua solve. Với quy tắc chặt ($\tau=0$), \Tref{T5}(d--f) đúng bất kể sai số solver; với quy tắc có dung sai (e), chúng đúng sai khác tối đa $\lambda\tau_g$ (cho $L$) và $\varepsilon\lambda\tau_h$ (cho KD bậc một).
\item Nếu $\cos(\hat d,d^*)\ge c$ thì $|g_O\T\hat d/\norm{\hat d}-g_O\T d^*/\norm{d^*}|\le\norm{g_O}\sqrt{2(1-c)}$.
\item $\norm{\hat d_{\mathrm{rep}}-\drep}\le(\norm{J\T HK_EHJ}\rho_1\norm{g_O}/\mu+\rho_2\norm{\hat b_2})/\mu$.
\item \emph{Quy tắc bỏ qua có dung sai.} Với chuỗi hai lần giải ($\hat x_1\approx A_\beta^{-1}b_1$, $b_1=J\T e$, residual $\rho_1$; $\hat z=HJ\hat x_1$, $\hat u=K_E\hat z$, $b_2=J\T H\hat u$; $\hat d\approx A_\beta^{-1}b_2$, residual $\rho_2$), đặt
\[
\hat\Delta=\norm{J\T HK_E}\norm{HJ}\rho_1\norm{b_1}/\mu^2+\rho_2\norm{b_2}/\mu,\quad\tau_h=\norm{g_O}\hat\Delta,\quad\tau_g=\norm g\hat\Delta.
\]
Khi đó $\norm{\hat d-d^*}\le\hat\Delta$, $|\hat h-h^*|\le\tau_h$, $|\hat g-g^*|\le\tau_g$. Nếu $h^*\ge0$ và $g^*\ge0$, quy tắc ``bỏ qua khi $\hat h<-\tau_h$ hoặc $\hat g<-\tau_g$'' không bao giờ bỏ nhầm; bước được áp dụng làm KD tăng bậc một không quá $\varepsilon\lambda\tau_h$.
\end{enumerate}
\end{namedthm}
\begin{proof}
(a) $\hat x-x=A^{-1}(A\hat x-b)$, $\norm{A^{-1}}\le\mu^{-1}$. (b) Định nghĩa. (c) Với vector đơn vị, $\norm{u_1-u_2}^2=2-2\cos$; Cauchy--Schwarz. (d) (a) hai lần và bất đẳng thức tam giác. (e) $\norm{\hat x_1-x_1}\le\rho_1\norm{b_1}/\mu$ nên $\norm{\hat z-z}\le\norm{HJ}\rho_1\norm{b_1}/\mu$; với $\tilde d=A_\beta^{-1}b_2$, $\norm{\tilde d-d^*}\le\norm{J\T HK_E}\norm{\hat z-z}/\mu$ và $\norm{\hat d-\tilde d}\le\rho_2\norm{b_2}/\mu$; cộng lại; hai hệ quả còn lại là Cauchy--Schwarz.
\end{proof}
\begin{hequa}
Bước S4: chứng chỉ không bao giờ tính qua solve; dung sai $\tau_h,\tau_g$ lấy từ residual cuối của CG (CIL có hai chuỗi, $\hat\Delta=\hat\Delta_{\mathrm{rep}}+\hat\Delta_N$). Chuẩn phổ ước lượng bằng lặp lũy thừa, nhân hệ số an toàn. Kiểm $\cos(\hat d,d_{\mathrm{exact}})\ge0.99$ ở các điểm kiểm (c).
\end{hequa}
\begin{kiem}
T9a (T9a-NEG: dừng theo residual ban đầu vi phạm $3.2\cdot10^5$ lần); T9b; T18a (0/60 vi phạm); T18a-NEG (bỏ $1/\mu$ vi phạm 26/60); T18b-WIT (quy tắc chặt bỏ nhầm 15/30 khi $V_O=0$); T18c (quy tắc có dung sai bỏ nhầm 0/30).
\end{kiem}

\begin{namedthm}{T18}{quy tắc bỏ qua có dung sai}
Nội dung là \Tref{T9}(e); xem phần kiểm số ở trên.
\end{namedthm}

\begin{namedthm}{T10}{phạm vi: BN, momentum, weight decay, augmentation}
\begin{enumerate}[label=(\alph*)]
\item \emph{BN train mode.} Các đồng nhất và chứng chỉ của \Tref{T1}--\Tref{T5}, \Tref{T7}(a--c, e--i), \Tref{T8}, \Tref{T9}, \Tref{T11}, \Tref{T12} chỉ dùng $J,H,W$ của loss trên cửa sổ, nên đúng chính xác dưới A1. Ngoại lệ: tính tối ưu của bước chọn (\Tref{T6}, \Tref{T7}(d)) cần A4$'$, với dạng riêng khi có BN (\Tref{T14}). Đại lượng tính ở eval mode dự báo sai.
\item Loại tham số BN khỏi $S$: mọi định lý vẫn đúng với $J_S$, và $\range(K_S)\subseteq\range(K)$.
\item A4 (tuyến tính toàn cục) sai khi BN đứng sau kênh ảo: $\varphi(2a)=\varphi(a)$.
\item \emph{Momentum.} Cộng ngoài buffer: \Tref{T5}(f) đúng, không phụ thuộc $\nu$. Ghi $\hat d$ vào buffer: $\hat d$ bị phát lại với tổng trọng số $1/(1-\nu)$ ở các điểm nó không được tính, nên dấu không còn được chứng chỉ.
\item \emph{Weight decay:} hiệu chỉnh thêm số hạng $-\varepsilon\lambda\omega\theta\T\hat d$ vào mục tiêu có regularization; phải log.
\item \emph{Augmentation:} bảo đảm chỉ cho loss thực nghiệm trên đúng các mẫu đã augment trong cửa sổ.
\end{enumerate}
\end{namedthm}
\begin{proof}
(a, b) Các định lý chỉ dùng tính chất đại số của $(J,H,W)$ tại $\theta$. (c) BN chuẩn hóa nên bất biến tỉ lệ. (d) \Tref{T5}(f); số hạng phát lại là $\sum_k\nu^k$. (e, f) Định nghĩa.
\end{proof}
\begin{hequa}
Bước S5: cộng $-\varepsilon\lambda\hat d$ trực tiếp vào tham số, \emph{không} qua momentum buffer; khôi phục buffer BN sau các forward của VNA (A1).
\end{hequa}
\begin{kiem}
T10a $8.4\cdot10^{-14}$; T10b (T10b-NEG: eval mode lệch 4.87 lần); T10d; T10f (\texttt{torch.optim.SGD} thật, momentum 0.9: cộng ngoài buffer dịch đúng $\mathrm{lr}\cdot d$, tỉ số 1.000000; NEG: qua buffer phát lại 10 lần).
\end{kiem}

\begin{namedthm}{T11}{các nhánh suy biến của quy tắc $\beta$}
Giả thiết A2, A3. Quy tắc $\beta$ ở bước S3 cho một $\beta$ mà, với solve chính xác, $h(\beta)\ge0$:
(a) $e_O=0$ (kể cả $c_{\mathrm{old}}=1$): $\beta=0$, $h\equiv0$.
(b) $V_O=0$, $e_O\ne0$: $\beta=0$, $h(0)=0$.
(c) Ngược lại: nếu $\leak(0)\le\varepsilon^*$ thì $\beta=0$; nếu không, bisection trên $\log\beta$ trong $[\beta_{\mathrm{suf}}10^{-k_{\mathrm{lo}}},\beta_{\mathrm{suf}}]$ trả về một đầu mút trên đã kiểm $\leak(\beta)\le\varepsilon^*$. Khoảng này luôn hợp lệ vì $\leak(\beta_{\mathrm{suf}})\le\varepsilon^*$.
\end{namedthm}
\begin{proof}
(a), (b) là \Tref{T3}(g)(iv), (iii). (c) \Tref{T3}(e): mọi $\beta$ có $\leak\le\varepsilon^*$ đều có chứng chỉ; bisection chỉ trả về điểm đã kiểm; tính đơn điệu của $\leak$ (\Tref{T3}(c)) bảo đảm $[\beta_{\mathrm{cert}},\infty)\ni\beta_{\mathrm{suf}}$.
\end{proof}
\begin{hequa}
Hàm \textsc{BetaRule} của \cref{alg:vna} cài đúng ba nhánh (a)--(c); không nhánh nào bỏ qua lần áp dụng, và không có ngưỡng tự do.
\end{hequa}
\begin{kiem}
T11a ($h(0)$ chuẩn hóa $1.6\cdot10^{-14}$ khi $V_O=0$); T11c-WIT (60/60 bài $V_O=0$ có $\leak(0)>0$, nên quy tắc cũ ``bỏ qua khi $V_O=0$'' vứt bỏ các bước an toàn); T11d, T11e; T11g (bisection log 12 bước: 0 sai trên 120 bài).
\end{kiem}

\begin{namedthm}{T15}{kiểm ở mức thuật toán (không phải định lý mới)}
Trên MLP tanh thật với CE + KD của PyCIL, chạy đúng bước S3--S4 rồi đi một bước $\theta-\varepsilon(g+\lambda\hat d)$: 24 lần áp dụng, 1 lần bỏ qua; $L$ giảm ở mọi lần áp dụng, dự báo bậc một đúng tới $1.1\cdot10^{-4}$ (tương đối) ở $\varepsilon=10^{-4}$; riêng bước hiệu chỉnh không làm $L_{KD}$ tăng ở lần nào khi $\hat h\ge0$ (0/24). \st{empirical}.
\end{namedthm}

\section{Một họ cập nhật: tách residual và phép chiếu phần ảo}\label{sec:unified}

Mục này trả lời câu hỏi: hiệu chỉnh $d$, bước Gauss--Newton có sẵn ở bước S1, optimizer gốc và GIR liên hệ với nhau thế nào? Kết luận (đã qua phản biện vòng 5) là:
\begin{itemize}
  \item $d$ là \emph{phép chiếu chính xác} của phần ảo của vận tốc mạng mở rộng, trên cotangent $z$ (\Tref{T19});
  \item bước GN có damping đóng vai \emph{bộ tách} $e=z+\rperp$ (\Tref{T20});
  \item optimizer gốc giữ nguyên động lực của chính nó; việc coi nó là ``phần $K$ của cùng phép chiếu'' chỉ là một \st{interpretation} xấp xỉ, sai khác $O(\mu\norm G)$ phụ thuộc phổ (\Tref{T20}(c), \Tref{T22});
  \item GIR tác động qua $\rperp$; vai trò của các thành phần là \emph{bổ sung} cho nhau dưới phép tách GN có damping, không tách bạch tuyệt đối (\Tref{T23}).
\end{itemize}

Định nghĩa phép chiếu có damping $P(t):=\argmin_x\norm{Jx-t}_H^2+\mu\norm x^2=GJ\T Ht$ (nghiệm duy nhất vì $\mu>0$). Mạng mở rộng có Jacobian $[J,J_E]$ và kernel $K+K_E$ tại trọng số ra $w=0$.

\begin{namedthm}{T19}{phân rã vận tốc mở rộng đã chiếu}
(a) $P((K+K_E)z)=P(Kz)+P(K_Ez)=x_K+d$, với $d=GJ\T HK_Ez$ (VNA ở $\beta=0$) và $x_K:=GJ\T HKz$.
(b) $x_K=GF(J\T z)=GF\cdot FGg=(FG)^2g$.
(c) $x_K$ chính là $d$ khi thay $K_E$ bằng $K$; arm này được gọi là \texttt{kernel\_self}.
(d) Nói chung $x_K\ne\xGN:=Gg$: khi $\mu\to0$ trên $\range(F)$, $x_K\to P_{\range F}\,g$ còn $\xGN\to F^+g$.
\end{namedthm}
\begin{proof}
(a) $P$ tuyến tính. (b) $Kz=J(J\T z)$ và $J\T z=J\T HJGg=FGg$; $F$, $G$ giao hoán. (c) Thay $K_E\to K$ trong $d$. (d) $FG\to P_{\range F}$ và $G\to F^+$ trên $\range(F)$.
\end{proof}
\begin{remark}
Dưới z-flow của mạng mở rộng ($\dot\theta=-J\T z$, $\dot w=-J_E\T z$), vận tốc thực của tham số cũ là $J\T z=FGg=g-\mu\xGN$; $x_K$ là vận tốc đó bị damping thêm một lần bởi phép chiếu có damping. Vì vậy việc gọi $x_K$ là ``phần $K$ trung thực'' phụ thuộc vào lựa chọn phép chiếu; điều được khẳng định chắc chắn là (a)--(c).
\end{remark}
\begin{hequa}
Bước S3 ở $\beta=0$ đúng là $P(K_Ez)$. Arm \texttt{kernel\_self} khác VNA \emph{chỉ ở kernel} đưa vào cùng phép chiếu, nên so sánh VNA với \texttt{kernel\_self} đo đúng đóng góp của kernel ảo.
\end{hequa}

\begin{namedthm}{T20}{tách residual bằng bước GN}
Với $\xGN=Gg$ (lần giải đầu tiên của bước S1):
(a) $z=HJ\xGN=HK(HK+\mu)^{-1}e$ và $\rperp=e-z=\mu(HK+\mu)^{-1}e$.
(b) $Ke-Kz=\mu J\xGN$ và $J\T\rperp=\mu\xGN$.
(c) Nếu $g\in\range(F)$ (đúng khi $e\in\range(H)$, tức đúng cho CE + KD theo \Tref{L5}) thì $\norm{Ke-Kz}/\norm{Ke}\to0$ khi $\mu\to0$. Hằng số không đều: sai khác cỡ $O(\mu\norm G)$ trên phần phổ của $F$ mà $g$ có năng lượng; với phổ lũy thừa có thể lớn ($\approx127\mu_{\mathrm{rel}}$ trong một thí nghiệm của phản biện).
(d) Với $\Delta^*=-J\xGN$: $z=-H\Delta^*$ là phần residual mà \emph{bước GN có damping} loại bỏ; $\rperp=e+H\Delta^*$ là gradient của mô hình bậc hai $e\T\Delta+\tfrac12\Delta\T H\Delta$ còn lại tại $\Delta^*$, tức phần bước có damping để lại.
\end{namedthm}
\begin{proof}
Push-through: $J(J\T HJ+\mu)^{-1}J\T=K(HK+\mu)^{-1}$ và $J\T(HK+\mu)^{-1}=GJ\T$. Do đó $z=HK(HK+\mu)^{-1}e$ và $e-z=\mu(HK+\mu)^{-1}e$; $Ke-Kz=\mu K(HK+\mu)^{-1}e=\mu JGJ\T e$; $J\T\rperp=\mu J\T(HK+\mu)^{-1}e=\mu Gg$. (c) Sai khác $\norm{2\mu Gg-\mu^2G^2g}$ tiến về $0$ khi $g$ nằm trong range của $F$; khi $e\notin\range(H)$ và $H$ suy biến, phản ví dụ cho sai khác 0.58 ở $\mu_{\mathrm{rel}}=10^{-8}$. (d) Thay $\Delta^*$.
\end{proof}
\begin{hequa}
Bước S1 tính $\xGN$ và lưu lại: từ đó có $z$ (cho VNA) và $\rperp$ (cho GIR) mà không cần giải thêm. Với $\mu>0$, $\rperp$ \emph{không} hoàn toàn ngoài tầm với của mạng hiện tại ($J\T\rperp=\mu\xGN\ne0$), và optimizer gốc vẫn tác động lên toàn bộ $e$; do đó vai trò ``VNA trên $z$, GIR trên $\rperp$'' là bổ sung, không tách bạch.
\end{hequa}

\begin{namedthm}{T21}{phân rã mức giảm bậc một; ý nghĩa của $\rho$}
Với $\beta=0$: $g\T x_K=\norm{FGg}^2\ge0$ và $g\T d=z\T K_Ez=V\ge0$. Dưới cập nhật thực sự chạy (bước của optimizer gốc theo $g$ cộng $\lambda d$), $\rho=V/\norm g^2$ là \emph{tỉ số} giữa mức giảm bậc một của mạng cố định do phần ảo và do phần nền, trên một đơn vị trọng số. Ngoài ra $\rho\le\rho_{\mathrm{proj}}:=V/\norm{FGg}^2$, tỉ số tương ứng cho z-flow đã chiếu.
\end{namedthm}
\begin{proof}
$g\T GF^2Gg=\norm{FGg}^2$; $g\T GJ\T HK_EHJGg=z\T K_Ez$ (\Tref{L4}); $\norm{FGg}\le\norm g$ vì $FG\preceq I$.
\end{proof}
\begin{remark}
$\rho$ là tỉ số (ảo trên nền), không phải ``tỉ phần''; tỉ phần sẽ là $V/(\norm{FGg}^2+V)$. Cả hai đại lượng là mức giảm của \emph{mạng cố định} sau khi chiếu, không phải tốc độ của mạng mở rộng: tốc độ giảm loss của mạng mở rộng, $e\T(K+K_E)z$, không có chứng chỉ (phần ảo $e\T K_Ez=V+\rperp\T K_Ez$ âm ở 79/400 bài ngẫu nhiên). Với $\beta>0$, $g\T d^*$ không còn xác định dấu (\Tref{T5}(g)).
\end{remark}
\begin{hequa}
Bước S7 log $\rho$ với đúng nghĩa trên. Cổng chẩn đoán dựa trên $\rho$ được giữ như đã đăng ký trước, nhưng phải đọc là ``tỉ số ảo trên nền'', không phải ``phần tiến bộ của mạng mở rộng đến từ neuron ảo''.
\end{hequa}

\begin{namedthm}{T22}{động lực GN của mạng mở rộng}
Bước $\theta$ của GN có damping cho mạng mở rộng là $\delta'_\theta=-J\T(HK'+\mu)^{-1}e$, $K'=K+K_E$. Tại $K_E=0$ nó bằng $-\xGN$, và đạo hàm bậc một theo $K_E$ là
\[
\partial\delta'_\theta/\partial K_E[K_E]=GJ\T HK_E(HK+\mu)^{-1}e=\tfrac1\mu GJ\T HK_E\rperp .
\]
Theo quy ước cập nhật $\theta\leftarrow\theta-x$, phần ảo của động lực này là $\xGN-\tfrac1\mu GJ\T HK_E\rperp$ ở bậc một, đúng khi $\norm{HK_E}/\mu\ll1$. Số hạng này khác $d$ (lệch tương đối $\ge43$ trên bài kiểm).
\end{namedthm}
\begin{proof}
$(HK'+\mu)^{-1}=(HK+\mu)^{-1}-(HK+\mu)^{-1}HK_E(HK+\mu)^{-1}+O(K_E^2)$, $J\T(HK+\mu)^{-1}=GJ\T$, và \Tref{T20}(a).
\end{proof}
\begin{remark}[phát biểu định lượng thay cho phát biểu định tính đã rút lại]
Khẳng định cũ ``$\xGN+d$ không phải phép chiếu của bất kỳ động lực mở rộng nào'' là \emph{sai}: tồn tại cotangent chung $c$ sao cho $P((K+K_E)c)=\xGN+d$ chính xác. Điều còn đúng là định lượng: $\xGN$ lệch khỏi phần $K$ của z-flow một lượng $O(1)$, còn $g$ lệch $O(\mu\norm G)$. Theo cùng tiêu chí, cập nhật ``optimizer gốc $+\lambda d$'' cũng là một tổ hợp (khối $\theta$ của gradient flow mở rộng cộng phần ảo của z-flow). Chúng tôi chọn tổ hợp này vì nó giữ nguyên optimizer gốc (plug-in) và có chứng chỉ \Tref{T5}; đó là một \emph{lựa chọn thiết kế}, không phải hệ quả bắt buộc.
\end{remark}
\begin{hequa}
Bước GN $\xGN$ là một baseline \emph{ngoài} (đổi optimizer gốc), không phải ablation; nó có sẵn miễn phí từ bước S1.
\end{hequa}

\begin{namedthm}{T23}{lợi của swap được điều khiển bởi $\rperp$}
Mở rộng span bằng cột $J_v$ làm giảm mô hình GN đúng $q^*_e(J)-q^*_e([J,J_v])=\tfrac12s\T S^{-1}s$, $s=J_v\T\rperp$, $S=J_v\T H_\perp J_v+\mu$ (cũng là \Tref{T7}(b)). Sau swap, $K'=K+J_vJ_v\T\succeq K$; ở lần áp dụng kế tiếp, bước S1 tính lại phép tách với $K'$, và phần của $\rperp$ mà neuron vừa lắp với tới được chuyển vào $z'$.
\end{namedthm}
\begin{proof}
\Tref{T7}(b) với $r=e$. Phần sau là \Tref{T7}(g).
\end{proof}
\begin{hequa}
Bước S6 dùng đúng $\rperp$ đã có từ bước S1 (không giải thêm); sau swap, lần áp dụng kế tiếp tự động đưa neuron vừa lắp vào kernel $K'$.
\end{hequa}
\begin{remark}
$S$ cũng phụ thuộc $K$ qua $H_\perp$; ``điều khiển bởi $\rperp$'' nghĩa là tử số của lợi chỉ phụ thuộc $\rperp$.
\end{remark}

\begin{namedthm}{T24}{cỡ bước dọc $d$}
Với $\theta\leftarrow\theta-s\,d$, mức giảm của mô hình GN $s\,g\T d-\tfrac12s^2d\T(F+\mu)d$ đạt cực đại tại $s^*=g\T d/d\T(F+\mu)d$.
\end{namedthm}
\begin{proof}
Hàm bậc hai lõm theo $s$.
\end{proof}
\begin{hequa}
Quy tắc $\lambda$ (thiết lập, phải đăng ký trước): $\lambda=\min(\lambda_{\mathrm{faithful}},s^*/\varepsilon)$ với $\lambda_{\mathrm{faithful}}=K$ dưới A13. Bốn lưu ý: (i) quy tắc bỏ qua phải chạy \emph{trước} khi tính $\lambda$, vì với $\beta>0$ có thể $s^*<0$; (ii) $s^*$ bỏ qua $K$ bước nền trong cùng khoảng; (iii) quy tắc momentum $\lambda=K/(1-m)$ là heuristic, không có dẫn xuất; (iv) các điều kiện bậc hai của \Tref{T5}(e,f) và độ lệch $O(\varepsilon\lambda\norm{\hat d})$ của GIR phải tính lại với $\lambda$ đã chọn (với $\lambda=K$ chúng có thể lớn hơn khoảng $K$ lần so với chuẩn hóa cùng norm). Quy tắc cùng norm $\lambda=\norm g/\norm d$ là một quy ước công bằng, báo cáo như thiết lập phụ. Mọi chứng chỉ bậc một bất biến theo $\lambda>0$.
\end{hequa}

\begin{kiem}
\texttt{check\_unified\_gn.py}: 15/15 đạt (U1--U9 kèm NEG và WIT; U1 $1.5\cdot10^{-13}$, U3 $7.5\cdot10^{-16}$, U2 $2.3\cdot10^{-12}$ với U2-WIT $\cos(x_K,\xGN)$ trung vị 0.50, U4 $2.9\cdot10^{-11}$, U6 sai phân trung tâm khớp $1.4\cdot10^{-6}$, U6-NEG lệch $\ge43$, U8 $6.2\cdot10^{-12}$). Phản biện vòng 5 tự viết 25 phép kiểm độc lập, gồm 8 nhân chứng chống lại các khẳng định cũ đã được sửa ở trên.
\end{kiem}

\subsection{T25: lợi của VNA bị chặn bởi phần residual mà mạng sống với tới được}

Mục này hình thức hóa câu ``VNA cần GIR''. Đặt $B(K):=g\T Gg$ với $g=J\T e$, $G=(F+\mu)^{-1}$; theo chứng minh của \Tref{T7}(a), $B(K)=-2q^*_e(J)$ là mức giảm của mô hình GN mà \emph{mạng sống} đạt được trên residual $e$. Viết $e=Rw$ với $R=H^{1/2}$, $w=R^+e$ (hợp lệ vì $e\in\range(H)$, \Tref{L5}) và $M=RKR$; khi đó $B(K)=w\T M(M+\mu)^{-1}w$, chỉ phụ thuộc $J$ qua $K=JJ\T$.

\begin{namedthm}{T25}{chặn trên theo khả năng với tới}
Giả thiết A2, A3, $\beta=0$, $e\in\range(H)$. Đặt $\lambda_E:=\lambda_{\max}(RK_ER)$.
\begin{enumerate}[label=(\alph*)]
\item $V=z\T K_Ez\le\lambda_E\,B(K)$.
\item $B$ không giảm theo thứ tự Loewner: $K\preceq K'\Rightarrow B(K)\le B(K')$. Đơn vị chết (cột của $J$ bằng $0$ trên cửa sổ) làm $K$ mất số hạng $J_uJ_u\T\succeq0$, nên $B$, và do đó cận ở (a), không tăng.
\item Chèn khối cột $J_v$ (GIR, $r=e$) làm $B$ tăng \emph{đúng} $2\,\mathrm{gain}(v\mid J)=s\T S^{-1}s$, $s=J_v\T\rperp$ (\Tref{T7}(b), \Tref{T23}). Vậy cận của VNA ở lần áp dụng kế tiếp tăng $2\lambda_E\,\mathrm{gain}(v\mid J)$.
\item \emph{Chỉ là điều kiện cần.} Bản thân $V$ không đơn điệu theo $K$; (a)--(c) nói: khi phần với tới được nhỏ thì VNA chắc chắn nhỏ, còn khi nó lớn thì VNA \emph{có thể} lớn.
\end{enumerate}
\end{namedthm}
\begin{proof}
(a) Push-through: $z=HJGg=RM(M+\mu)^{-1}w$ (vì $(RRK+\mu)^{-1}R=R(RKR+\mu)^{-1}$). Đặt $X=M(M+\mu)^{-1}$, $0\preceq X\preceq I$. Khi đó $V=w\T X(RK_ER)Xw\le\lambda_E\,w\T X^2w\le\lambda_E\,w\T Xw=\lambda_E B(K)$, dùng $X^2\preceq X$. Cuối cùng $w\T Xw=e\T JGJ\T e=g\T Gg$ theo cùng đồng nhất push-through.
(b) $x\mapsto x/(x+\mu)$ là hàm đơn điệu toán tử trên $[0,\infty)$ (Löwner); $K\preceq K'\Rightarrow M\preceq M'$.
(c) $B([J,J_v])-B(J)=-2[q^*_e([J,J_v])-q^*_e(J)]=2\,\mathrm{gain}(v\mid J)$.
(d) Phản ví dụ số ở mục kiểm.
\end{proof}
\begin{hequa}
Chuỗi lập luận ``GIR mở không gian, VNA khai thác'' nay là hệ quả: đơn vị chết làm co cận trên của mức giảm bậc một mà VNA có thể cho (b); GIR là thao tác \emph{duy nhất} trong thuật toán làm cận đó tăng, và tăng đúng bằng lợi của chính nó (c). Hệ quả kiểm được: khi tỉ lệ đơn vị chết cao và không có GIR, VNA bị chặn bởi một đại lượng nhỏ (B5: 77--95\% đơn vị chết ở 50 task cuối); lợi của VNA khi có GIR phải lớn dần theo độ rộng mạng. Log $B(K)$ và $\lambda_E$ ở bước S7 để kiểm cận (a) trên dữ liệu thật.
\end{hequa}
\begin{kiem}
\texttt{check\_t25.py}, 300 bài ngẫu nhiên (gồm 30\% $H$ suy biến): (a) 0 vi phạm, độ chặt $V/(\lambda_EB)$ trung vị 0.39, lớn nhất 0.985 (cận không rỗng); (b) bỏ ngẫu nhiên khoảng 40\% cột, 0 vi phạm; (c) $B([J,J_v])-B(J)=2\,\mathrm{gain}$ tới $10^{-8}$; T25d-WIT: $V$ \emph{tăng} sau khi một đơn vị chết ở 119/300 bài, nên chỉ cận là đơn điệu.
\end{kiem}
\begin{remark}[bằng chứng hiện có, \st{empirical}]
B5, 3 seed: VNA không có GIR kém BP (1.315 so với 1.209), nhất quán với (b). VNA có GIR so với GIR: 0.548 so với 0.561, thắng 84/150 task ghép cặp, chưa phân biệt được với nhiễu. Như vậy phần ``cần'' của (d) được dữ liệu ủng hộ, phần ``bổ sung'' chưa được chứng minh bằng thực nghiệm. Trong CIL online có BN (O1) không có đơn vị chết và không có KD, nên cả T25 lẫn \Tref{T3} không có tác dụng; kết quả đó không kiểm định phương pháp.
\end{remark}

\paragraph{Phát biểu trung thực về ``một cơ chế''.} Đây là \emph{một họ cập nhật} trong đó kernel đưa vào phép chiếu là thành phần duy nhất thay đổi giữa các arm; khi $\mu$ nhỏ so với phổ của $F$ trên $g$, họ cập nhật này xấp xỉ phép chiếu của một z-flow mở rộng duy nhất.

\section{Đóng các lựa chọn tự do: trust region, GIR có bảo vệ, swap theo lợi ròng}\label{sec:closure}

Bộ khung (\cref{sec:core,sec:select,sec:unified}) không có hằng số tự do, nhưng phần nối các bước còn bốn lựa chọn tay: cỡ bước (cap $\mathrm{lr}\norm g/\norm{\hat d}$, chặn ở 99.5\% lần áp dụng trên B5), damping $\mu$ (chọn từ lưới, điều kiện preflight không đạt), residual của GIR ($\eO$ hay $e$), và slot nào được thay (chỉ slot chết). Mục này thay cả bốn bằng hệ quả của cùng một bài toán: \emph{mô hình GN có bảo vệ, trong vùng mà mô hình đó còn đáng tin}.

Quy ước: $\theta\leftarrow\theta-x$; mô hình GN của loss cửa sổ là $m(x)=L-g\T x+\tfrac12x\T Fx$; $\xGN(\mu)=(F+\mu)^{-1}g$.

\subsection{T26: damping là nhân tử của trust region}

\begin{namedthm}{T26}{trust region thay cho $\mu_{\mathrm{rel}}$ và cap bước}
Giả thiết A2, $g\ne0$.
\begin{enumerate}[label=(\alph*)]
\item \emph{(Moré--Sorensen \cite{more1983}.)} Với mọi $\mu>0$, $\xGN(\mu)$ là nghiệm duy nhất của $\min m(x)$ với $\norm x\le\Delta(\mu):=\norm{\xGN(\mu)}$, và $\mu$ là nhân tử Lagrange. $\Delta(\mu)$ giảm chặt theo $\mu$, nên $\mu$ và bán kính tin cậy $\Delta$ xác định lẫn nhau.
\item Tương tự, $P_\mu(t)=GJ\T Ht$ là nghiệm của $\min\norm{Jx-t}_H^2$ với $\norm x\le\norm{P_\mu(t)}$. Vậy hiệu chỉnh $\hat d$ là \emph{phép chiếu trong vùng tin cậy} của vận tốc ảo; khi có bảo vệ, $\hat d$ là nghiệm KKT với hai nhân tử: $\beta$ (ràng buộc output cũ) và $\mu$ (ràng buộc vùng tin cậy).
\item \emph{Ý nghĩa của phép tách.} $z(\mu)=HJ\xGN(\mu)$ là phần residual mà bước GN \emph{trong vùng tin cậy} loại được; $\norm{\rperp(\mu)}_{H^+}$ tăng và $B(\mu)=g\T Gg$ giảm theo $\mu$. Mô hình càng kém tin cậy thì càng nhiều residual được giao cho GIR.
\item \emph{Đo độ tin cậy.} Với một forward thêm, $\varrho=[L(\theta)-L(\theta-\xGN)]/[m(0)-m(\xGN)]$. Quy tắc cập nhật chuẩn của Levenberg--Marquardt \cite{madsen2004}: nếu $\varrho>0$ thì $\mu\leftarrow\mu\max(\tfrac13,1-(2\varrho-1)^3)$, ngược lại $\mu\leftarrow2\mu$. Các hằng số là hằng số giáo khoa; lý thuyết hội tụ của trust region đúng cho cả một khoảng giá trị \cite{nocedal2006}. Vì các lần áp dụng xen giữa các bước SGD, ở đây \emph{không} khẳng định hội tụ; $\varrho$ chỉ được dùng như độ trung thành đo được của mô hình GN tại điểm đang xét.
\item \emph{Cỡ bước.} $s=\min(s_F^*,\Delta/\norm{\hat d})$, $s_F^*=\hat g/\hat d\T F\hat d$ ($s_F^*=\infty$ nếu mẫu số bằng $0$). Khi đó $\norm{s\hat d}\le\Delta$ và $m(s\hat d)\le m(0)-\tfrac12s\hat g$.
\item \emph{Chấp nhận chính xác trên cửa sổ.} Sau khi áp dụng, một forward cho $L(\theta')$, $L_{KD}(\theta')$; chấp nhận khi cả hai không tăng, ngược lại hoàn tác và coi như $\varrho<0$. Khi đó mỗi lần áp dụng \emph{không làm tăng} $L$ và $L_{KD}$ trên cửa sổ một cách chính xác, không chỉ ở bậc một.
\end{enumerate}
\end{namedthm}
\begin{proof}
(a) $m$ lồi ($F\succeq0$); KKT của bài toán có ràng buộc: $(F+\mu)x=g$, $\mu\ge0$, $\mu(\norm x-\Delta)=0$; với $\mu>0$ ràng buộc chặt và $x=\xGN(\mu)$ thỏa. Lồi nên KKT đủ; nghiệm duy nhất vì $F+\mu\succ0$. $\Delta(\mu)^2=\sum_ig_i^2/(\lambda_i+\mu)^2$ trong cơ sở riêng của $F$, giảm chặt khi $g\ne0$.
(b) Như (a) với mục tiêu $\norm{Jx-t}_H^2$, gradient $2(Fx-J\T Ht)$. Với bảo vệ, thêm số hạng $\beta\norm{C(x-\drep)}_W^2$ (\Tref{T1}); KKT có hai nhân tử.
(c) Trong tọa độ $w$ của \Tref{T25}: $\rperp\leftrightarrow\mu(M+\mu)^{-1}w$, mỗi hệ số $\mu/(\lambda_i+\mu)$ tăng theo $\mu$; $B=\sum\lambda_iw_i^2/(\lambda_i+\mu)$ giảm.
(d) Định nghĩa và \cite{madsen2004}.
(e) $\norm{s\hat d}\le\Delta$ theo định nghĩa. $m(s\hat d)-m(0)=-s\hat g+\tfrac12s^2\hat d\T F\hat d\le-s\hat g+\tfrac12s\hat g$ vì $s\le s_F^*$.
(f) Theo định nghĩa của phép kiểm.
\end{proof}
\begin{hequa}
$\mu$ không còn được chọn: nó là đối ngẫu của bán kính tin cậy, và bán kính được điều chỉnh theo $\varrho$ đo được. Cap $\mathrm{lr}\norm g/\norm{\hat d}$ (đơn vị là ``một bước SGD'') được thay bằng $\Delta$ (đơn vị là ``bước GN mà mạng tin được''). Quy tắc $\lambda=K$, giả thiết A13 và các arm \texttt{\_lf} không còn cần. Phép kiểm (f) cho chứng chỉ chính xác trên cửa sổ; chứng chỉ bậc một của S4 và dung sai $\tau$ chỉ còn là bộ lọc rẻ trước khi thử.
\end{hequa}
\begin{kiem}
\texttt{check\_t26\_28.py}, 200 bài (H suy biến kiểu softmax): T26a so với SLSQP có ràng buộc, 0 vi phạm; T26b ($\Delta$ giảm chặt) 0 vi phạm; T26c ($B$ giảm theo $\mu$) 0 vi phạm.
\end{kiem}

\subsection{T27: GIR trên bài toán có bảo vệ}

Gọi $x_R=G g_O$ là bước GN trên riêng residual KD và $\delta_O=-Cx_R$ là độ dịch logit cũ mà bước đó tạo ra (đích của phần sửa, tương tự $C\drep$ trong \Tref{T1}). Mô hình có bảo vệ trên độ dịch output $y$:
\[
\hat Q(y)=e\T y+\tfrac12y\T Hy+\tfrac\beta2\norm{Oy-\delta_O}_W^2,\qquad \hat q^*(\Jcal)=\min_x\hat Q(\Jcal x)+\tfrac\mu2\norm x^2 .
\]

\begin{namedthm}{T27}{GIR có bảo vệ}
Giả thiết A2, A3, $\beta\ge0$ do \textsc{BetaRule} trả về. Đặt $H_\beta=H+\beta O\T WO$, $e_\beta=e-\beta O\T W\delta_O$.
\begin{enumerate}[label=(\alph*)]
\item $\hat q^*(\Jcal)=-\tfrac12(\Jcal\T e_\beta)\T(\Jcal\T H_\beta\Jcal+\mu)^{-1}(\Jcal\T e_\beta)+\mathrm{const}$: một mô hình GN thường với metric $H_\beta$ và residual $e_\beta$. Với $\Jcal=J$, ma trận hệ là $J\T H_\beta J+\mu=F+\beta B+\mu=A_\beta$, \emph{đúng ma trận} mà bước S3 giải cho VNA.
\item Với mọi khối cột mới $J_v$: $\hat q^*(J)-\hat q^*([J,J_v])=\tfrac12s\T S^{-1}s\ge0$, với $s=J_v\T\rperp^\beta$, $\rperp^\beta=\mu(H_\beta K+\mu)^{-1}e_\beta$, $S=J_v\T H_\perp^\beta J_v+\mu$, $H_\perp^\beta=H_\beta-H_\beta JG_\beta J\T H_\beta$. Cận dưới không cần giải hệ và hệ số xấp xỉ giữ như \Tref{T7}(d) với $H\to H_\beta$.
\item $\beta=0$: $H_\beta=H$, $e_\beta=e$, trở về đúng quy tắc của B5 ($r=e$). $\beta$ lớn: lợi đo phần học mới mà neuron đem lại \emph{trong khi} logit cũ được giữ ở đích $\delta_O$.
\item Không có tham số mới: $\beta$ là $\beta$ của \textsc{BetaRule}. Quy tắc cũ ``$r=\eO$ trong CIL'' không phải giới hạn nào của họ này và được rút lại.
\end{enumerate}
\end{namedthm}
\begin{proof}
(a) Khai triển $\hat Q(\Jcal x)=x\T\Jcal\T(e-\beta O\T W\delta_O)+\tfrac12x\T\Jcal\T(H+\beta O\T WO)\Jcal x+\tfrac\beta2\norm{\delta_O}_W^2$; cực tiểu của hàm bậc hai lồi chặt. $J\T O\T WOJ=C\T WC=B$.
(b) $e\in\range(H)$ (\Tref{L5}) và $O\T W\delta_O\in\range(O\T WO)$; vì $H,O\T WO\succeq0$, $\range(H_\beta)=\range(H)+\range(O\T WO)$, nên $e_\beta\in\range(H_\beta)$. Áp \Tref{T7}(b) với $(H,r)\to(H_\beta,e_\beta)$. Không âm vì $\hat q^*$ là cực tiểu của \emph{cùng một} hàm trên span lớn hơn ($\delta_O$ cố định theo $J$ hiện tại, không phụ thuộc $J_v$).
(c) Thay $\beta=0$; giới hạn $\beta\to\infty$ như \Tref{T2}(b) với ràng buộc $\Lambda(Oy-\delta_O)=0$.
(d) Định nghĩa.
\end{proof}
\begin{hequa}
Bước S6 dùng $\rperp^\beta$ thay cho $\rperp(\eO)$ hay $\rperp(e)$. VNA và GIR nay giải \emph{cùng một} mô hình (metric $H_\beta$, ma trận $A_\beta$): VNA tìm hệ số trên span hiện tại cho vận tốc ảo, GIR mở rộng span; $\beta$ là một con số duy nhất cân giữa học mới và giữ cũ cho cả hai. \Tref{T25} đúng nguyên văn với $(H,e)\to(H_\beta,e_\beta)$.
\end{hequa}
\begin{kiem}
T27a (dạng đóng so với cực tiểu trực tiếp) và T27b ($A_\beta=J\T H_\beta J+\mu$): 0 vi phạm; T27c (Schur so với hiệu $\hat q^*$ tính trực tiếp, và $\ge0$): 0/200 vi phạm.
\end{kiem}

\subsection{T28: thay slot theo lợi ròng}

Slot sống $u$ có đóng góp $f_u$ và các cột $J_u$; thay bằng $v$ ($B_v=0$) thì output thành $f-f_u$ và span thành $J'=[J_{\setminus u},J_v]$.

\begin{namedthm}{T28}{lợi ròng của swap; slot chết là trường hợp riêng}
Giả thiết như \Tref{T27}.
\begin{enumerate}[label=(\alph*)]
\item Mô hình sau swap và hồi phục trong span mới:
\[\hat q_{\mathrm{sau}}=\hat Q(-f_u)-\tfrac12(J'^{\top}e')\T(J'^{\top}H_\beta J'+\mu)^{-1}(J'^{\top}e'),\qquad e'=e_\beta-H_\beta f_u .\]
Lợi ròng $\Gamma(u,v)=\hat q^*(J)-\hat q_{\mathrm{sau}}$ gồm ba phần: giá của bước nhảy $\hat Q(-f_u)-\hat Q(0)$ (đã gồm phạt trên logit cũ), mất span của $u$, và lợi span của $v$.
\item Slot chết ($f_u=0$, $J_u=0$): $\Gamma(u,v)$ bằng đúng lợi của \Tref{T27}(b). Quy tắc slot chết là trường hợp riêng.
\item \emph{Quy tắc.} Swap khi $\Gamma(u,v)>0$ \emph{và} bước nhảy đo chính xác thỏa $j_u^L\le0$, $j_u^{KD}\le0$ (\Tref{T7}(i.5)). Điều kiện thứ hai là chứng chỉ chính xác tức thời; điều kiện thứ nhất là dự báo của mô hình GN về lợi sau hồi phục. Với slot chết cả hai tự động đúng khi $s\ne0$.
\item Thứ tự xét ứng viên $u$ theo $\norm{f_u}_{H_\beta}$ (đóng góp đo trong cùng metric) chỉ là thứ tự tìm kiếm, không phải tiêu chí quyết định.
\end{enumerate}
\end{namedthm}
\begin{proof}
(a) $\hat Q(-f_u+J'x)=\hat Q(-f_u)+x\T J'^{\top}\nabla\hat Q(-f_u)+\tfrac12x\T J'^{\top}H_\beta J'x$ với $\nabla\hat Q(-f_u)=e-Hf_u-\beta O\T W(Of_u+\delta_O)=e_\beta-H_\beta f_u$; cực tiểu theo $x$ cùng số hạng $\tfrac\mu2\norm x^2$.
(b) Thay $f_u=0$; các cột bằng $0$ không đổi $\hat q^*$; còn lại là \Tref{T27}(b).
(c) \Tref{T7}(i.1, i.5).
(d) Định nghĩa.
\end{proof}
\begin{hequa}
Bước S6 không còn phân biệt ``slot chết'' và ``slot sống'' như hai cơ chế: một quy tắc duy nhất, slot chết là trường hợp $f_u=0$. Đây là điều kiện để GIR hoạt động trên mạng có BN. Phần chứng nhận được vẫn chỉ là tức thời; $\Gamma>0$ là dự báo, phụ thuộc độ trung thành của mô hình hồi phục. Trust region của \Tref{T26} làm dự báo này đáng tin hơn: ở damping lớn, dấu dự báo khớp bước GN thực hiện ở 194/198 cặp, so với 87/198 ở $\mu_{\mathrm{rel}}=10^{-2}$ (\st{empirical}, \Tref{T16}).
\end{hequa}
\begin{kiem}
T28a (dạng đóng so với BFGS trực tiếp) 0/200 vi phạm; T28b (slot chết trùng \Tref{T27}) 0/200 vi phạm.
\end{kiem}

\subsection{Các lựa chọn nhỏ còn lại}

\begin{itemize}
\item \textbf{Số swap.} Mọi slot thỏa \Tref{T28}(c), không giới hạn $S$; cập nhật phần bù tuần tự theo \Tref{T13}(b); tổng lợi là telescoping (\Tref{T13}(a)).
\item \textbf{Bias của neuron mới.} Chọn $b$ cực đại $\ell_\beta(v)$ trên các quantile của $\sigma Xa_k$; cùng tiêu chí gain, không thêm tiêu chí.
\item \textbf{Làm mới $A$.} Mỗi lần áp dụng (tiny\_V rẻ hơn một lần giải CG); bỏ hằng số ``10 lần''.
\item \textbf{Dung sai $\tau$.} Không còn quyết định gì: phép kiểm chính xác \Tref{T26}(f) thay vai trò chứng chỉ.
\end{itemize}

\paragraph{Sau mục này, thuật toán còn đúng ba tham số:} $r$ (độ rộng ảo, tham số khái niệm của chính câu chuyện), ngân sách compute (quyết định $K$), và base method cùng cửa sổ của nó. Mọi đại lượng khác ($\mu$, cỡ bước, $\beta$, residual của GIR, slot được thay, bias) là nghiệm của mô hình GN có bảo vệ trong vùng tin cậy.

\begin{tabularx}{\linewidth}{@{}lXl@{}}
\toprule
Lựa chọn & Trước & Sau \\
\midrule
damping $\mu$ & lưới $\mu_{\mathrm{rel}}$ + preflight & nhân tử trust region, cập nhật theo $\varrho$ (T26) \\
cỡ bước & $\min(s^*,\mathrm{lr}\norm g/\norm{\hat d})$ & $\min(s_F^*,\Delta/\norm{\hat d})$ + kiểm chính xác (T26) \\
chứng chỉ & bậc một, dung sai $\tau$ & chính xác trên cửa sổ (T26f) \\
residual GIR & $\eO$ (CIL) / $e$ (B5) & $e_\beta$, metric $H_\beta$ (T27) \\
slot được thay & chỉ slot chết, $S=1$ & mọi slot có $\Gamma>0$ và $j\le0$ (T28) \\
bias neuron mới & $\{0,-\mathrm{median}\}$ & cực đại $\ell_\beta$ \\
làm mới $A$ & mỗi 10 lần & mỗi lần \\
\bottomrule
\end{tabularx}

\section{Giải chính xác và vòng tin cậy trong mỗi lần áp dụng}\label{sec:exact}

\paragraph{Cập nhật cài đặt sau profiling T4 (21/10/2026).}
T29 vẫn là đồng nhất đại số và bản Gram là tham chiếu cho bài nhỏ.
Cài đặt sản xuất dùng CG trên hệ đối xứng $A=J\T MJ+\mu I$,
không dựng $G_{\mathrm{out}}$, $L$ toàn cục hay phân rã output.
T30, T31 dùng cùng tuyến tính hóa và $\mu^*$; nghiệm số được kiểm residual
như bổ đề dưới đây. Đây là đổi bộ giải có sai số kiểm được, không phải
khẳng định bitwise hay mọi quyết định rời rạc đều trùng bản Gram.


\paragraph{Vì sao cần mục này.} Lần chạy M1 đầu tiên (noBN, seed 0, commit \texttt{d88abc4}) có ba lần áp dụng, và VNA không được dùng lần nào:
\begin{itemize}
\item Lần 1 bị bỏ qua vì CG cho $z$ không đạt $0.1$ sau 200 vòng.
\item Lần 2 bị bỏ qua vì $\varrho=-23.9$.
\item Lần 3 có $\varrho=\mathrm{NaN}$, lọt qua cổng kiểm vì phép so sánh $\varrho\le0$ trả về sai, rồi bị phép kiểm cửa sổ bác ($L$ từ $5.93$ lên $7.49$).
\end{itemize}
Nguyên nhân chung là $\mu\approx7\cdot10^{-8}$, trong khi $\lambda_E\approx30$. Bước GN gần như không có damping ($\norm{\xGN}=289$ đến $656$). Quy tắc của \Tref{T26}(d) đẩy $\mu$ lên quá chậm: mỗi lần chỉ nhân $2$ và chỉ có hiệu lực ở lần áp dụng sau, trong khi mỗi task chỉ có $2$ đến $3$ lần áp dụng.

GIR vẫn chạy: 55 swap slot sống. Các $\Gamma$ được tính với $\mu$ đó và với CG còn residual $0.01$ đến $0.05$. GIR chiếm $85$ đến $98\%$ thời gian của mỗi lần áp dụng.

Mục này sửa ba điểm bằng định lý chứ không bằng tinh chỉnh:
\begin{itemize}
\item mọi hệ được giải chính xác trong không gian output (T29);
\item $\mu$ được chứng nhận \emph{trong} mỗi lần áp dụng, và VNA cùng GIR dùng chung $\mu$ đó (T30);
\item mọi quyết định của GIR trong một lần áp dụng dùng \emph{một} mô hình (T31).
\end{itemize}

Ký hiệu: $J\in\R^{m\times P}$ là Jacobian của logit trên cửa sổ, với $m=N\cdot C$. $G_{\mathrm{out}}=JJ\T\in\R^{m\times m}$. $M\succeq0$ là metric trên output, khối $C\times C$ theo từng mẫu, với $M=H$ hoặc $M=H_\beta=H+\beta O\T WO$. $A=J\T MJ+\mu I$.

\subsection{T29: mọi hệ của thuật toán được giải trong không gian output}

\begin{namedthm}{T29}{giải chính xác bằng ma trận Gram của output}
Cho $\mu>0$ và $M\succeq0$, viết $M=LL\T$.
\begin{enumerate}[label=(\alph*)]
\item \emph{Vế phải nằm trong ảnh của $J\T$.} Với $b=J\T c$, nghiệm duy nhất của $Ax=b$ là $x=J\T y$ với $y=(\mu I+MG_{\mathrm{out}})^{-1}c$. Ma trận $\mu I+MG_{\mathrm{out}}$ luôn khả nghịch.
\item \emph{Dạng đối xứng.} $y=\tfrac1\mu\bigl[c-L(\mu I+L\T G_{\mathrm{out}}L)^{-1}L\T G_{\mathrm{out}}c\bigr]$. Hệ cần giải là đối xứng xác định dương, cỡ $m\times m$.
\item \emph{Mọi đại lượng của S1 đến S6 tính được từ $y$ và $G_{\mathrm{out}}$:}
\begin{itemize}
\item $Jx=G_{\mathrm{out}}y$ và $z=MG_{\mathrm{out}}y$;
\item $\norm x^2=y\T G_{\mathrm{out}}y$, nên $\Delta=\norm{\xGN}$ tính được;
\item $x\T J\T MJx=y\T G_{\mathrm{out}}MG_{\mathrm{out}}y$;
\item $B(K)=c\T G_{\mathrm{out}}y$;
\item $\rperp^\beta=\mu(MG_{\mathrm{out}}+\mu)^{-1}e_\beta$, chính là công thức trong \Tref{T27}(b).
\end{itemize}
Không gian tham số chỉ xuất hiện khi dựng $G_{\mathrm{out}}$ và ở một VJP cuối, khi cần áp $x=J\T y$ vào mạng.
\item \emph{Thêm hoặc bớt một khối cột.} Đặt $\hat J=L\T J$ và $K=\hat J\hat J\T$.
\begin{itemize}
\item Thêm các cột $J_v$, với $\hat J_v=L\T J_v$. Phần bù Schur là $S=\mu\bigl(I+\hat J_v\T(K+\mu I)^{-1}\hat J_v\bigr)$, trong đó $K$ là Gram \emph{không} chứa $J_v$.
\item Bớt các cột $J_u$: $[A^{-1}]_{uu}=\tfrac1\mu\bigl(I-\hat J_u\T(K+\mu I)^{-1}\hat J_u\bigr)$, trong đó $K$ là Gram \emph{có} chứa $J_u$.
\item Đổi $J_u$ lấy $J_v$: $G_{\mathrm{out}}'=G_{\mathrm{out}}-J_uJ_u\T+J_vJ_v\T$, một cập nhật hạng thấp.
\end{itemize}
\item \emph{Swap slot chết là cập nhật chính xác.} Giả sử $u$ chết trên cửa sổ (tiền kích hoạt $\le0$ ở mọi mẫu, với quy ước $\sigma'(0)=0$) và neuron mới được cài với $B_v=0$. Khi đó hàm của mạng và mọi cột Jacobian khác không đổi, $J_u=0$, cột theo trọng số vào của $v$ bằng $0$, và $G_{\mathrm{out}}'=G_{\mathrm{out}}+J_{B_v}J_{B_v}\T$ \emph{chính xác}, không phải xấp xỉ tuyến tính.
\item \emph{Ổn định số.} Giả sử $G_{\mathrm{out}}$ được dựng có sai số $\delta G$ với $\norm{M}\norm{\delta G}<\mu$. Khi đó $y$ tính được là nghiệm chính xác của mô hình GN với Gram $G_{\mathrm{out}}+\delta G$. Phép đo $\varrho$ và phép kiểm cửa sổ của \Tref{T26}(f) dùng loss thật, nên chúng không phụ thuộc vào $\delta G$.
\end{enumerate}
\end{namedthm}
\begin{proof}
(a) Thử $x=J\T y$: $Ax=J\T MG_{\mathrm{out}}y+\mu J\T y=J\T(MG_{\mathrm{out}}+\mu)y=J\T c$. Vì $A\succ0$ nên nghiệm là duy nhất, và đó chính là $x$. Các trị riêng khác $0$ của $MG_{\mathrm{out}}=LL\T G_{\mathrm{out}}$ trùng với các trị riêng của $L\T G_{\mathrm{out}}L\succeq0$, nên $MG_{\mathrm{out}}+\mu$ có mọi trị riêng $\ge\mu>0$.
(b) Đẳng thức Woodbury: $(\mu I+LL\T G_{\mathrm{out}})^{-1}=\tfrac1\mu\bigl[I-L(\mu I+L\T G_{\mathrm{out}}L)^{-1}L\T G_{\mathrm{out}}\bigr]$. Nhân hai vế với $(\mu I+LL\T G_{\mathrm{out}})$ để kiểm.
(c) Thay $x=J\T y$ vào từng biểu thức. Với $\rperp^\beta$, áp (a) cho $c=e_\beta$ và dùng $e_\beta-MJx=(I-MG_{\mathrm{out}}(\mu+MG_{\mathrm{out}})^{-1})e_\beta=\mu(\mu+MG_{\mathrm{out}})^{-1}e_\beta$.
(d) Với $A_0=\hat J\T\hat J+\mu$, ta có $S=\hat J_v\T\hat J_v+\mu-\hat J_v\T\hat JA_0^{-1}\hat J\T\hat J_v$. Lại có $\hat JA_0^{-1}\hat J\T=K(K+\mu)^{-1}=I-\mu(K+\mu)^{-1}$, nên $S=\mu I+\mu\hat J_v\T(K+\mu)^{-1}\hat J_v$. Khối $[A^{-1}]_{uu}$ là khối con của Woodbury $A^{-1}=\tfrac1\mu[I-\hat J\T(K+\mu)^{-1}\hat J]$. Dòng cuối của (d) là định nghĩa của $G_{\mathrm{out}}$.
(e) Neuron $u$ chết nên đầu ra của nó bằng $0$ trên cửa sổ. Do đó $\partial f/\partial B_u=0$, và $\partial f/\partial a_u=0$ vì $\sigma'=0$. Neuron $v$ có $B_v=0$ nên không đóng góp vào forward: hàm mạng, mọi kích hoạt và mọi đạo hàm theo tham số khác giữ nguyên. Đạo hàm theo $a_v$ đi qua $B_v=0$ nên bằng $0$. Cột duy nhất khác $0$ mới xuất hiện là $J_{B_v}$.
(f) Áp (a) cho $G_{\mathrm{out}}+\delta G$. Điều kiện $\norm M\norm{\delta G}<\mu$ giữ cho ma trận hệ khả nghịch. Tính ổn định ngược của Cholesky hoặc LU cho hệ $m\times m$ là kết quả chuẩn \cite{higham2002}.
\end{proof}
\begin{hequa}
\begin{itemize}
\item Bản Gram không cần CG. Bản sản xuất dùng CG--IR với residual double $10^{-10}$, tối đa 200 vòng CG cho mỗi hệ và warm start chỉ trong một lần áp dụng (chi tiết phía dưới). Không đạt dung sai là thất bại của lần thử LM (nếu đang ở LM), không phải nghiệm được chấp nhận. Các bước sau LM thất bại thì hủy và khôi phục toàn bộ lần áp dụng.
\item Ở bản Gram tham chiếu, mọi hệ của S1 đến S6 (gồm $g$, $g_O$, $d_{\mathrm{rep}}$, $d_N$ với mọi $\beta$, $\rperp^\beta$ và mọi $\Gamma$) dùng chung một $G_{\mathrm{out}}$. Mỗi $\beta$ cần một phân rã cỡ $m\times m$. Mỗi ứng viên của GIR chỉ cần các lát Jacobian $J_u$, $J_v$.
\item Phép đồng nhất này chính là đồng nhất push-through mà \Tref{T25} và \Tref{T27} đã dùng: công thức trong định lý trở thành chính code.
\item Chi phí dựng $G_{\mathrm{out}}$ là khoảng $m$ cặp VJP và JVP, nhưng bộ nhớ là $O(m^2)$; các thừa số và workspace còn tăng thêm bộ nhớ. Bản sản xuất bỏ chi phí này bằng matvec CG--IR; các phép đo dung lượng phải bao gồm snapshot double và $D_{\mathrm{site}}$, không chỉ bộ giải.
\end{itemize}
\end{hequa}
\begin{kiem}
\texttt{check\_t29\_31.py}, 200 bài, $H$ và $H_{\mathrm{old}}$ suy biến kiểu softmax, $\beta\in[0,5]$, $\mu\in[10^{-4},1]$:
\begin{itemize}
\item T29a ($x=J\T y$): 0 vi phạm;
\item T29b (dạng đối xứng): 0 vi phạm;
\item T29c ($\norm x^2=y\T G_{\mathrm{out}}y$): 0 vi phạm;
\item T29d (Schur trong không gian output so với phần bù trực tiếp): 0 vi phạm.
\end{itemize}
\end{kiem}


\subsection{Bổ đề sai số cho bộ giải CG}

Đặt $A=J\T MJ+\mu I\succeq\mu I$, $b=J\T c$, $x^*=A^{-1}b$,
và residual thật $r=b-A\widetilde x$. Nếu $\norm r\le\varepsilon\norm b$ thì
\[
 \norm{\widetilde x-x^*}\le\frac{\norm r}{\mu}
 \le\frac{\varepsilon\norm b}{\mu}.
\]
Với đại lượng mà code dùng $\widetilde q=-\tfrac12 b\T\widetilde x$,
\[
 |\widetilde q-q^*|\le\frac{\norm b\,\norm r}{2\mu}
 \le\frac{\varepsilon\norm b^2}{2\mu}.
\]
Do đó, nếu hai hệ trong $\Gamma$ có cùng $\mu^*$ và cùng dung sai,
\[
 |\widetilde\Gamma-\Gamma|
 \le\frac{\varepsilon}{2\mu^*}(\norm{b_{\mathrm{trước}}}^2+
                                    \norm{b_{\mathrm{sau}}}^2).
\]
Đây là $O(\varepsilon)$ khi dữ liệu và $\mu^*>0$ cố định; không phải cận
đều theo $\mu$, và không bảo đảm dấu hay tập swap giữ nguyên khi lợi gần $0$.
Nếu đánh giá hàm bậc hai đầy đủ tại $\widetilde x$ thì sai số giá trị là
$\tfrac12r\T A^{-1}r\le\norm r^2/(2\mu)$; code dùng dạng bilinear ở trên.
\begin{proof}
$\widetilde x-x^*=-A^{-1}r$ và $\norm{A^{-1}}\le1/\mu$.
Áp Cauchy--Schwarz cho hiệu bilinear, rồi cộng sai số hai $q$ trong $\Gamma$.
\end{proof}

\paragraph{Cập nhật CG--IR: residual thật ở float64.}
Tại điểm tham số đóng băng, dựng một snapshot model double bằng cách nâng
\emph{các trọng số hiện tại} lên float64. Giữ $M$ và $c$ của bài toán đã chọn
cố định; đặt $A_{64}=J_{64}\T M J_{64}+\mu I$, $b_{64}=J_{64}\T c$.
Các cột thay thế trong $J'_{64}$ cũng phải được tính bằng model double;
không được chỉ cast cột float32 sang double rồi gọi đó là residual float64.
Cài đặt tính $r_k=b_{64}-A_{64}x_k$ bằng JVP/VJP của snapshot double, giải
hiệu chỉnh bằng CG trên toán tử native float32 ở dung sai trong $10^{-4}$,
rồi cập nhật $x_{k+1}=x_k+d_k$ bằng float64. Dừng chỉ khi
$\norm{r_k}\le10^{-10}\norm{b_{64}}$; tổng tối đa 200 vòng CG và tám lần
hiệu chỉnh mỗi hệ. Stagnation hay hết ngân sách là lỗi số của lần thử LM.
Đây là CG hiệu chỉnh trong khung mixed-precision IR \cite{carsonhigham2018},
không phải khẳng định dùng nguyên thuật toán LU/GMRES--IR của bài đó.

\paragraph{Bổ đề hội tụ có điều kiện.}
Giả sử $A_{64}\succ0$, $A_{32}=A_{64}+E$,
$\delta=\norm{A_{64}^{-1}}\norm E<1$ và mỗi hiệu chỉnh có residual
$\norm{r_k-A_{32}d_k}\le\eta\norm{r_k}$.
Với $\kappa=\norm{A_{64}}\norm{A_{64}^{-1}}$, ta có
\[
 \norm{x_{k+1}-x^*}\le
 \frac{\delta+\eta\kappa}{1-\delta}\norm{x_k-x^*}.
\]
Do đó $2\delta+\eta\kappa<1$ là một điều kiện đủ cho co sai số.
Nếu sai số toán tử có cỡ $u_{32}\norm{A_{64}}$ thì $\delta$ có cỡ
$\kappa u_{32}$, với các hằng số và giả thiết ổn định tương ứng;
riêng $\kappa u_{32}<1$ chưa là chứng minh đủ cho CG hiệu chỉnh hữu hạn.
\begin{proof}
Đặt $e_k=x_k-x^*$ và $s_k=r_k-A_{32}d_k$.
Khi đó $e_{k+1}=A_{32}^{-1}Ee_k-A_{32}^{-1}s_k$.
Dùng $\norm{A_{32}^{-1}}\le\norm{A_{64}^{-1}}/(1-\delta)$ và
$\norm{r_k}\le\norm{A_{64}}\norm{e_k}$.
\end{proof}
Khi residual cuối đạt dung sai, bổ đề sai số phía trên áp dụng cho $A_{64}$,
bất kể điểm khởi đầu. Ở hữu hạn dung sai, warm start có thể làm các bit cuối
khác nhau; nó không cho đẳng thức bitwise giữa các cách khởi đầu.
Nếu không hội tụ khi $\mu$ nhỏ thì đánh dấu lỗi số và thử LM tiếp theo;
không đặt cutoff $\mu$ tùy ý. Muốn kết luận $\mu^*$ không đổi phải đối chiếu
các lần lỗi số với $\varrho$ tham chiếu tại cùng $\mu$ trên checkpoint đó.

\paragraph{Warm start trong một lần áp dụng.}
Các nghiệm ở $\mu$, $\beta$ lân cận và trước swap chỉ làm điểm khởi đầu;
mọi nghiệm mới đều phải đạt residual double nói trên. Cache hệ trùng vế phải
và ma trận chỉ được dùng khi version của residual/metric/span vẫn khớp.
Mọi cache và snapshot được tạo mới trong lần áp dụng, không vào checkpoint
engine và không truyền sang lần áp dụng sau. $\varrho$ và kiểm cửa sổ vẫn
đo loss thật trên model native; không bỏ hay thay các cổng kiểm này.
Log cận sai số hai $q$ để đánh giá nhiễu $\Gamma$, không thêm cổng chấp nhận:
không có bảo đảm chung rằng sai số nhỏ hơn mọi $\Gamma$ khi $\Gamma$ gần $0$.

\paragraph{Swap dưới dạng matvec.}
Giữ $J$ tại mạng đóng băng; che các tọa độ incoming/outgoing của slot bỏ đi,
và đặt cột $J_v$ vào tọa độ outgoing của slot mới. Cột incoming mới bằng $0$
vì $B_v=0$. Đặt toán tử này là $J'$; CG giải $J'\T H_\beta J'+\mu^*I$.
Các tọa độ đã che có vế phải $0$ nên giữ hệ số $0$ dưới damping, tương đương
bỏ hẳn chúng. Với các site kề nhau, lần thay sau ghi đè đúng cột đã thay trước;
không trừ hai lần. Không dựng các cột native $J_u$ hay cập nhật Woodbury.

\subsection{T30: vòng tin cậy trong mỗi lần áp dụng, và một $\mu^*$ dùng chung}

\begin{namedthm}{T30}{vòng Levenberg--Marquardt bên trong}
Giả thiết: $L$ khả vi tại $\theta$ và $g\ne0$. Với mạng ReLU, điều này đúng hầu khắp nơi. Đặt $x(\mu)=(F+\mu)^{-1}g$, độ giảm dự báo $\mathrm{pred}(\mu)=g\T x-\tfrac12x\T Fx$ và $\varrho(\mu)=[L(\theta)-L(\theta-x(\mu))]/\mathrm{pred}(\mu)$. Nếu $L(\theta-x)$ không hữu hạn thì đặt $\varrho=-\infty$.
\begin{enumerate}[label=(\alph*)]
\item $\mathrm{pred}(\mu)\ge\tfrac12g\T x>0$ và $\varrho(\mu)\to1$ khi $\mu\to\infty$.
\item \emph{Vòng lặp.} Từ $\mu_0>0$ bất kỳ (mang sang từ lần áp dụng trước), lặp: nếu $\varrho(\mu)>0$ thì dừng với $\mu^*=\mu$; ngược lại đặt $\mu\leftarrow\nu\mu$ và $\nu\leftarrow2\nu$. Vòng lặp dừng sau hữu hạn bước. Số lần thử là $O\bigl(\sqrt{\log(\bar\mu/\mu_0)}\bigr)$, với $\bar\mu$ là ngưỡng mà từ đó $\varrho>0$.
\item \emph{Cổng kiểm.} Lần áp dụng bị bỏ qua khi và chỉ khi $g=0$ hoặc vòng lặp vượt ngân sách thử. Giá trị $\varrho$ không hữu hạn được tính là thất bại, không bao giờ là chấp nhận.
\item \emph{Một mô hình được chứng nhận cho cả VNA và GIR.} Với $\mu^*$:
\begin{itemize}
\item VNA dùng $\Delta=\norm{x(\mu^*)}$ và $\beta$ của \textsc{BetaRule} ở $\mu^*$;
\item GIR dùng $\ell_\beta$, $\rperp^\beta$ và $\Gamma$ ở cùng $\mu^*$.
\end{itemize}
Mỗi số hạng $\hat q^*$ trong $\Gamma$ (\Tref{T28}) là độ giảm của bước GN có damping $\mu^*$ trên một span. Với span hiện tại, bước đó chính là bước vừa được $\varrho>0$ chứng nhận.
\item Sau khi dừng, quy tắc Nielsen $\mu\leftarrow\mu^*\max(\tfrac13,1-(2\varrho-1)^3)$ cho giá trị khởi đầu của lần sau. Kết quả (b) không phụ thuộc $\mu_0$, nên $\mu_0$ không còn là tham số cần chọn.
\end{enumerate}
\end{namedthm}
\begin{proof}
(a) Ta có $g=(F+\mu)x$, nên $g\T x=x\T Fx+\mu\norm x^2\ge x\T Fx$. Suy ra $\mathrm{pred}\ge\tfrac12g\T x$, và $g\T x\ge\norm g^2/(\lambda_{\max}+\mu)>0$.
Viết $L(\theta-x)=L(\theta)-g\T x+r(x)$ với $|r(x)|\le\epsilon(\norm x)\norm x$, trong đó $\epsilon(t)\to0$ khi $t\to0$ (khả vi). Khi đó
\[\varrho-1=\frac{\tfrac12x\T Fx-r(x)}{\mathrm{pred}}.\]
Vì $x\T Fx\le\lambda_{\max}\norm x^2$ và $g\T x\ge\mu\norm x^2$, số hạng đầu bị chặn bởi $\lambda_{\max}/\mu$. Vì $\norm x\le\norm g/\mu$ và $\mathrm{pred}\ge\tfrac12\norm g^2/(\lambda_{\max}+\mu)$, số hạng thứ hai bị chặn bởi $2\epsilon(\norm x)(1+\lambda_{\max}/\mu)/\norm g$. Khi $\mu\to\infty$ thì $\norm x\to0$, nên cả hai số hạng tiến về $0$. Khi $\mu$ đủ lớn, $\theta-x$ đủ gần $\theta$ để $L$ hữu hạn (liên tục), nên cách đặt $\varrho=-\infty$ không ảnh hưởng tới giới hạn.
(b) Theo (a), tồn tại $\bar\mu$ sao cho $\varrho(\mu)>0$ với mọi $\mu\ge\bar\mu$. Sau $k$ lần thất bại, $\mu_k=\mu_0\,2^{k(k+1)/2}$, tăng nhanh hơn cấp số nhân. Do đó $k\le\lceil\sqrt{2\log_2(\bar\mu/\mu_0)}\,\rceil+1$.
(c) Định nghĩa. Phép so sánh $\varrho>0$ trả về sai khi $\varrho=\mathrm{NaN}$, nên NaN không được chấp nhận.
(d) Theo \Tref{T27}(a), $\hat q^*(\Jcal)=-\tfrac12b\T(\Jcal\T H_\beta\Jcal+\mu)^{-1}b$ là giá trị của mô hình tại bước GN có damping $\mu$ trên span $\Jcal$. Với $\Jcal=J$ và $\beta=0$, bước đó là $x(\mu)$.
(e) \cite{madsen2004}. Nếu $\varrho$ đủ lớn thì thừa số nhỏ hơn $1$.
\end{proof}
\begin{hequa}
\begin{itemize}
\item \Tref{T26}(d) được thay: $\mu$ không còn trễ một lần áp dụng. Đây là vòng lặp chuẩn của Moré \cite{more1978}, và \Tref{T26}(a)--(f) giữ nguyên với $\mu=\mu^*$.
\item Với T29, mỗi lần thử chỉ tốn một phân rã $m\times m$ và một forward. Ở $\beta=0$, một phân rã riêng của $L\T G_{\mathrm{out}}L$ dùng được cho mọi $\mu$.
\item Kiểm trên MLP ReLU nhỏ (50 bài, $\mu_0=10^{-8}$): vòng lặp luôn dừng, sau đúng 7 lần thử, phù hợp với cận ở (b).
\item \emph{Phần chưa được chứng nhận.} $\varrho$ chỉ chứng nhận mô hình dọc theo $x(\mu^*)$. Một swap của GIR đi theo hướng khác, nên hiệu ứng thật của nó vẫn được kiểm bằng các bước nhảy chính xác $j_L,j_{KD}$ (\Tref{T28}(c)). T30 bảo đảm $\Gamma$ được tính trong cùng mô hình và cùng damping mà bước GN đã được kiểm, không bảo đảm dấu của $\Gamma$ luôn đúng.
\end{itemize}
\end{hequa}

\subsection{T31: GIR tuần tự trong một mô hình}

\begin{namedthm}{T31}{một tuyến tính hóa cho mỗi lần áp dụng}
Cố định $\theta$, $J$, $H_\beta$, $e_\beta$, $\mu^*$ tại đầu lần áp dụng.
\begin{enumerate}[label=(\alph*)]
\item Các swap được nhận lần lượt $(u_1,v_1),(u_2,v_2),\dots$ Trạng thái mô hình sau $k$ swap gồm:
\begin{itemize}
\item span $J^{(k)}$, thu được từ $J$ bằng cách thay các cột $J_{u_i}$ bằng $J_{v_i}$;
\item độ dịch output $-\sum_{i\le k}f_{u_i}$;
\item residual $e^{(k)}=e_\beta-H_\beta\sum_{i\le k}f_{u_i}$.
\end{itemize}
Đặt $\hat q_k=\hat Q(-\sum_{i\le k}f_{u_i})+\min_x[\dots]$ như trong \Tref{T28}(a), và $\Gamma_k=\hat q_{k-1}-\hat q_k$. Khi đó $\hat q_0>\hat q_1>\dots$, và tổng lợi của mô hình bằng $\sum_k\Gamma_k$.
\item Swap slot chết không đổi hàm mạng và cập nhật $G_{\mathrm{out}}$ chính xác (\Tref{T29}(e)): mô hình sau swap là mô hình tuyến tính hóa đúng của mạng thật.
\item Swap slot sống đổi hàm mạng thật. Mô hình giữ tuyến tính hóa tại $\theta$, còn hiệu ứng trên mạng thật được kiểm bằng $j_L,j_{KD}\le0$, đo tuần tự trên mạng đã sửa. Lần áp dụng sau tuyến tính hóa lại, giống mỗi bước của LM.
\item Ở bản Gram tham chiếu, chi phí mỗi ứng viên: các lát Jacobian $J_u$ và $J_v$ (khoảng $|u|+|v|$ JVP; với slot chết thì $J_u=0$), cộng một cập nhật Woodbury hạng $|u|+|v|$ trên phân rã của $K+\mu^*$ (\Tref{T29}(d)). Không cần CG và không cần tuyến tính hóa lại mạng. Các giới hạn top-2 ứng viên và 4 slot sống mỗi site chỉ tồn tại vì chi phí. Profiling quyết định có bỏ được chúng không; chúng không phải một phần của định lý.
\end{enumerate}
\end{namedthm}
\begin{proof}
(a) Mỗi swap được nhận có $\Gamma_k>0$ theo quy tắc của \Tref{T28}(c). Áp \Tref{T28}(a) tại trạng thái $k-1$: gradient của $\hat Q$ tại độ dịch hiện tại là $e^{(k-1)}$, và phép thay cột là định nghĩa của $J^{(k)}$. Tổng là telescoping.
(b) \Tref{T29}(e).
(c) Định nghĩa. Bước nhảy được đo bằng forward trên mạng thật.
(d) \Tref{T29}(d) áp hai lần, một lần bớt và một lần thêm.
\end{proof}
\begin{hequa}
Code hiện tại tuyến tính hóa lại mạng thật cho \emph{từng} ứng viên. Đó là một mô hình khác với \Tref{T28}, đắt đỏ, và không nhất quán với $\ell_\beta$ (vốn tính tại $\theta$). T31 đưa code về đúng phát biểu của \Tref{T28}: trong một lần áp dụng, VNA và GIR dùng \emph{cùng} $J$, $H_\beta$, $e_\beta$, $\mu^*$.
\end{hequa}
\begin{kiem}
T31a (dãy $\hat q_k$ giảm khi chỉ nhận $\Gamma>0$, ba swap ngẫu nhiên mỗi bài): 0/200 vi phạm.
\end{kiem}

\subsection{R1: chuẩn hóa của chỉ số reach bị hỏng, và cách thay}

\begin{namedthm}{R1}{mẫu số $\norm w^2$ bị vài mẫu chi phối}
Với CE, mẫu $i$ có $e_i=(p_i-\mathbf 1_{y_i})/n$ và $H_i=(\operatorname{diag}p_i-p_ip_i\T)/n$. Khi đó $e_i\T H_i^+e_i=(1/p_{i,y_i}-1)/n$.
\end{namedthm}
\begin{proof}
$w=-\mathbf 1_{y_i}/p_{i,y_i}$ thỏa $H_iw=e_i$, vì $\operatorname{diag}(p)\mathbf 1_y-pp_y=\mathbf 1_yp_y-pp_y$. Mặt khác $e_i\perp\mathbf 1=\ker H_i$, nên $e_i\T H_i^+e_i=e_i\T w=(1-p_y)/(np_y)$.
\end{proof}
\begin{hequa}
\begin{itemize}
\item Một mẫu lớp cũ bị quên với $p_y\approx0$ đóng góp $1/p_y$ vào mẫu số. Trên baseline, $\norm w^2$ tăng từ $18$ lên $3\cdot10^{11}$ qua 20 task, nên reach $=B/\norm w^2$ tụt về $0$ chỉ vì vài mẫu, không phải vì mạng mất khả năng với tới. Thêm vào đó, CG của phép chẩn đoán có residual tới $1.9$.
\item Thay thế, không thêm tham số: log $B$ và $V$ theo đơn vị loss (độ giảm dự báo của bước GN có damping $\mu^*$), cùng tỉ số kiểm của \Tref{T25}. Bỏ reach khỏi các chỉ số chính.
\end{itemize}
\end{hequa}
\begin{kiem}
R1 so với $H^+$ tính bằng giả nghịch đảo: 0/200 vi phạm.
\end{kiem}

\section{Thuật toán}\label{sec:algorithm}

\subsection{Siêu tham số}

$\mu_{\mathrm{rel}}$ (damping), $K$ (khoảng áp dụng, tính theo iteration), $\lambda$ (theo \Tref{T24}), $r$ (số neuron ảo), $n_b$ (số bước bisection), $k_{\mathrm{lo}}$ (độ rộng khoảng log), $S$ (số swap tối đa mỗi lần áp dụng), $\rho_{\mathrm{rel}}=10^{-6}$ (ridge của $S_x$), \texttt{LIVE\_SWAP}$\in\{$\texttt{off} (mặc định), \texttt{measured}$\}$. \textbf{Cập nhật theo \cref{sec:closure}:} $\mu_{\mathrm{rel}}$, quy tắc $\lambda$ và $S$ được thay bằng trust region (\Tref{T26}) và quy tắc lợi ròng (\Tref{T28}); residual của GIR là $e_\beta$ (\Tref{T27}); \texttt{LIVE\_SWAP} được thay bằng \Tref{T28}(c). Bảng giả mã dưới đây giữ phiên bản đã chạy pilot để đối chiếu log. \textbf{Không có} $\beta$, $\gamma$, $\eta$ hay ngưỡng tự do nào khác: $\beta$ được tính (\Tref{T3}, \Tref{T11}), dung sai $\tau_h,\tau_g$ được tính từ CG (\Tref{T9}).

\begin{algorithm}[p]
\caption{VNA--GIR (một lần áp dụng, mỗi $K$ iteration; giữa các lần áp dụng, optimizer gốc chạy bình thường)}\label{alg:vna}
\small
\begin{algorithmic}[1]
\Require cửa sổ $\Wcal$ (cùng batch, train mode), tham số $\theta$, $A$ hiện tại, cờ \texttt{CIL}
\Statex \textbf{S0.} Mỗi iteration: $\theta\leftarrow\theta-\varepsilon m_t$ (optimizer gốc, không đổi)
\Statex \textbf{S1. Residual và tách} \hfill\Tref{T20}, \Tref{L5}
\State $e=e_N+\eO$ từ loss; $g=J\T e$, $g_O=J\T\eO$ \Comment{một backward}
\State $\xGN\leftarrow\mathrm{CG}(A_0,g)$; $z=HJ\xGN$; $z_O=z(\eO)$; $z_N=z-z_O$ \Comment{2 lần giải}
\State $\rperp=e-z$; với CIL, $\rperp^O=\eO-z_O$ \Comment{không giải thêm}
\Statex \textbf{S2. Chọn neuron ảo} \hfill\Tref{T6}, \Tref{T14}, \Tref{L9}
\If{đầu task hoặc đã 10 lần áp dụng từ lần làm mới trước}
  \State $A_0\leftarrow$ top-$r$ eigvec suy rộng của $(N_zN_z\T,S_x)$ \Comment{tiny\_V$(z)$}
  \State $\beta_0\leftarrow\textsc{BetaRule}(A_0)$; $y=HJG_{\beta_0}g$ \Comment{1 lần giải}
  \State $M_S=\sym(N_zN_{z_O}\T)+\sym(N_yN_{z_N}\T)$; $A\leftarrow$ eigvec suy rộng dương lớn nhất của $(M_S,S_x)$, tối đa $r$
\EndIf
\Statex \textbf{S3. Hiệu chỉnh} \hfill\Tref{T19}, \Tref{T1}--\Tref{T3}, \Tref{T11}
\State $u_R=K_Ez_O$, $u_N=K_Ez_N$, $V=\norm{J_E\T z}^2$, $V_O=\norm{J_E\T z_O}^2$
\If{không phải CIL (B5)} \State $\beta=0$; $\hat d\leftarrow\mathrm{CG}(A_0,J\T HK_Ez)$ \Comment{1 lần giải}
\Else
  \State $\drep\leftarrow\mathrm{CG}(A_0,J\T Hu_R)$; $d_N^0\leftarrow\mathrm{CG}(A_0,J\T Hu_N)$ \Comment{2 lần giải}
  \State $\beta\leftarrow\textsc{BetaRule}(A)$; $\hat d=\drep+\dN$
\EndIf
\Statex \textbf{S4. Chứng chỉ trên vector áp dụng} \hfill\Tref{T5}, \Tref{T9}, \Tref{T18}
\State $\hat h=g_O\T\hat d$, $\hat g=g\T\hat d$; $\hat\Delta$ từ residual cuối của các chuỗi CG; $\tau_h=\norm{g_O}\hat\Delta$, $\tau_g=\norm g\hat\Delta$
\If{$\hat g<-\tau_g$ \textbf{or} (CIL \textbf{and} $\hat h<-\tau_h$)} \State log ``bỏ qua''; \textbf{chuyển tới S6}
\EndIf
\Statex \textbf{S5. Áp dụng} \hfill\Tref{T10}(d), \Tref{T24}
\State $s^*=\hat g/\hat d\T(F+\mu)\hat d$; $\lambda=\min(K,s^*/\varepsilon)$ \Comment{sau S4, nên $s^*\ge0$ trừ dung sai}
\State $\theta'\leftarrow\theta-\varepsilon\lambda\hat d$ \emph{ngoài} momentum buffer; khôi phục buffer BN
\Statex \textbf{S6. GIR} (sau S5) \hfill\Tref{T7}, \Tref{T13}, \Tref{T17}, \Tref{T23}
\State $r\leftarrow\eO$ (CIL) hoặc $e$ (B5); dùng $\rperp(r)$ tính tại $\theta$
\State Ứng viên $\Vcal$: hướng $a_k$ = top-$k$ eigvec của $(N_{\rperp}N_{\rperp}\T,S_x)$; $\sigma\in\{\pm1\}$; bias $b\in\{0,-\operatorname{median}(\sigma Xa_k)\}$ (có BN: $\gamma_v=\sigma$, $\beta_v\in\{0,-\operatorname{median}(\sigma\hat x_k)\}$); bỏ ứng viên có $\varphi\equiv0$
\State $\ell(v)=\tfrac12\norm{J_v\T\rperp}^2/(\lambda_{\max}(J_v\T HJ_v)+\mu)$
\For{mỗi slot hợp lệ (A12) \emph{còn chết tại $\theta'$} (kiểm lại A5), tối đa $S$}
  \State $v=\argmax_{\Vcal}\ell$; đặt $a_v$ (và $b_v$ hoặc $\gamma_v,\beta_v$) vào slot; $B_v=0$
  \State có BN: running\_mean$=0$, running\_var$=1$; đặt trạng thái optimizer của slot bằng $0$
  \State cập nhật phần bù: chính xác theo \Tref{T13}(b), hoặc bản rẻ (giữ $\rperp$, hướng kế tiếp $S_x$-trực giao)
\EndFor
\If{\texttt{LIVE\_SWAP}$=$\texttt{measured}} \State thí nghiệm slot sống: swap chỉ khi $j_u^L\le0$, $j_u^{KD}\le0$, $\Delta L_{\mathrm{pred}}<0$, $\Delta KD_{\mathrm{pred}}<0$ (\Tref{T7}(i), \Tref{T16}); log kết quả thực
\EndIf
\Statex \textbf{S7. Log} \hfill\Tref{T3}, \Tref{T8}, \Tref{T21}
\State $\rho$, $V$, $V_O$, $\hat h$, $\hat g$, $\leak$, $\varepsilon^*$, $\beta$, nhánh của \textsc{BetaRule}, $\Phi_N$, số lần bỏ qua, residual CG, $\cos(\hat d,d_{\mathrm{exact}})$ ở điểm kiểm, defect $\delta_m$, $-\varepsilon\lambda\omega\theta\T\hat d$, số slot chết trước/sau S5, số swap, $\ell(v)$, gain
\Statex
\Function{BetaRule}{$A$} \hfill\Tref{T3}(e--g), \Tref{T11}
  \State $\chi^2=\frac1N\sum_i\chi^2(p_{t,i}\|p_{s,i})$ \Comment{dạng đóng, \Tref{L6}}
  \If{$\chi^2=0$ \textbf{or} $e_O=0$} \Return $0$ \Comment{\Tref{T3}(g)(iv)}
  \ElsIf{$V_O=0$} \Return $0$ \Comment{\Tref{T3}(g)(iii)}
  \EndIf
  \State $\varepsilon^*=V_O^2/\chi^2$; \textbf{if} $\leak(0)\le\varepsilon^*$ \textbf{then return} $0$
  \State $\beta_{\mathrm{suf}}=\chi^2\norm{u_N}_H^2/V_O^2$; bisection trên $\log\beta$ trong $[\beta_{\mathrm{suf}}10^{-k_{\mathrm{lo}}},\beta_{\mathrm{suf}}]$, $n_b$ bước, warm start; $\leak(\beta)=\norm{\Lambda C\dN}^2$ (một JVP)
  \State \Return đầu mút trên đã kiểm $\leak(\beta)\le\varepsilon^*$
\EndFunction
\end{algorithmic}
\end{algorithm}

\subsection{Hai biến thể theo bài toán}

\begin{tabularx}{\linewidth}{@{}lXX@{}}
\toprule
 & CIL (iCaRL, ResNet32) & B5 (mất plasticity) \\
\midrule
Residual GIR nhắm tới & $e_\beta$, metric $H_\beta$ (\Tref{T27}; bản pilot dùng $\eO$) & $e$ (trường hợp $\beta=0$ của \Tref{T27}) \\
$\beta$ & \textsc{BetaRule} (có lớp cũ) & $0$ (không có lớp cũ) \\
Chứng chỉ ở S4 & $\hat g$ và $\hat h$ & $\hat g$ \\
GIR & không kích hoạt với cấu hình mặc định (0\% đơn vị chết; slot sống tắt) & cơ chế chính (nhiều đơn vị chết) \\
Số lần giải mỗi lần áp dụng & $6+n_b$ (tốt nhất $4$); $9+2n_b$ khi làm mới $A$ & $2$ \\
\bottomrule
\end{tabularx}

\subsection{Chi phí}

\begin{tabularx}{\linewidth}{@{}Xccc@{}}
\toprule
Phần (số lần giải hệ tuyến tính, trường hợp xấu nhất) & CIL & CIL, làm mới $A$ & B5 \\
\midrule
$z$, $z_O$ & 2 & 2 & 1 \\
\textsc{BetaRule} cho $A_0$ & --- & $2+n_b$ & --- \\
$y=HJG_{\beta_0}g$ & --- & 1 & --- \\
$\drep$, $d_N^0$ & 2 & 2 & 1 \\
bisection cho $A$ đang dùng (khi $\leak(0)>\varepsilon^*$), có thể $+1$ cho $d_N(\beta_{\mathrm{suf}})$ & $2+n_b$ & $2+n_b$ & --- \\
\midrule
\textbf{Tổng} & $\mathbf{6+n_b}$ & $\mathbf{9+2n_b}$ & $\mathbf 2$ \\
\bottomrule
\end{tabularx}

\medskip
Chứng chỉ, $\leak$, $\varepsilon^*$, $\chi^2$, $\rperp$, $\tau_h$, $\tau_g$ không cần giải thêm. GIR slot chết: $\ell(v)$ không cần giải; deflation chính xác tốn $q+1$ lần giải mỗi slot thêm. Chi phí phải đo trên phần cứng thật và đưa vào quy tắc chọn $K$.

\subsection{Mỗi dòng của \cref{alg:vna} và cơ sở của nó}

\begin{longtable}{@{}p{0.24\linewidth}p{0.30\linewidth}p{0.40\linewidth}@{}}
\toprule
Dòng / bước & Cơ sở & Bảo đảm hoặc vai trò \\
\midrule
\endhead
S1: $\xGN$, $z$, $\rperp$ & \Tref{T20}, \Tref{T1}(d) & Tách chính xác; dạng không cần range, đúng cả khi softmax bão hòa. \\
S2: tiny\_V$(z)$ & \Tref{T6}(c); A4$'$ & Tối ưu bậc một chính xác ở $\beta=0$ trong vùng mặt nạ. \\
S2: $A$ từ $(M_S,S_x)$ & \Tref{T6}(a, b, e); \Tref{L9} & Tối ưu có điều kiện theo $\beta_0$; $S_x\succ0$ nhờ ridge. \\
S3: $\hat d$ ở B5 & \Tref{T19}, \Tref{T2}(a) & Phép chiếu chính xác phần ảo; $g\T d=V\ge0$. \\
S3: $\drep$, $\dN$ & \Tref{T1}(a, b) & Phân rã duy nhất; $\beta$ chỉ tác động lên phần học mới. \\
\textsc{BetaRule} & \Tref{T3}(c, e--h), \Tref{T11}, \Tref{L6} & $\beta$ trả về luôn có chứng chỉ $h\ge0$; là nhân tử KKT khi $\leak(0)>\varepsilon^*$. \\
S4: bỏ qua & \Tref{T5}(d, f, g), \Tref{T9}(b, e) & Không làm tăng $L$, $L_{KD}$ ở bậc một (sai khác $\lambda\tau_g$, $\varepsilon\lambda\tau_h$); không bỏ nhầm. \\
S5: $\lambda$ & \Tref{T24}; A13 & Cực đại mô hình GN; $\lambda_{\mathrm{faithful}}=K$ chỉ là giả thiết diễn giải. \\
S5: ngoài momentum & \Tref{T10}(d); A9 & Giữ dấu của chứng chỉ. \\
S6: kiểm lại A5 & \Tref{T7}(f), \Tref{T13}; A5 & Bảo toàn hàm chính xác chỉ khi slot còn chết tại $\theta'$. \\
S6: chọn $v$ theo $\ell$ & \Tref{T7}(b, d), \Tref{T23} & Lợi $\ge\ell(v)\ge0$; hệ số $\mu/(c_H+\mu)$ so với tốt nhất trong $\Vcal$. \\
S6: sau S5 & \Tref{T7}(g) & Swap làm $K$ lớn lên, không giao hoán với VNA. \\
S6: BN, optimizer & \Tref{T17}, \Tref{T13}(d) & Neuron mới không chết vì $\beta_{BN}$; không có dịch chuyển không gradient. \\
S6: slot sống & \Tref{T7}(i), \Tref{T16} & Chỉ chứng chỉ tức thời; dự báo không đáng tin ở damping thật, nên tắt mặc định. \\
S7: $\rho$, defect & \Tref{T21}, \Tref{T8} & Nghĩa của chẩn đoán được chứng minh. \\
\bottomrule
\end{longtable}

\section{Ablation là trường hợp riêng; baseline ngoài}\label{sec:ablation}

\subsection{Một họ cập nhật}

Gọi $\mathcal U(\lambda,\kappa,\beta,\text{chọn},S)$ là cập nhật: mỗi iteration optimizer gốc chạy; mỗi $K$ iteration cộng $-\varepsilon\lambda\,P(\kappa z)$ với phân rã bảo vệ của \Tref{T1}, rồi GIR với tối đa $S$ swap. Mọi arm dưới đây là thành viên của họ này; chúng chỉ khác một thành phần, nên so sánh cặp đo đúng đóng góp của thành phần đó.

\begin{tabularx}{\linewidth}{@{}llX@{}}
\toprule
Arm & Thiết lập & Câu hỏi được trả lời \\
\midrule
baseline (iCaRL / BP / CBP) & $\lambda=0$ & mốc \\
\texttt{kernel\_self} & $\kappa=K$ (thay $K_E$) & kernel ảo có hơn kernel của chính mạng không? (\Tref{T19}(c)) \\
VNA & $\kappa=K_E$ (tiny\_V), $\beta=0$ & tác dụng của phần ảo không bảo vệ \\
VNA + bảo vệ & $\beta=\beta_{\mathrm{cert}}$ & tác dụng của số hạng bảo vệ \\
A ngẫu nhiên & $\kappa=K_E(A_{\mathrm{rand}})$, cùng vết & tiêu chí grow có quan trọng không? \\
không swap & $S=0$ & tác dụng của GIR (B5) \\
GIR, $A$ ngẫu nhiên & swap với hướng ngẫu nhiên, cùng slot & ``reset thành gì'' có quan trọng không? \\
\bottomrule
\end{tabularx}

\medskip
Hai so sánh quyết định là \textbf{VNA so với \texttt{kernel\_self}} và \textbf{tiny\_V so với $A$ ngẫu nhiên}: cả hai chỉ khác kernel đưa vào cùng một phép chiếu, nên kết quả không thể đọc là ``một optimizer khác đã thắng''.

\subsection{Baseline ngoài}

\begin{tabularx}{\linewidth}{@{}lX@{}}
\toprule
Baseline & Vai trò \\
\midrule
GN lười ($\xGN$) & đổi optimizer gốc; có sẵn miễn phí từ bước S1 (\Tref{T22}); không phải ablation \\
K-FAC / Shampoo \cite{martens2015kfac,gupta2018shampoo} & preconditioner độ cong của chính mạng, cùng wall-clock; báo $\cos(d,d_{\mathrm{KFAC}})$ \\
C-Flat \cite{bian2024cflat} & plug-in optimizer cùng setting PyCIL; mốc hiệu quả và chi phí \\
CBP \cite{dohare2024}, ReDo \cite{sokar2023redo} & reset ngẫu nhiên với cùng tiêu chí chọn đơn vị \\
Reset khởi tạo kiểu GradMax \cite{evci2022gradmax} & cùng slot, hướng theo gradient trên $e$ thay vì $\rperp$; trả lời phản biện ``GIR = GradMax áp vào ReDo'' \\
GPM / EFC \cite{saha2021gpm,magistri2024efc} & bảo vệ lớp cũ không có dung lượng ảo (nếu CIL là track chính) \\
\bottomrule
\end{tabularx}

\subsection{Bằng chứng toy hiện có (chỉ để định hướng)}

Trên toy nhỏ, float64, 12 seed, ghép cặp (\st{empirical}): ở bài kiểu CIL, VNA hơn SGD $+0.22$ điểm accuracy trung bình (10/12 seed, kiểm định dấu hai phía $p\approx0.04$), hơn \texttt{kernel\_self} $+0.28$ điểm (8/12, $p\approx0.39$), hơn GN $+0.17$ điểm (8/12, mức nhiễu). Ở bài plasticity không có đơn vị chết, VNA \emph{không} hơn \texttt{kernel\_self} (4/12) hay GN (3/12); tổ hợp $\xGN+d$ có loss cuối trung bình tốt nhất. Các arm toy dùng $\lambda$ chuẩn hóa cùng norm, không phải quy tắc của \Tref{T24}. Những số này không thay đổi lý thuyết và không trả lời câu hỏi quyết định trên ResNet32 và B5.

\section{Khẳng định đã rút lại, giới hạn và câu hỏi mở}\label{sec:limits}

\subsection{Các khẳng định đã bị thay thế}

Những khẳng định sau xuất hiện ở các bản trước và \textbf{không còn được dùng}:
\begin{enumerate}
\item \textbf{Núm $\gamma$} (trọng số riêng cho phần sửa): $\gamma=1$ chính là VNA hiện tại, còn $\gamma\ne1$ phá bảo đảm giảm loss (\Tref{T5}(c)).
\item \textbf{Dạng $\norm{Cd+\eta e_{\mathrm{old}}}^2$}: sai dấu và sai metric; thay bằng đích $C\drep$ (\Tref{T1}).
\item \textbf{Phạt về đích $0$ hay metric $H+\beta H_{KD}$ làm đề xuất chính}: làm mất phần sửa (\Tref{T3}(i)).
\item \textbf{Giới hạn cứng $Cd=0$}: sai; đúng là $\Lambda C(d-\drep)=0$ (\Tref{T2}(c)).
\item \textbf{Số hạng bậc một $+\varepsilon e_O\T Cd$}: dấu đúng là $-\varepsilon$ (\Tref{T3}(d)).
\item \textbf{Hessian KD dùng xác suất teacher}: phải dùng student (\Tref{L8}).
\item \textbf{``VNA-NS là chiếu của growth kiểu DER'' / ``growth với hàm cũ đóng băng''}: bỏ hẳn; không có định lý nào. Phát biểu hình thức duy nhất là \Tref{T2}(d).
\item \textbf{Độ khớp tăng đơn điệu theo $\beta$}: sai; chỉ $\Phi_N$ đơn điệu (\Tref{T4}).
\item \textbf{Các arm $\beta\in\{0,3,\infty\}$}: thay bằng $\beta$ có chứng chỉ.
\item \textbf{``Một Lagrangian duy nhất $\Psi$'' như một định lý}: rút lại; VNA và swap chỉ có cùng \emph{dạng} hàm chiếu, không cùng thang đo (\Tref{T7}(a)).
\item \textbf{$h$ đơn điệu, bisection trên $h$}: sai (\Tref{T3}(h)); bisection trên $\leak$.
\item \textbf{Chọn neuron swap theo toàn bộ $e$}: bị dữ liệu bác bỏ; đúng là $\rperp$ (\Tref{T7}).
\item \textbf{Swap giao hoán với VNA}: sai (\Tref{T7}(g)).
\item \textbf{$\beta_{\mathrm{suf}}$ làm mặc định}: quá bảo thủ (trung vị gấp 29 lần $\beta_{\mathrm{cert}}$).
\item \textbf{``$V_O=0\Rightarrow$ bỏ qua''}: sai; $\beta=0$ có chứng chỉ (\Tref{T3}(g)(iii)).
\item \textbf{``Slot sống không có chứng chỉ vì $L_3$ không kiểm được''}: sai, $L_3$ có dạng đóng (\Tref{L10}); nhưng \textbf{``slot sống có chứng chỉ trên cửa sổ''} cũng là nói quá; quy tắc Net$_{KD}$ mô hình hồi phục chỉ trên KD, trái A10. Slot sống nay là \st{measured}, tắt mặc định (\Tref{T16}).
\item \textbf{$L_3\le0.708(1+T^{-3})/N$}: đúng nhưng lỏng 2.6 lần; thay bằng $c_3\approx0.2722$ chặt.
\item \textbf{Quy tắc chặt $\hat h<0$ khi dùng CG}: bỏ nhầm khoảng một nửa số bước khi $V_O=0$; thay bằng dung sai (\Tref{T9}(e)).
\item \textbf{``$K_E=0$ cho bước GN''}: sai; $K_E\to K$ cho $x_K$, không phải $\xGN$ (\Tref{T19}).
\item \textbf{``$\rho_{\mathrm{proj}}$ là phần tiến bộ của mạng mở rộng đến từ neuron ảo''}: sai; $\rho$ là tỉ số ảo trên nền của mạng cố định (\Tref{T21}).
\item \textbf{``$\xGN+d$ không phải phép chiếu của bất kỳ động lực nào''}: sai; thay bằng phát biểu định lượng (\Tref{T22}).
\item \textbf{``$\rperp$ nằm ngoài tầm với; chỉ GIR tác động lên nó; vai trò tách bạch''}: sai khi $\mu>0$ (\Tref{T20}(b)).
\item \textbf{Cap bước $\mathrm{lr}\norm g/\norm{\hat d}$ và lưới $\mu_{\mathrm{rel}}$}: thay bằng trust region (\Tref{T26}); cap cũ chặn 99.5\% lần áp dụng trên B5, nên T24 gần như không được dùng.
\item \textbf{Residual GIR $r=\eO$ trong CIL}: không suy ra từ bài toán nào; thay bằng $e_\beta$ (\Tref{T27}).
\item \textbf{Hai cơ chế slot chết / slot sống}: gộp thành một quy tắc lợi ròng (\Tref{T28}).
\item \textbf{``SGD $+\lambda d$ đã là phép chiếu rời rạc của một động lực mở rộng duy nhất''}: chỉ là \st{interpretation} xấp xỉ.
\end{enumerate}

\subsection{Giới hạn}

\begin{enumerate}
\item \textbf{Cục bộ.} Mọi bảo đảm là bậc một (hoặc bậc hai với hằng số A6 không kiểm được), cho từng lần áp dụng, trên loss của cửa sổ. Không có bảo đảm cho cả quỹ đạo (trừ khi defect đo được nhỏ, \Tref{T8}), về tổng quát hóa hay accuracy trên test.
\item \textbf{Tối ưu của bước chọn chỉ trong một vùng mặt nạ} (A4$'$); tối ưu toàn cục là A4-glob.
\item \textbf{Bổ đề chiếu lồng nhau.} VNA chỉ phân bổ lại tốc độ trong $\range(K)$; với chuẩn hóa cùng norm, một lần áp dụng không hơn một bước SGD ở bậc một. Lợi ích, nếu có, đến từ thông tin độ cong.
\item \textbf{Bổ đề nội suy.} Khi model fit hết cửa sổ ($e\to0$) thì $z,\rperp\to0$, nên mọi thành phần cùng tắt; lợi ích dự kiến tập trung ngay sau khi chuyển task.
\item \textbf{GIR trong CIL} không hoạt động với cấu hình mặc định.
\item \textbf{$\beta_{\mathrm{cert}}$ bảo thủ}: trung vị gấp 4.5 lần $\beta$ nhỏ nhất cho $h\ge0$ trên lưới.
\item \textbf{Tỉ lệ bỏ qua}: khoảng 4\% trên toy với $\beta_{\mathrm{cert}}$; tới 42\% với VNA không bảo vệ trong toy kiểu CIL. Phải log trên ResNet32.
\item \textbf{Chi phí} CIL tới $6+n_b$ lần giải mỗi lần áp dụng.
\item \textbf{Kiểm số chỉ trên mô hình nhỏ, float64.} Ở float32 phải kiểm lại các đồng nhất chính (T1, T3a, T3e) tại điểm kiểm.
\item \textbf{A10 không mô tả hồi phục thật}; mọi diễn giải ``lợi về sau'' của GIR là \st{interpretation}; chỉ bảo toàn hàm và chứng chỉ là chính xác.
\item \textbf{Slot hợp lệ} (A12) loại kênh đi vào phép cộng residual.
\item \textbf{T25 chỉ là chặn trên}: nó giải thích vì sao VNA một mình thất bại khi nhiều đơn vị chết, nhưng không bảo đảm VNA cộng thêm lợi khi đã có GIR.
\end{enumerate}

\subsection{Câu hỏi thực nghiệm mở}\label{sec:open}

Lý thuyết không trả lời được, và phải được trả lời bằng thí nghiệm có đăng ký trước:
\begin{enumerate}
\item \textbf{$K_E$ so với $K$:} VNA có hơn \texttt{kernel\_self} trên ResNet32 (CIL) và trên B5 không? Đây là so sánh quyết định của toàn bộ phương pháp.
\item \textbf{tiny\_V so với $A$ ngẫu nhiên:} ở phép đo nhanh CIL trước đây, tiny\_V thua random ở 7/9 dòng. Nếu điều này lặp lại khi train thật, phải bỏ khẳng định ``tiêu chí grow'' và chỉ giữ khẳng định về phép chiếu.
\item \textbf{GIR so với reset ngẫu nhiên và reset kiểu GradMax} trên B5 với cùng slot, cùng thời điểm.
\item \textbf{Chi phí thực tế} (GPU-giờ, hệ số so với baseline) và tỉ lệ bỏ qua trên ResNet32.
\item \textbf{Bổ sung của VNA khi có GIR} (\Tref{T25}): khoảng cách (VNA+GIR) $-$ GIR phải tăng theo độ rộng mạng; kiểm bằng B5 với 5, 20, 100 đơn vị.
\item \textbf{Dự đoán kiểm được:} $\rho$ giảm trong mỗi task; lợi tập trung sau khi chuyển task; không có lợi trên bài toán đã nội suy (Phase 0 là đối chứng âm được dự đoán).
\end{enumerate}


\appendix
\section{Kiểm số và lịch sử phản biện}\label{app:verification}

\subsection{Bộ kiểm số}

\begin{tabularx}{\linewidth}{@{}lXl@{}}
\toprule
Script & Phạm vi & Kết quả \\
\midrule
\texttt{final\_check\_v3.py} & L5--L10, T0--T18; float64; chạy \texttt{python3 -I} từ thư mục trống, khoảng 45 giây & 144/144 đạt \\
\texttt{check\_unified\_gn.py} & T19--T24 (U1--U9) & 15/15 đạt \\
\texttt{live\_slot\_198.py} & thí nghiệm 198 cặp slot sống của phản biện vòng 3 & xem \Tref{T16} \\
phản biện vòng 5 (độc lập) & T19--T24 viết lại từ đầu, gồm trường hợp suy biến & 25/25 đạt, 8 nhân chứng chống khẳng định cũ \\
\bottomrule
\end{tabularx}

\medskip
Ba loại phép kiểm phụ được phân biệt:
\begin{itemize}
\item \textbf{NEG (29): đối chứng đột biến.} Cùng phép kiểm áp cho một công thức \emph{cố ý sai}; chỉ đạt khi công thức sai bị loại với biên độ rõ. Ví dụ: Hessian KD dùng teacher (lệch 0.69), giới hạn $Cx=0$ (cách $\ge0.265$), $s=J_v\T r$ thay $J_v\T\rperp$ (lệch trung vị 1.95), hằng số $L_3$ nhân $0.1$, bỏ hệ số $1/\mu$ trong dung sai CG (vi phạm 26/60).
\item \textbf{WIT (22): nhân chứng tồn tại.} Chứng tỏ một hiện tượng có xảy ra, nên một khẳng định phủ định là đúng hoặc một quy tắc là cần thiết. Ví dụ: $h$ không đơn điệu (62/300), VNA thuần làm KD tăng (102/300), đơn vị chết sống lại sau bước hiệu chỉnh (18/40).
\item \textbf{[structural]:} phép kiểm đúng theo cách dựng; ghi rõ để không bị đọc như bằng chứng độc lập (T12c, T13c, T16e).
\end{itemize}

\subsection{Lịch sử phản biện}

Mỗi vòng do một agent độc lập thực hiện, không thấy quá trình dẫn xuất, tự dẫn lại bằng tay và viết mã kiểm riêng.

\begin{longtable}{@{}p{0.11\linewidth}p{0.2\linewidth}p{0.63\linewidth}@{}}
\toprule
Vòng & Kết quả & Phát hiện chính đã được sửa \\
\midrule
\endhead
1 & 55 mục: 34 đúng, 13 cần sửa, 6 sai, 2 không kiểm được & ``Một nguyên lý $\Psi$'' chỉ là phép tương tự; giới hạn $\beta\to\infty$ sai ($Cd=0$ thay vì $W^{1/2}Cd=0$); quy tắc swap mâu thuẫn giữa các tài liệu; ``độ khớp đơn điệu'' sai; lỗi dấu (giá swap, số hạng KD bậc một, hệ số $\eta$); Hessian KD dùng teacher; nhiều phép kiểm rỗng. \\
2 & 48 mục: 30 đúng, 13 cần sửa, 3 sai, 2 không kiểm được & ``$V_O=0\Rightarrow$ bỏ qua'' sai; slot sống \emph{có} hằng số $L_3$ dạng đóng; GIR: kiểm lại slot chết sau hiệu chỉnh, đặt lại optimizer, quy tắc nhiều slot, dấu/bias ReLU; A4 sai với ReLU/BN (thay bằng A4$'$); $S_x$ cần ridge; chi phí đếm thiếu. \\
3 & 46 mục: 36 đúng, 8 cần sửa, 1 nói quá, 1 sai & Quy tắc slot sống Net$_{KD}<0$ sai logic dưới A10 (lệch dấu 104/198 cặp); ``có chứng chỉ trên cửa sổ'' là nói quá; $L_3$ lỏng 2.6 lần; BN của slot ($\gamma,\beta$); dung sai CG cho $V_O=0$. \\
4 & không còn mục sai & Chỉnh câu chữ: đếm lần giải ở bước 4; T9(b) dưới quy tắc có dung sai; T16(c) quy về A10; hai phép kiểm làm hai phía. \\
5 (T19--T24) & 21 mục: 11 đúng, 5 cần sửa, 4 sai/nói quá, 1 giả thiết & Cách đọc $\rho$ (tỉ số, không phải tỉ phần; không phải tốc độ của mạng mở rộng); ``$\xGN+d$ không phải phép chiếu của bất kỳ động lực nào'' sai; ``$\rperp$ ngoài tầm với, vai trò tách bạch'' sai khi $\mu>0$; ``SGD $+\lambda d$ đã là một phép chiếu duy nhất'' chỉ là diễn giải. Tất cả đã sửa trong \cref{sec:unified}. \\
\bottomrule
\end{longtable}

\subsection{Bảng trạng thái tóm tắt}

\begin{longtable}{@{}p{0.13\linewidth}p{0.52\linewidth}p{0.29\linewidth}@{}}
\toprule
ID & Khẳng định & Trạng thái \\
\midrule
\endhead
L5--L10 & range; $\chi^2$; $\Lambda$; Hessian student; ridge; $L_3$ chặt & \st{proven} \\
T1 & dạng đóng duy nhất, phân rã, dạng không cần range & \st{proven} \\
T2 & $\beta=0$ là VNA, $g\T d=V$; giới hạn cứng $\Lambda Cx=0$; chiếu lồng nhau & \st{proven} \\
T3 & sửa đúng chiều; rò rỉ đơn điệu; chứng chỉ; KKT; suy biến & \st{proven}; $h$ đơn điệu: \textbf{sai} \\
T4 & $\Phi_N$ không giảm; cái giá $\le\beta\leak(0)$ & \st{proven}; độ khớp đơn điệu: \textbf{sai} \\
T5 & cận chuẩn; không có $\gamma$; bậc một/bậc hai có quy tắc bỏ qua & \st{proven} (bậc hai: A6) \\
T6, T14 & lợi $=\tr(A\T M_SA)$; eigvec suy rộng; tiny\_V tối ưu ở $\beta=0$ & \st{proven} (A4$'$) \\
T7 & Schur theo $\rperp$; chọn theo $\ell$; slot chết bảo toàn hàm; không giao hoán & \st{proven}; diễn giải lợi về sau: A10 \\
T7(i), T12 & slot sống: $j_u$ đo chính xác; cận Taylor; chứng chỉ tức thời & \st{proven} / \st{measured} \\
T16 & điểm khai triển $f-f_u$; dự báo đáng tin ở damping thật & \st{proven}; \textbf{sai} (\st{empirical}) \\
T7(j) & GIR không hoạt động trong CIL (mặc định) & \st{empirical} \\
T8 & Grönwall theo defect thực hiện & \st{proven} (A7) \\
T9, T18 & sai số CG; chứng chỉ trên vector áp dụng; dung sai & \st{proven} (A8) \\
T10 & BN train mode; momentum; weight decay; augmentation & \st{proven} (A1, A9) \\
T11 & quy tắc $\beta$ và các nhánh suy biến & \st{proven} \\
T13, T17 & nhiều slot; optimizer; BN của slot & \st{proven} \\
T15 & kiểm ở mức thuật toán & \st{empirical} \\
T19--T21, T23, T24 & phân rã chiếu; tách residual; tỉ số $\rho$; lợi swap; cỡ bước & \st{proven} \\
T22 & đạo hàm của GN mở rộng; phát biểu định lượng & \st{proven}; phát biểu định tính cũ: \textbf{sai} \\
--- & optimizer gốc là phần $K$ của cùng phép chiếu & \st{interpretation} \\
--- & $\lambda_{\mathrm{faithful}}=K$ & \st{assumption} (A13) \\
--- & nguyên lý Gauss, điều khiển ưu tiên, OMP, BUG & \st{interpretation} \\
\bottomrule
\end{longtable}

\section{C3: tham số đối thủ cố định, không chọn cấu hình}\label{app:c3-defaults}

C1 bị loại khỏi kế hoạch thực nghiệm. C3 chạy CBP, ReDo, S\&P và UPGD
với ba seed $\{0,1,2\}$ (12 job), trên cùng thiết lập iCaRL noBN,
thứ tự superclass và 20 task. Không quét tham số, không chọn bằng seed~0,
và không cần tệp \texttt{best\_config.json} của C1.
Theo quyết định của người thiết kế thực nghiệm, mọi phương pháp dùng
learning rate iCaRL/PyCIL $0.1\times0.5=0.05$, ở cả task đầu và các task sau,
với lịch giảm learning rate chung. Đây không phải một nửa learning rate
riêng trong thí nghiệm RL hoặc Adam của từng bài gốc.

\begin{tabularx}{\linewidth}{@{}lXX@{}}
\toprule
Phương pháp & Tham số cố định & Nguồn cụ thể \\
\midrule
CBP & Replacement rate $10^{-5}$; maturity threshold $1000$;
utility contribution; decay $0.99$ & Cấu hình incremental CIFAR của tác giả
và hàm khởi tạo thí nghiệm, commit \texttt{a6b7958}. \\
ReDo & $\tau=0.025$; chu kỳ $1000$ bước optimizer; tối đa $64$ mẫu
để tính score & Bảng~1, Phụ lục~A của Sokar và cộng sự (ICML 2023),
thiết lập DQN mặc định. Bảng phân biệt $0.025$ cho default và $0.1$ cho các
thiết lập khác; ở đây cố định trường hợp default. \\
S\&P & Shrink $0.4$; perturb $0.1$ & Giá trị mặc định của
\texttt{warm\_start/run.py}, commit \texttt{896176b}. Perturb nhân một
khởi tạo mạng mới, không phải nhiễu Gaussian phương sai đơn vị. \\
UPGD & $\beta_{\rm utility}=0.999$; $\sigma=0.001$;
weight decay $0.001$ & Constructor được tác giả khuyến nghị trong README,
commit \texttt{b75e90a}; chạy lớp native \texttt{FirstOrderGlobalUPGD}. \\
\bottomrule
\end{tabularx}

\paragraph{Phạm vi chuyển sang iCaRL.}
Giữ backbone, bộ nhớ, loss, số epoch và đánh giá chung của benchmark.
ReDo đếm số bước optimizer thực tế, dùng tối đa 64 activation của minibatch
hiện tại trước bước cập nhật; giữ bộ đếm và lần reset cuối trong trạng thái
learner để tiếp tục chính xác qua task/checkpoint. Giữ bản chuyển NCHW
hiện có, gồm khởi tạo Xavier cho trọng số vào, đặt trọng số ra bằng không
và xóa trạng thái optimizer tương ứng. S\&P giữ lịch can thiệp theo task
của adapter. UPGD dùng công thức cập nhật của lớp native đã pin, không
chép hệ số cập nhật khác trong ví dụ README, và không thêm momentum SGD.
Đây là các bản chuyển sang iCaRL, không phải tái lập toàn bộ giao thức gốc.

\paragraph{Truy xuất nguồn và kết quả.}
Mỗi protocol lưu đủ dictionary tham số với nhãn
\texttt{original\_defaults\_no\_search\_v1}. Không đổi nhãn checkpoint
C1 cũ thành cấu hình mới; loại kết quả đã chọn qua C1 khỏi bảng C3.
Danh sách đường dẫn nguồn bất biến, khác biệt giữa các mặc định của
bài gốc và giới hạn của bản chuyển được ghi trong
\texttt{docs/vna/C3\_DEFAULTS\_APPENDIX.md} của repository.
Bài ReDo: \url{https://proceedings.mlr.press/v202/sokar23a/sokar23a.pdf}.


\section*{Ghi chú về tài liệu tham khảo}
Các mục không có chú thích đã được kiểm tiêu đề, venue và năm trong các vòng survey của dự án (thường kèm URL arXiv hoặc trang hội nghị). Danh sách tác giả đầy đủ của một số mục được ghi theo trí nhớ và phải đối chiếu lại trước khi nộp. Các mục có chú thích chưa được kiểm tại nguồn.

\begin{thebibliography}{99}\small

\bibitem{lubich2008book} C.~Lubich. \emph{From Quantum to Classical Molecular Dynamics: Reduced Models and Numerical Analysis}. EMS, 2008.\unv

\bibitem{kochlubich2007} O.~Koch, C.~Lubich. Dynamical low-rank approximation. \emph{SIAM J. Matrix Anal. Appl.}, 2007.\unv

\bibitem{bruna2024neuralgalerkin} J.~Bruna, B.~Peherstorfer, E.~Vanden-Eijnden. Neural Galerkin schemes with active learning for high-dimensional evolution equations. \emph{J. Comput. Phys.}, 2024. arXiv:2203.01360.

\bibitem{ceruti2022bug} G.~Ceruti, J.~Kusch, C.~Lubich. A rank-adaptive robust integrator for dynamical low-rank approximation. \emph{BIT}, 2022.\unv

\bibitem{schotthofer2022dlrt} S.~Schotthöfer, E.~Zangrando, J.~Kusch, G.~Ceruti, F.~Tudisco. Low-rank lottery tickets: finding efficient low-rank neural networks via matrix differential equations. \emph{NeurIPS}, 2022.

\bibitem{amari1998} S.~Amari. Natural gradient works efficiently in learning. \emph{Neural Computation}, 1998.\unv

\bibitem{martens2020} J.~Martens. New insights and perspectives on the natural gradient method. \emph{JMLR}, 2020.\unv

\bibitem{muller2023engd} J.~Müller, M.~Zeinhofer. Achieving high accuracy with PINNs via energy natural gradient descent. \emph{ICML}, 2023.

\bibitem{jacot2018ntk} A.~Jacot, F.~Gabriel, C.~Hongler. Neural tangent kernel: convergence and generalization in neural networks. \emph{NeurIPS}, 2018.\unv

\bibitem{martens2015kfac} J.~Martens, R.~Grosse. Optimizing neural networks with Kronecker-factored approximate curvature. \emph{ICML}, 2015.\unv

\bibitem{gupta2018shampoo} V.~Gupta, T.~Koren, Y.~Singer. Shampoo: preconditioned stochastic tensor optimization. \emph{ICML}, 2018.\unv

\bibitem{purifyingshampoo} Purifying Shampoo: investigating Shampoo's heuristics by decomposing its preconditioner. arXiv:2506.03595, 2025.

\bibitem{verbockhaven2024tiny} M.~Verbockhaven, T.~Rudkiewicz, S.~Chevallier, G.~Charpiat. Growing tiny networks: spotting expressivity bottlenecks and fixing them optimally. \emph{TMLR}, 2024. arXiv:2405.19816.

\bibitem{evci2022gradmax} U.~Evci, B.~van Merriënboer, T.~Unterthiner, M.~Vladymyrov, F.~Pedregosa. GradMax: growing neural networks using gradient information. \emph{ICLR}, 2022. arXiv:2201.05125.

\bibitem{wu2020firefly} L.~Wu, B.~Liu, P.~Stone, Q.~Liu. Firefly neural architecture descent: a general approach for growing neural networks. \emph{NeurIPS}, 2020. arXiv:2102.08574.

\bibitem{dohare2024} S.~Dohare, J.~F.~Hernandez-Garcia, Q.~Lan, P.~Rahman, A.~R.~Mahmood, R.~S.~Sutton. Loss of plasticity in deep continual learning. \emph{Nature}, 2024. arXiv:2306.13812.

\bibitem{sokar2023redo} G.~Sokar, R.~Agarwal, P.~S.~Castro, U.~Evci. The dormant neuron phenomenon in deep reinforcement learning. \emph{ICML}, 2023.

\bibitem{rebuffi2017icarl} S.-A.~Rebuffi, A.~Kolesnikov, G.~Sperl, C.~H.~Lampert. iCaRL: incremental classifier and representation learning. \emph{CVPR}, 2017.\unv

\bibitem{zhou2023pycil} D.-W.~Zhou, F.-Y.~Wang, H.-J.~Ye, D.-C.~Zhan. PyCIL: a Python toolbox for class-incremental learning. \emph{Science China Information Sciences}, 2023.\unv

\bibitem{saha2021gpm} G.~Saha, I.~Garg, K.~Roy. Gradient projection memory for continual learning. \emph{ICLR}, 2021. arXiv:2103.09762.

\bibitem{wang2021nscl} S.~Wang, X.~Li, J.~Sun, Z.~Xu. Training networks in null space of feature covariance for continual learning. \emph{CVPR}, 2021.

\bibitem{lin2022trgp} S.~Lin, L.~Yang, D.~Fan, J.~Zhang. TRGP: trust region gradient projection for continual learning. \emph{ICLR}, 2022. arXiv:2202.02931.

\bibitem{saha2023sgp} G.~Saha, K.~Roy. Continual learning with scaled gradient projection. \emph{AAAI}, 2023. arXiv:2302.01386.

\bibitem{kao2021ncl} T.-C.~Kao, K.~T.~Jensen, G.~M.~van de Ven, A.~Bernacchia, G.~Hennequin. Natural continual learning: success is a journey, not (just) a destination. \emph{NeurIPS}, 2021.

\bibitem{pan2020fromp} P.~Pan, S.~Swaroop, A.~Immer, R.~Eschenhagen, R.~E.~Turner, M.~E.~Khan. Continual deep learning by functional regularisation of memorable past. \emph{NeurIPS}, 2020.

\bibitem{magistri2024efc} S.~Magistri, T.~Trinci, A.~Soutif-Cormerais, J.~van de Weijer, A.~D.~Bagdanov. Elastic feature consolidation for cold start exemplar-free incremental learning. \emph{ICLR}, 2024. arXiv:2402.03917.

\bibitem{delange2023gap} M.~De Lange, G.~M.~van de Ven, T.~Tuytelaars. Continual evaluation for lifelong learning: identifying the stability gap. \emph{ICLR}, 2023. arXiv:2205.13452.

\bibitem{harun2024gap} M.~Y.~Harun, C.~Kanan. Overcoming the stability gap in continual learning. \emph{TMLR}, 2024. arXiv:2306.01904.

\bibitem{bian2024cflat} Make continual learning stronger via C-Flat. arXiv:2404.00986, 2024 (venue NeurIPS 2024 chưa xác nhận).

\bibitem{cflatpp} C-Flat++: towards a more efficient and powerful framework for continual learning. arXiv:2508.18860, 2025.

\bibitem{udwadia1992} F.~E.~Udwadia, R.~E.~Kalaba. A new perspective on constrained motion. \emph{Proc. R. Soc. A}, 1992.\unv

\bibitem{udwadia2005} F.~E.~Udwadia. Reflections on the Gauss principle of least constraint. \emph{J. Optim. Theory Appl.}, 2005.

\bibitem{nakamura1987} Y.~Nakamura, H.~Hanafusa, T.~Yoshikawa. Task-priority based redundancy control of robot manipulators. \emph{Int. J. Robotics Research}, 1987.\unv

\bibitem{siciliano1991} B.~Siciliano, J.-J.~E.~Slotine. A general framework for managing multiple tasks in highly redundant robotic systems. \emph{ICAR}, 1991.\unv

\bibitem{chiaverini1997} S.~Chiaverini. Singularity-robust task-priority redundancy resolution for real-time kinematic control of robot manipulators. \emph{IEEE Trans. Robotics and Automation}, 1997.\unv

\bibitem{wang2025lorapro} Z.~Wang, J.~Liang, R.~He, Z.~Wang, T.~Tan. LoRA-Pro: are low-rank adapters properly optimized? \emph{ICLR}, 2025. arXiv:2407.18242.

\bibitem{cai2022netaug} H.~Cai, C.~Gan, J.~Lin, S.~Han. Network augmentation for tiny deep learning. \emph{ICLR}, 2022. arXiv:2110.08890.

\bibitem{vwn2025} ByteDance Seed. Virtual width networks. arXiv:2511.11238, 2025 (báo cáo kỹ thuật).

\bibitem{fpe2025} Expand neurons, not parameters. arXiv:2510.04500, 2025 (preprint).

\bibitem{dehaan2026control} R.~de Haan, Taoudi-Benchekroun, P.~Vilimelis Aceituno, B.~F.~Grewe. Continual learning through control minimization. \emph{ICML}, 2026.

\bibitem{richter1957} H.~Richter. Parameterfreie Abschätzung und Realisierung von Erwartungswerten. \emph{Blätter der DGVFM}, 1957.\unv

\bibitem{rogosinski1958} W.~W.~Rogosinski. Moments of non-negative mass. \emph{Proc. R. Soc. A}, 1958.\unv

\bibitem{more1983} J.~J.~Moré, D.~C.~Sorensen. Computing a trust region step. \emph{SIAM J. Sci. Stat. Comput.}, 4(3), 1983.
\bibitem{madsen2004} K.~Madsen, H.~B.~Nielsen, O.~Tingleff. \emph{Methods for Non-Linear Least Squares Problems}, 2nd ed. IMM, DTU, 2004.
\bibitem{nocedal2006} J.~Nocedal, S.~J.~Wright. \emph{Numerical Optimization}, 2nd ed. Springer, 2006.
\bibitem{more1978} J.~J.~Moré. The Levenberg--Marquardt algorithm: implementation and theory. In \emph{Numerical Analysis}, LNM 630, Springer, 1978.
\bibitem{higham2002} N.~J.~Higham. \emph{Accuracy and Stability of Numerical Algorithms}, 2nd ed. SIAM, 2002.
\bibitem{frangella2023} Z.~Frangella, J.~A.~Tropp, M.~Udell. Randomized Nyström preconditioning. \emph{SIAM J. Matrix Anal. Appl.} 44(2), 2023.
\bibitem{fan1949} K.~Fan. On a theorem of Weyl concerning eigenvalues of linear transformations. \emph{PNAS}, 1949.\unv

\bibitem{carsonhigham2018} E.~Carson, N.~J.~Higham.
Accelerating the Solution of Linear Systems by Iterative Refinement in Three Precisions.
\emph{SIAM J. Sci. Comput.}, 40(2), 2018. DOI: 10.1137/17M1140819.

\end{thebibliography}


\end{document}
